% Canonical entry point for Arithmetic Expression Geometry V.
\documentclass[11pt]{article}

\usepackage[letterpaper,margin=1in]{geometry}
\usepackage[T1]{fontenc}
\usepackage[utf8]{inputenc}
\usepackage[english]{babel}
\usepackage{lmodern}
\usepackage{amsmath,amssymb,amsthm,mathtools}
\usepackage{booktabs,longtable,array}
\newcolumntype{L}[1]{>{\raggedright\arraybackslash}p{#1}}
\usepackage{graphicx}
\usepackage{microtype}
\usepackage{xcolor}
\usepackage{tikz}
\usetikzlibrary{arrows.meta,positioning,calc,fit,decorations.pathreplacing}
\usepackage{enumitem}
\usepackage{url}
\usepackage[
  colorlinks=true,
  linkcolor=blue!45!black,
  citecolor=green!35!black,
  urlcolor=blue!55!black,
  pdfencoding=auto
]{hyperref}
\hypersetup{
  pdftitle={Arithmetic Expression Geometry V: Triadic Characteristics and Threaded Apertures},
  pdfauthor={Mingli Yuan},
  pdfsubject={Three Readings, Observation Fibres, and the Syntax of Apertures},
  pdfkeywords={arithmetic expression geometry, triadic characteristic, observation fibre, residual, chirality, through-crossing, aperture, linear logic, normalization}
}

\newtheorem{theorem}{Theorem}[section]
\newtheorem{proposition}[theorem]{Proposition}
\newtheorem{lemma}[theorem]{Lemma}
\newtheorem{corollary}[theorem]{Corollary}
\theoremstyle{definition}
\newtheorem{definition}[theorem]{Definition}
\newtheorem{example}[theorem]{Example}
\newtheorem{convention}[theorem]{Convention}
\newtheorem{construction}[theorem]{Construction}
\newtheorem{calibration}[theorem]{Calibration}
\newtheorem{proposal}[theorem]{Proposal}
\newtheorem{conjecture}[theorem]{Conjecture}
\newtheorem{nogo}[theorem]{No-go}
\theoremstyle{remark}
\newtheorem{remark}[theorem]{Remark}
\newtheorem{warning}[theorem]{Warning}
\newtheorem{problem}[theorem]{Open problem}

\newcommand{\R}{\mathbb{R}}
\newcommand{\C}{\mathbb{C}}
\newcommand{\F}{\mathbb{F}}
\newcommand{\Z}{\mathbb{Z}}
\newcommand{\N}{\mathbb{N}}
\newcommand{\Q}{\mathbb{Q}}
\newcommand{\PP}{\mathbb{P}}
\newcommand{\No}{\mathbf{No}}
\newcommand{\PGL}{\operatorname{PGL}}
\newcommand{\PSL}{\operatorname{PSL}}
\newcommand{\GL}{\operatorname{GL}}
\newcommand{\Aff}{\operatorname{Aff}}
\newcommand{\Hist}{\operatorname{Hist}}
\newcommand{\Obs}{\operatorname{Obs}}
\newcommand{\Fib}{\operatorname{Fib}}
\newcommand{\Fill}{\operatorname{Fill}}
\newcommand{\Sect}{\operatorname{Sect}}
\newcommand{\Fix}{\operatorname{Fix}}
\newcommand{\im}{\operatorname{im}}
\newcommand{\rank}{\operatorname{rank}}
\newcommand{\id}{\operatorname{id}}
\newcommand{\sgn}{\operatorname{sgn}}
\newcommand{\lcm}{\operatorname{lcm}}
\newcommand{\dd}{\,\mathrm{d}}
\newcommand{\defeq}{\mathrel{:=}}
\newcommand{\cA}{\mathcal A}
\newcommand{\cB}{\mathcal B}
\newcommand{\cC}{\mathcal C}
\newcommand{\cD}{\mathcal D}
\newcommand{\cE}{\mathcal E}
\newcommand{\cF}{\mathcal F}
\newcommand{\cG}{\mathcal G}
\newcommand{\cH}{\mathcal H}
\newcommand{\cI}{\mathcal I}
\newcommand{\cK}{\mathcal K}
\newcommand{\cL}{\mathcal L}
\newcommand{\cM}{\mathcal M}
\newcommand{\cN}{\mathcal N}
\newcommand{\cO}{\mathcal O}
\newcommand{\cP}{\mathcal P}
\newcommand{\cQ}{\mathcal Q}
\newcommand{\cR}{\mathcal R}
\newcommand{\cS}{\mathcal S}
\newcommand{\cT}{\mathcal T}
\newcommand{\cU}{\mathcal U}
\newcommand{\cV}{\mathcal V}
\newcommand{\cW}{\mathcal W}
\newcommand{\cX}{\mathcal X}
\newcommand{\cZ}{\mathcal Z}
\newcommand{\bW}{\mathbb{W}}
\newcommand{\bE}{\mathbb{E}}
\newcommand{\bH}{\mathbb{H}}
\newcommand{\bP}{\mathbb{P}}

\setlist[itemize]{topsep=0.35em,itemsep=0.2em,leftmargin=1.8em}
\setlist[enumerate]{topsep=0.35em,itemsep=0.2em,leftmargin=2.1em}

\title{\textbf{Arithmetic Expression Geometry V:\\Triadic Characteristics and Threaded Apertures}\\[0.45em]
\large Three Readings, Observation Fibres, and the Syntax of Apertures}

\author{Mingli Yuan\\Swarma Research\\Beijing 100083, China}
\date{October 20, 2026\\[0.25em]
\small Draft manuscript under mathematical review\\[0.35em]
\small
\textbf{MSC 2020:} 03F52 (primary); 03F05, 15A66;
14D05, 18B10, 68Q55}

\begin{document}
\maketitle

\begin{quote}\small
\textbf{Draft notice.}
This manuscript is a working draft in the editorial and publication sense, and it is
the first manuscript of a slot in the series that the earlier governance files did
not yet cover; its governing scope is
\texttt{governance/00c-paper-V-triadic-aperture-scope.md}.  Here ``draft'' is not a
mathematical claim-status label in the repository governance system and does not
alter the stated status of any theorem, proposition, calibration, proposal,
conjecture, or open problem.
\end{quote}

% Abstract provenance (compressed 2026-10-29 from the 730-word version).
% Every sentence is a selection from sentences of the long abstract; no claim was
% added or strengthened.
%  S1  <- P1 sents 2-3; Secs 1-2
%  S2  <- P2 sents 2-3; Sec 2 and the catalogue of Sec 10
%  S3  <- P2 sents 4-6; Prop 3.2 (prop:quadratic-structure), Prop 3.6
%         (prop:negative-inverse), Prop 3.9 (prop:order-two-four)
%  S4  <- P2 sents 7-8; Thm 4.9 (thm:endogenous-clifford), Thm 4.10
%         (thm:interchange-insufficient)
%  S5  <- P2 sents 12-13; Prop 7.10 (prop:legendre-three-readings), Thm 8.6
%         (thm:monodromy-difference), Prop 8.11 (prop:shear-transvection)
%  S6  <- P2 sent 15; Thm 7.17 (thm:defect-decomposes)
%  S7  <- P3 sents 1-2; Prop 11.8 (prop:seven-faces), Thm 12.4 (thm:adequacy)
%  S8  <- P3 sents 3-6; Thm 13.3 (thm:soundness), Thm 13.5 (thm:incompleteness),
%         Thm 14.3 (thm:substitution-splice), Thm 14.9 (thm:root-transport),
%         Thm 14.12 (thm:strong-normalization), Thm 15.4 (thm:strong-diamond),
%         Cor 15.5 (cor:confluence), Thm 15.10 (thm:subformula), Thm 15.12
%         (thm:decidability)
%  S9  <- P3 sent 8; Thm 16.3 (thm:counting-bound); relative initiality as
%         stated in Sec 16 and Sec 17
%  S10 <- P4, the two-part discipline
%
% Source restoration of 2026-10-29: Sections 13-15 had drifted from the Adva
% research notes they cite.  The grading is again the free abelian resource
% grading of note 0084 (def:grade, prop:grade-invariant,
% cor:graded-nonderivability, ex:graded-witness); the descent measures are again
% the ledger counts of note 0086 (def:proof-ledger, thm:ledger-descent); and the
% discharged occurrence is again the right-boundary occurrence (def:tnd0-rules).
% The abstract sentence above was adjusted to match.
\begin{abstract}This paper develops a theory of three coordinated typed readings of one carrier,
temporal, spatial, and constructive, and of the syntax of the apertures through
which they are filled and checked.  A finite observer contract fixes an
observation map, an equivalence, projectors, and a budget; a characteristic is then
an operational quotient together with a residual, and the recurring theorem shape is
negative: values, gradients, spectra, and energies do not recover history.

In the quadratic family $A_\kappa=\mathbb R[j]/(j^2-\kappa)$ every member satisfies
rank-one recovery, conjugation, and a multiplicative norm, and a globally defined
negative-inverse involution, equivalently an oriented order-four lift with
$J^2=-I$ over $[J]^2=[I]$, selects the complex member;
chirality is the central sign that projectivization forgets.  On a finite cut
interval the six-state carrier $V\oplus\Lambda^2V$ is not postulated: cut-membership
cochains with oriented interchange generate the Clifford anticommutator, whereas
interchange alone provably cannot.  In the Legendre family one cusp is at once a
branch-factor collision, a vanishing class carrying a Dehn transvection, and a
period monodromy; a through-crossing written as a cospan with a resolution torsor
derives $M_{\mathrm{Legendre}}=T_\delta^2$ and realizes the Picard--Lefschetz
transvection on a finite cellular annulus.  On the finite Legendre
prism the total defect $(-I,[\infty,1,0])$ decomposes into two channels, so no
irreducible triadic $3$-cocycle is exhibited.

Part~II takes the typed hole as primitive and proves threaded logic adequate on a
finite terminal-world carrier under four explicit hypotheses.  For the ordered
linear fragment $\mathrm{TND}_0$ it proves soundness and, through a resource grading,
incompleteness for Boolean semantics; an ordered one-hole substitution theorem; strong normalization, via ledger transport lifting through every
context with two decreasing measures; a one-step strong diamond, hence local and
global confluence and a unique normal form per fixed proof; the subformula property;
and a terminating sound and complete focused decision procedure.  A six-port whole-cut kernel answers ``why six'' by the counting bound
$n\ge 2m$, hence $n_{\min}=6$, under a stated contract, universality being weakened
to relative initiality.  Throughout, a quotient is admissible only with a declared
residual and a calibration is never promoted to a theorem; triadicity is certified,
not assumed.\end{abstract}

\noindent\textbf{Keywords:}
arithmetic expression geometry; triadic characteristic; observation fibre;
operational quotient; residual; chirality; oriented exchange lift; Clifford relation;
through-crossing; resolution torsor; monodromy; satisfaction form; aperture; typed
hole; linear logic; strong normalization; confluence; focused derivability; six-port
kernel

\setcounter{tocdepth}{2}
\tableofcontents
\newpage

\part{Triadic characteristic geometry}

\section{Introduction and programme}\label{sec:introduction}
% Section 1: Introduction and programme
\label{sec:intro-body}

\subsection{The question this paper asks}

Papers~0--IV of this series follow one expression through four successively coarser
stages: an expression tree, a marked history, the operator that the history induces,
and the endpoint value where it lands \cite{Yuan2026AEGPaperZero,Yuan2026AEGFoundations}.
Paper~IV then asks when a quotient of that data is still sufficient for future
continuation, and separates a semantic residual from an operational trace
\cite{Yuan2026AEGCondensation}.

This paper asks a different question, which the earlier papers deliberately do not
answer.  Suppose one and the same carrier is to be read in three coordinated ways ---
say as a process, as a space, and as a construction --- and suppose each reading is
a typed interpretation with its own direction of information flow.  What is the
minimal structure that lets such readings be compared at all?  What is necessarily
lost in the comparison, and how is that loss to be accounted for?  And what is the
syntax in which the places of loss --- the places where a reading is *not yet*
determined --- can be written down, filled, and checked?

The answer developed here has two halves, and they are the two parts of the paper.

\subsection{Part I in one paragraph}

Part~I develops a theory of \emph{triadic characteristics}.  Its primitive is not a
value but an \emph{observer contract}: a finite datum fixing what is seen, when two
things are declared indistinguishable, which projections are available, and what
budget is granted.  A characteristic of a process is then an operational quotient of
that process together with a residual, and the residual is not a bookkeeping
afterthought: it is the part of the construction that the quotient cannot see.

Given that framing, several classical-looking questions become sharp.

\begin{itemize}
\item Among the one-parameter quadratic scalar families
  $A_\kappa=\R[j]/(j^2-\kappa)$, which one does the geometry force?  The answer is
  that nothing in rank-one characteristic recovery, conjugation, or a multiplicative
  quadratic norm forces any of them; what selects $\C$ is a global negative-inverse
  involution on the punctured carrier, and this is equivalent, under one explicit
  bridge hypothesis, to the existence of an oriented order-four exchange lift
  $J$ satisfying $J^2=-I$ whose projective class satisfies $[J]^2=[I]$.
  \emph{Chirality} is precisely the central sign that projectivization forgets.
\item Where does the six-state carrier $V\oplus\Lambda^2 V$ come from?  Not from a
  postulate: on a declared finite cut interval, cut-membership cochains together
  with oriented interchange generate the Clifford anticommutator, whereas
  interchange alone provably cannot.
\item What survives when one finite characteristic algebra is read temporally,
  spatially, and constructively?  A stable variance signature, inference as
  constrained intersection of observation fibres, and a hard separation between four
  numeric invariants that a naive theory conflates.
\item What is the right notion of \emph{crossing}?  Not a transport arrow, but a
  cospan with a resolution torsor; monodromy is then the difference of two
  resolutions, and the Legendre monodromy is \emph{derived} from the quadratic base
  change rather than inserted by hand.
\end{itemize}

Part~I also contains the paper's principal \emph{negative} calibration.  On the
finite Legendre prism the total defect decomposes into a central sign channel and a
history-suffix channel; therefore no irreducible triadic $3$-cocycle is exhibited
there.  We report this as a failure of the finite model, not as a success in
disguise.

\subsection{Part II in one paragraph}

Part~II takes seriously the place where a reading is undetermined.  It makes the
typed hole --- the \emph{aperture} --- the primitive: syntactically empty,
semantically a family of fillings, and with open and close operations that are
deliberately not inverse.

On that primitive it builds a finite semantics (a terminal-world carrier, filling
fibres, and an adequacy theorem for threaded logic under explicit hypotheses), an
ordered linear natural-deduction fragment that is sound but provably incomplete for
Boolean semantics, and a proof theory strong enough to be interesting: an ordered
one-hole substitution law with a ledger splice, conditional fresh-cut admissibility,
strong normalization via two decreasing measures, local and global confluence with a
unique normal form per fixed proof, and finally a focused normal-form characterization
with the subformula property that yields a terminating, sound, and complete decision
procedure for derivability.

Part~II closes with the six-port whole-cut kernel, which answers the question ``why
six ports?'' by a counting theorem $n\ge 2m$ under a stated contract, and which
weakens universality to relative initiality rather than claiming it.

\subsection{The discipline the two parts share}

Three rules govern every claim in this paper, and the reader is invited to hold the
paper to them.

\begin{enumerate}
\item \textbf{A quotient is admissible only with a declared residual.}  Every
  quotient result must either be accompanied by an exact reconstruction statement or
  name a residual whose contents are characterized.  The phrase ``the quotient loses
  nothing important'' is not an available move in this paper.
\item \textbf{A calibration is never promoted to a theorem.}  Where a statement is
  established only by an exhaustive check of a declared finite fixture, it is
  labelled a calibration, and the fixture is named.  Where it is established by an
  argument, the argument appears in the paper or in an appendix.
\item \textbf{A property is certified, not named.}  Three readings are triadic only
  if the pairwise crossings are given and the cyclic residual is exhibited.
  Universality is only relative initiality unless a stronger theorem is supplied.
  A boundary-return word is not a periodic orbit.
\end{enumerate}

These rules are the reason Part~I contains as many no-go results as positive ones.
They are also the reason Part~II spends effort on metatheorems: soundness alone is
cheap, and the interesting question is always what fails.

\subsection{Reading guide and dependencies}

A reader interested only in the geometry of readings may read Part~I, taking
\S\ref{sec:observers}--\S\ref{sec:triadic-inference} as the core and
\S\ref{sec:classical}--\S\ref{sec:through} as the calibrations that constrain it.
A reader interested in the syntax of apertures may read Part~II from
\S\ref{sec:apertures}, taking \S\ref{sec:ordered-linear}--\S\ref{sec:confluence-decidability}
as its proof-theoretic content.  The two parts meet at the notion of a filling fibre:
in Part~I a residual is what a quotient cannot see, and in Part~II an aperture is
where a partial reading records that it cannot yet see.

The paper imports, and does not redefine:

\begin{itemize}
\item from Paper~0, the hyperbolic model, the affine cocycle and its relative defect,
  the M\"obius involution $z\mapsto -1/z$, and the cut-and-glue language
  \cite{Yuan2026AEGPaperZero};
\item from Paper~I, the separation of literal syntax, marked history, induced
  operator, and endpoint value, and the projective-versus-affine distinction
  \cite{Yuan2026AEGFoundations};
\item from Paper~IV, the requirement that a quotient be judged against a declared
  continuation class and observation, and the separation of a semantic bound from an
  operational trace \cite{Yuan2026AEGCondensation}.
\end{itemize}

\subsection{Conventions and notation}

Throughout, $\R$ and $\C$ are the real and complex fields, $\F_2$ the field with two
elements, and $\PP^1$ the projective line.  For a finite set $X$ we write
$\cP(X)$ for its powerset and $\mathbf 2=\{\bot,\top\}$ for the Sierpinski lattice
with $\bot<\top$.

An \emph{observer} is written $Q$; its components are introduced in
\S\ref{sec:observers} and are never abbreviated away.  A \emph{characteristic} is
written $\chi$, and $\Fib_{D,Q}(x)$ denotes an observation fibre; the subscript
always names both the domain and the observer, because omitting either is the most
common source of error in this material.

Crossings are written with subscripts naming the two readings being compared, so
that $J_{\cK\cX}$ and $J_{\cX t}$ are different maps and composition order matters.
The orientation expression of the six-state carrier is $\Omega$, and it is a
\emph{derived} object (\S\ref{sec:six-state}) whenever the derivation is available.

In Part~II, a hole is written $h_d$ with $d$ its domain slot; $\Fill_{\cM,\rho}$ is
the filling family in model $\cM$ under assignment $\rho$; $\mathrm{TND}_0$ names the
ordered linear natural-deduction fragment of \S\ref{sec:ordered-linear}.  The
dedicated environments \textsf{Calibration}, \textsf{Proposal}, and
\textsf{Conjecture} carry the epistemic weight that the words suggest, and a
\textsf{Theorem} or \textsf{Proposition} is never used for a finite fixture check.

\subsection{What this paper does not establish}

For emphasis, and because the material is exploratory in places, the following are
\emph{not} claimed anywhere in this paper.

\begin{itemize}
\item That the three named readings are canonically identified.  They are compared
  through declared crossings with declared residuals; a canonical identification is
  an open problem (\S\ref{sec:boundaries}).
\item That arithmetic programs must use complex foundational scalars.  Only that,
  within the tested quadratic family and under a stated bridge hypothesis, one
  property selects $\C$ (\S\ref{sec:scalar-trichotomy}).
\item That the six-state carrier is a spacetime decomposition or a physical claim of
  any kind.  The derivation is finite and combinatorial (\S\ref{sec:six-state}).
\item That any finite calibration demonstrates universality, non-computability, or a
  physical realization.
\item That order can be inferred from an energy, a grade, or a scalar.  The opposite
  is proved in several places (\S\ref{sec:satisfaction-logic},
  \S\ref{sec:residual-catalogue}).
\item That the triadic programme has been completed.  Its decisive missing theorem
  is named in \S\ref{sec:boundaries}.
\end{itemize}

\section{Observers, fibres, and the residual principle}\label{sec:observers}
% Section 2: Observers, fibres, and the residual principle
\label{sec:observers-body}

Everything in this paper is relative to what some observer is allowed to see.  This
section fixes that relativity once, so that later statements about "what cannot be
recovered" have a definite scope.  The material is developed on one deliberately
small worked model, the binary affine monoid; nothing in \S\ref{sec:triadic-inference}
or later depends on that model being the only one available.

\subsection{The observer contract}

\begin{definition}[Observer contract]\label{def:observer-contract}
Let $\cD$ be a class of processes.  A \emph{finite observer contract} on $\cD$ is a
tuple
\[
  Q=(F_Q,\equiv_Q,\Pi_Q,B_Q)
\]
consisting of
\begin{enumerate}
\item an \emph{observation map} $F_Q\colon\cD\to\cO_Q$ into a codomain of readable
  data, possibly many-sorted;
\item an \emph{equivalence} $\equiv_Q$ on $\cO_Q$, declaring when two readable data
  are indistinguishable to $Q$;
\item a family of \emph{projectors} $\Pi_Q$, each specifying which part of an
  observation is retained by a given reading;
\item a \emph{budget} $B_Q$, bounding the finite computations $Q$ may perform.
\end{enumerate}
\end{definition}

The budget is part of the contract, not an implementation detail.  A claim of the
form "no observer of class $\cO$ can determine $X$" has no content unless $\cO$ is
fixed, and $B_Q$ is where the finiteness of the observation enters.

\begin{remark}
The four components are not independent.  Enlarging a budget can refine
$\equiv_Q$, and adding a projector can shrink $F_Q$'s image.  Every statement below
that quantifies over observers therefore names the contract it uses.
\end{remark}

\subsection{Fibres and characteristics}

\begin{definition}[Observation fibre]\label{def:observation-fibre}
For $o\in\cO_Q$ and a reading $D\in\{t,X,\cK\}$ (temporal, spatial, constructive), the
\emph{fibre} of $o$ under $D$ is
\[
  \Fib_{D,Q}(o)=\{\,p\in\cD : F_{D,Q}(p)\equiv_Q o\,\}.
\]
\end{definition}

\begin{definition}[Characteristic]\label{def:characteristic}
A \emph{characteristic} of $p$ for the contract $Q$ is a triple
\[
  \chi_Q(p)=\bigl(\bar\chi_Q(p),\,\pi_p,\,R_{Q}(p)\bigr),
\]
where $\bar\chi_Q(p)$ is the operational quotient (the readable value),
$\pi_p$ is a certificate that $p$ maps to it, and $R_Q(p)$ is the \emph{residual}:
the data retained about $p$ that the quotient does not determine.
\end{definition}

\begin{remark}[A characteristic is not a value]
Definition~\ref{def:characteristic} is where this paper departs from the usual
usage.  A spectrum, a gradient, an eigenvalue, or a scalar is at best
$\bar\chi_Q$; the residual is part of the object, and propositions below exhibit
cases where discarding it destroys exactly the information a later reading needs.
\end{remark}

\subsection{The residual principle}

The following principle is used as an admissibility test throughout; it is the
Part~I counterpart of the rule that no quotient may be used without a declared
residual.

\begin{proposal}[Residual principle]\label{prop:residual-principle}
A quotient $\bar\chi_Q$ of a class $\cD$ may be used as a state for a declared class
of continuations and observations only when accompanied by
\begin{enumerate}
\item an exact reconstruction statement for the continuations in question, or
\item a residual $R_Q$ whose contents are characterized, together with a stated
  sense in which the continuation cannot depend on what $R_Q$ omits.
\end{enumerate}
\end{proposal}

This is a \textsf{Proposal} rather than a theorem because options (i) and (ii) are
requirements on presentation, not mathematical assertions.  Its bite is in what it
forbids: a quotient presented without either is \emph{not} eligible as a state.

\subsection{Identifiability is relative to a hypothesis class}

The worked model is the binary affine monoid.  Let
\[
  f_{k,b}(x)=2^{k}x+b,\qquad k,b\in\Z,
\]
with composition in execution order.  Writing $S$ for $f_{0,1}$ (successor) and $D$
for $f_{1,0}$ (doubling), every element has a normal form $D^{k};S^{b}$.

\begin{lemma}[Affine product law]\label{lem:affine-product}
For $f_{i}=f_{k_i,b_i}$ executed in the order $f_1$ then $f_2$,
\[
  f_2\odot f_1
  =\bigl(k_1+k_2,\ 2^{k_2}b_1+b_2\bigr),
\]
with identity $f_{0,0}$.  In particular
\[
  \chi(S;D)=(1,2)\neq(1,1)=\chi(D;S).
\]
\end{lemma}

\begin{proof}
$f_2(f_1(x))=2^{k_2}(2^{k_1}x+b_1)+b_2=2^{k_1+k_2}x+2^{k_2}b_1+b_2$.  The two
composites differ in their second coordinate because $2\cdot1+0=2$ while
$2^{0}\cdot 0+1=1$.
\end{proof}

Identifiability of $(k,b)$ is not an absolute property; it is relative to a declared
class of hypotheses and to the probes the contract allows.

\begin{proposition}[Two-probe temporal recovery]\label{prop:two-probe}
Probing the temporal reading at $x=0$ and $x=1$ recovers $(k,b)$ uniquely within the
affine class: with $y_i=f_{k,b}(i)$ one has $b=y_0$ and $2^{k}=y_1-y_0$.
\end{proposition}

\begin{proof}
Immediate from $f_{k,b}(0)=b$ and $f_{k,b}(1)=2^{k}+b$.
\end{proof}

\begin{remark}[One probe is not enough]\label{rem:one-probe}
A single probe at $x=1$ cannot separate $S$ from $D$, since $f_{0,1}(1)=2=f_{1,0}(1)$.
The two-probe result is therefore a statement about the contract, not about the
process.
\end{remark}

The same phenomenon appears on the spatial side, and the contrast is instructive:
there a single \emph{exact interval} suffices, whereas a radius alone does not.

\begin{proposition}[Spatial recovery from one exact interval]\label{prop:interval-recovery}
Let the spatial reading act on closed intervals by inverse image,
$\rho_X(k,b)\bigl(I(m,r)\bigr)=I\bigl((m-b)/2^{k},\,r/2^{k}\bigr)$.  From the single
observation $I(m,r)\mapsto I(m',r')$ one recovers $2^{k}=r/r'$ and hence
$b=m-2^{k}m'$.
\end{proposition}

\begin{proof}
The observed image has midpoint $(m-b)/2^{k}$ and radius $r/2^{k}$.  Dividing gives
$2^{k}$, and substituting gives $b$.
\end{proof}

\begin{warning}
Proposition~\ref{prop:interval-recovery} does not say that radii determine
translations.  With the radius alone the observation is $r\mapsto r/2^{k}$, which is
blind to $b$ entirely.  The informative observer is doing the work, not the
mathematics of affine maps.
\end{warning}

\subsection{The four-valued outcome, and what ``unknown'' means}

Inference in this setting is not a function into answers; it is a computation that
may fail to conclude.  We therefore insist on four outcomes.

\begin{definition}[Outcome discipline]\label{def:outcome-discipline}
An inference query returns exactly one of \emph{unique}, \emph{ambiguous},
\emph{inconsistent}, or \emph{unknown}.  An empty candidate set is reported as
\emph{inconsistent} only when exhaustion of the declared search space is certified;
a budget expiry or an incomplete enumeration is reported as \emph{unknown}.
\end{definition}

\begin{warning}[Failure to find is not a proof of absence]
\label{warn:unknown-not-nonexistence}
Definition~\ref{def:outcome-discipline} forbids the inference
``search found nothing, therefore nothing exists''.  This prohibition recurs in
Part~II, where an exhausted search and a certified empty fibre are deliberately
distinguished (\S\ref{sec:adequacy}).
\end{warning}

\subsection{The quotient is lossy, and the loss is the point}

It is tempting to treat a characteristic as a compressed description of a process.
The affine model shows immediately that compression and loss are the same phenomenon.

\begin{proposition}[Quotient without section]\label{prop:lossy-quotient}
Let $\chi_\cK$ be the constructive characteristic that records the parameter pair of
a normal form.  Then
\[
  \chi_\cK(S;D)=\chi_\cK(D;S;S)=(1,2),
\]
so $\chi_\cK$ does not distinguish the raw histories, although
$\chi_\cK\circ\operatorname{NF}=\id$ holds on parameter pairs.  Consequently
$\operatorname{NF}\circ\chi_\cK\neq\id$ on raw histories.
\end{proposition}

\begin{proof}
The composites are $f_{1,0}\circ f_{0,1}=f_{1,2}$ and
$f_{0,1}\circ f_{0,1}\circ f_{1,0}=f_{1,2}$ by Lemma~\ref{lem:affine-product}, while
the words $S;D$ and $D;S;S$ are distinct.
\end{proof}

\begin{proposition}[Compression is inseparable from loss]\label{prop:compression-loss}
The characteristic of $S;D^{n}$ is the pair $(n,2^{n})$, whose canonical word
$D^{n};S^{2^{n}}$ has $2^{n}$ successor symbols, while the parameter pair requires
$O(n)$ bits.  The compression ratio is therefore exponential, and it is purchased
exactly by the identification in Proposition~\ref{prop:lossy-quotient}.
\end{proposition}

\begin{proof}
By induction on $n$ using Lemma~\ref{lem:affine-product}: $f_{n,2^{n}}\odot S=f_{n,2^{n}+1}$
and doubling increments the first coordinate.  The normal form of $(n,2^{n})$ is
$D^{n};S^{2^{n}}$ by construction.
\end{proof}

\begin{remark}
Propositions~\ref{prop:lossy-quotient} and~\ref{prop:compression-loss} are the
smallest instance of the pattern that organises the whole paper: a quotient is
useful in proportion to what it forgets, and admissible in proportion to whether the
forgotten part is accounted for.
\end{remark}

\subsection{Summary of what this section fixes}

The section has introduced the observer contract, the fibre, the characteristic with
its residual, and the outcome discipline.  Two things should be carried forward.

First, \emph{identifiability is a property of a contract}, not of a process
(Proposition~\ref{prop:two-probe} and Remark~\ref{rem:one-probe}).

Second, \emph{a quotient must be presented with a residual}
(Proposition~\ref{prop:residual-principle} and
Proposition~\ref{prop:lossy-quotient}).  Every later section that computes a
characteristic states which of the two options in
Proposition~\ref{prop:residual-principle} it uses.

\section{The quadratic trichotomy and the central sign}\label{sec:scalar-trichotomy}
% Section 3: The quadratic trichotomy and the central sign
\label{sec:scalar-body}

A recurring temptation in this material is to declare the complex field as the
foundational scalar and then to derive the geometry.  This section shows why that
order is unavailable.  Within the natural one-parameter family of two-dimensional real
algebras, the coarse requirements that a "characteristic" theory might impose are
satisfied by all three members; only a genuinely geometric condition separates them,
and that condition has a precise meaning: it is the retention of a central sign.

\subsection{The quadratic family}

\begin{definition}[Quadratic scalar family]\label{def:quadratic-family}
For $\kappa\in\{-1,0,1\}$ let
\[
  A_\kappa=\R[j]/(j^{2}-\kappa),
\]
so that $a+bj$ with $a,b\in\R$ multiplies by
\[
  (a+bj)(c+dj)=(ac+\kappa bd)+(ad+bc)j .
\]
\end{definition}

\begin{proposition}[Elementary structure]\label{prop:quadratic-structure}
In each $A_\kappa$: the quadratic norm $N_\kappa(x)=x\bar x$ with
$\overline{a+bj}=a-bj$ is multiplicative, and conjugation is a multiplicative
involution.  Moreover for every $u\in A_\kappa$ the rank-one map $L_u(x)=ux$
satisfies
\[
  L_u(1)=u,\qquad L_u\circ L_v=L_{uv},
\]
so a scalar is recovered from its own characteristic action without coordinates, and
the three aspect units $\tau,\sigma,\rho\in A_\kappa$ compose cyclically to a single
characteristic $L_{\rho\sigma\tau}$.
\end{proposition}

\begin{proof}
Multiplicativity of $N_\kappa$ is the identity
$(ac+\kappa bd)^{2}-\kappa(ad+bc)^{2}=(a^{2}-\kappa b^{2})(c^{2}-\kappa d^{2})$,
which is a direct expansion.  The rest is immediate from associativity of $A_\kappa$
and $L_u(1)=u$.
\end{proof}

The three cases are conventionally named by the sign of the norm:

\begin{center}
\begin{tabular}{@{}ccll@{}}
\toprule
$\kappa$ & algebra & $N_\kappa(a,b)$ & type \\
\midrule
$-1$ & complex numbers & $a^{2}+b^{2}$ & elliptic \\
$0$ & dual numbers & $a^{2}$ & parabolic \\
$1$ & split-complex numbers & $a^{2}-b^{2}$ & hyperbolic \\
\bottomrule
\end{tabular}
\end{center}

\begin{remark}[What Proposition~\ref{prop:quadratic-structure} does not do]
\label{rem:not-enough}
It does not select $\kappa=-1$.  Rank-one recovery of a characteristic from its own
action, multiplicative conjugation, and a multiplicative quadratic norm are all
satisfied by the three members.  Any argument that begins "the scalar must be
complex, since only then can one conjugate and take norms" is therefore invalid.
\end{remark}

\subsection{The three nondegeneracy conditions, and why they are not derivations}

Three further requirements do separate the cases.

\begin{proposition}\label{prop:three-separators}
Within $\{A_\kappa\}$ each of the following conditions holds for $\kappa=-1$ alone:
\begin{enumerate}
\item every nonzero element is invertible;
\item $N_\kappa$ is positive definite;
\item the unit-norm group $\{x:N_\kappa(x)=1\}$ is compact.
\end{enumerate}
\end{proposition}

\begin{proof}
For $\kappa=-1$ the inverse is $(a+bj)^{-1}=(a-bj)/(a^{2}+b^{2})$ and the unit group
is $S^{1}$.  For $\kappa=0$, every $bj$ with $b\neq0$ is a noninvertible nonzero
element.  For $\kappa=1$, $(1+j)(1-j)=0$ exhibits nonzero zero divisors and the norm
is indefinite with noncompact unit locus.
\end{proof}

\begin{warning}
Proposition~\ref{prop:three-separators} is a separation result, not a derivation.
None of the three conditions is forced by the geometry of readings.  They are
\emph{available} selection principles, and choosing one is a mathematical decision
that must be declared.
\end{warning}

\subsection{The negative-inverse criterion}

A better candidate comes from the geometry itself.  Paper~0 already contains the
M\"obius involution $S(z)=-1/z$ as the reciprocal reading
\cite{Yuan2026AEGPaperZero}; the question is which quadratic scalars admit it
globally.

\begin{proposition}[Negative inverse in the quadratic family]\label{prop:negative-inverse}
For $z=a+bj\in A_\kappa$ with $a^{2}-\kappa b^{2}\neq0$,
\[
  S_\kappa(a,b)=\Bigl(\frac{-a}{a^{2}-\kappa b^{2}},\ \frac{b}{a^{2}-\kappa b^{2}}\Bigr)
\]
satisfies $z\,S_\kappa(z)=-1$.  For $\kappa=-1$ the formula
\begin{enumerate}
\item is defined on the whole punctured plane $\R^{2}\setminus\{0\}$;
\item is an involution;
\item satisfies $zS_{-1}(z)=-1$ identically;
\item preserves the upper half-plane $\{b>0\}$.
\end{enumerate}
For $\kappa=0$ the entire nilpotent axis $a=0$ is excluded.  For $\kappa=1$ the two
light rays $a=\pm b$ are excluded, and the sign of the second coordinate need not be
preserved across the components they separate.
\end{proposition}

\begin{proof}
$1/z=(a-bj)/(a^{2}-\kappa b^{2})$, so $-1/z=S_\kappa(a,b)$ as displayed, and
$zS_\kappa(z)=-1$ follows.  For $\kappa=-1$ the denominator $a^{2}+b^{2}$ never
vanishes off the origin, and $b/(a^{2}+b^{2})>0$ whenever $b>0$; involutivity is the
usual $(S\circ S)(z)=z$ for a M\"obius involution on a chart where it is defined.  The
excluded loci for $\kappa=0,1$ are exactly the zero divisors.
\end{proof}

\begin{proposal}[The AEG inversion criterion]\label{prop:inversion-criterion}
Among the tested family, $\C$ is selected by the requirement that a globally defined
negative-inverse involution exist on every nonzero scalar and preserve one connected
real half-plane.
\end{proposal}

This is labelled \textsf{Proposal} because the requirement is imported from the
observed transition between two real grids, not derived from program geometry; what
is proved is only that it selects.

\subsection{The oriented lift, and chirality as a central sign}

The inversion criterion has a sharper form which explains \emph{why} it selects, and
which supplies the paper's working definition of chirality.

\begin{definition}[Time--space exchange and its lift]\label{def:exchange-lift}
On a homogeneous pair $(t,s)$ put
\[
  J(t,s)=(-s,t),\qquad J=\begin{pmatrix}0&-1\\ 1&0\end{pmatrix}.
\]
On the projective chart $z=t/s$ this induces $[t:s]\mapsto[-s:t]$, i.e.
$z\mapsto-1/z$.  We call $J$ the \emph{oriented lift} of the projective exchange.
\end{definition}

\begin{proposition}[Order two versus order four]\label{prop:order-two-four}
$[J]^{2}=[I]$ on the projective line, while $J^{2}=-I$ on the oriented lift.  Hence
$J^{-1}=-J$, and the forward and reverse exchanges induce the same projective
involution although they are distinct on the lift.
\end{proposition}

\begin{proof}
$J^{2}=\begin{pmatrix}-1&0\\0&-1\end{pmatrix}$.  Projectively $[-\id]=[\id]$, so
$[J]^{2}=[\id]$, while on the lift the scalar $-1$ is retained.
\end{proof}

\begin{definition}[Chirality]\label{def:chirality}
\emph{Chirality} is the central sign distinguishing the forward and reverse oriented
lifts of one projective exchange: the datum that $J$ and $J^{-1}=-J$ are different
lifts of the same projective involution.
\end{definition}

This definition is deliberately austere.  It says nothing about the orientation of the
ambient plane, and the following computation shows that it must not.

\begin{proposition}[Chirality is not ambient orientation reversal]\label{prop:not-orientation}
In real coordinates $S(x,y)=\bigl(-x/(x^{2}+y^{2}),\,y/(x^{2}+y^{2})\bigr)$ has
Jacobian determinant
\[
  \det DS=\frac{1}{(x^{2}+y^{2})^{2}}>0 .
\]
Thus the involution is holomorphic and orientation-preserving on the plane; the
chirality change is not a reversal of that plane but the difference between the two
lifts of one projective map.
\end{proposition}

\begin{proof}
$S$ is the M\"obius map $-1/z$ with derivative $1/z^{2}$, so
$\det DS=|z|^{-4}=(x^{2}+y^{2})^{-2}>0$.
\end{proof}

\subsection{The generator action under one bridge hypothesis}

The central sign becomes a selection mechanism for the algebra when the same oriented
two-dimensional carrier is asked to represent both the exchange lift and the action of
the quadratic generator.

\begin{proposition}[Generator action]\label{prop:generator-action}
Let $L_j$ be multiplication by $j$ on $(a,b)$, so that
\[
  L_j=\begin{pmatrix}0&\kappa\\ 1&0\end{pmatrix},\qquad L_j^{2}=\kappa I .
\]
Then:
\begin{enumerate}
\item for $\kappa=-1$, $L_j=J$ is orientation-preserving, invertible, of order four
  before projectivization;
\item for $\kappa=1$, $L_j^{2}=+I$, $\det L_j=-1$, and the generator is already its own
  inverse on the lift;
\item for $\kappa=0$, $L_j$ is singular and nilpotent.
\end{enumerate}
\end{proposition}

\begin{proof}
Direct computation for each $\kappa$; $\det L_j=-\kappa$ and $L_j^{2}=\kappa I$.
\end{proof}

\begin{calibration}[The bridge is a hypothesis]\label{cal:bridge-hypothesis}
Requiring an oriented time--space exchange whose projective class has order two while
its lift distinguishes forward from reverse selects $j^{2}=-1$, \emph{provided} the
same oriented two-dimensional carrier is declared to represent both the exchange lift
and the generator action.  Both halves are calibrated, not proved: the selection is a
finite algebraic check, and the bridge is a declared identification.
\end{calibration}

\begin{warning}
The matrix $J\in SL(2,\R)$ already defines the real-coefficient M\"obius map $-1/z$;
complex notation is therefore not by itself evidence that the underlying scalars are
complex.  What is new is that retaining the lift rather than quotienting to
$\PSL(2,\R)$ exposes the central sign.  Calling that sign chirality requires the
oriented carrier to remember it.
\end{warning}

\subsection{The conjugated presentation}

The last piece of the trichotomy is what the involution does to the two AEG
generators, because it shows the central sign in operational form.

\begin{proposition}[Conjugated generators]\label{prop:conjugated-generators}
Write $A(z)=z+1$ and $M_\lambda(z)=\lambda z$ for $\lambda>0$ in the upper
half-plane chart.  Then
\[
  S M_\lambda S=M_{\lambda^{-1}},\qquad
  S A S(z)=\frac{z}{1-z}.
\]
Thus conjugation inverts the scale direction while turning the translation grid into a
parabolic grid organised around the other fixed point.  For a Baumslag--Solitar-type
relation $MAM^{-1}=A^{n}$ the conjugated presentation has scale generator $M^{-1}$ and
ripple generator $SAS$.
\end{proposition}

\begin{proof}
For the first identity, $S M_\lambda S(z)=-1/(\lambda\cdot(-1/z))=z/\lambda
=M_{\lambda^{-1}}(z)$.  For the second,
$SAS(z)=-1/((-1/z)+1)=z/(1-z)$, which is parabolic with the single fixed point $1/2$.
\end{proof}

\begin{remark}[Baumslag--Solitar background]
The relation $MAM^{-1}=A^{n}$ is the Baumslag--Solitar presentation
\cite{BaumslagSolitar1962NonHopfian}; Paper~0 records the exact routing of
copy/branch/recombine programs through these relations.  The content here is only that
conjugation exchanges the forward and reverse presentations of the selected scale
direction, which is the visible consequence of lift-level chirality.
\end{remark}

\subsection{What the trichotomy settles, and what it leaves open}

\begin{calibration}[Bounded red-team calibration]\label{cal:trichotomy-bounded}
Within the tested quadratic family: two real dimensions per aspect, cyclic rank-one
action, recovery of a characteristic from its unit action, conjugation, and a
multiplicative norm do \emph{not} select $\C$; a global negative-inverse involution
preserving a connected real AEG half-plane \emph{does} select $\C$; and the
mechanism is that the projective exchange has order two while its
orientation-preserving lift has order four.
\end{calibration}

Three interpretations of the six real degrees of freedom remain live and are not
separated by anything proved here.

\begin{proposal}[Three readings of $6=2+2+2$]\label{prop:three-readings}
The cases $\kappa=-1,0,1$ may respectively encode phase/orientation, infinitesimal
extension, and causal scale; alternatively the three pairs may be three complex
lines; alternatively $\C$, dual, and split behaviour may be charts or strata of one
expression-valued scalar geometry.  No canonical assignment of the three cases to
temporal, spatial, and relational roles is established.
\end{proposal}

\begin{remark}[The Jordan shoulder]
The dual case is not a pathology to be eliminated.  A nilpotent extension naturally
records a first-order shoulder that survives complexification, which suggests that a
faithful spectrum may require both a semisimple character and a residual extension
object.  \S\ref{sec:residual-catalogue} returns to this.
\end{remark}

\section{The six-state carrier and the derived Clifford relation}\label{sec:six-state}
% Section 4: The six-state carrier and the derived Clifford relation
\label{sec:six-state-body}

Three readings, each two-dimensional over the reals, give six real dimensions.  The
temptation is to equip each pair with a complex structure and declare the result
$\C^{3}$.  This section shows that the correct carrier is not a sum of three copies of
anything: it is the six-dimensional space
\[
  W=\Lambda^{1}E\oplus\Lambda^{2}E,\qquad \dim_{\R}E=3,
\]
in which each reading is paired with the \emph{oriented relation of the other two},
and in which the complex unit appears exactly once.  The section then shows that on a
declared finite cut interval this structure is not postulated but generated.

\subsection{The carrier and its paired planes}

\begin{definition}[Six-state carrier]\label{def:six-state-carrier}
Let $E$ be a three-dimensional real vector space with basis $T,S,R$ representing the
temporal, spatial, and relational readings.  Put
\[
  W=\Lambda^{1}E\oplus\Lambda^{2}E=\langle T,S,R\rangle\oplus\langle SR,RT,TS\rangle ,
\]
so $\dim_{\R}W=6$.  The three \emph{aspect--opposite-relation planes} are
\[
  W_T=\langle T,SR\rangle,\qquad W_S=\langle S,RT\rangle,\qquad W_R=\langle R,TS\rangle .
\]
\end{definition}

The pairing is the content of the definition: a reading is not paired with a second
copy of itself but with the oriented plane spanned by the other two.  This is forced
by the dimension count once one asks that each reading be recoverable from its own
opposite relation, and it has an immediate consequence for the complex structure.

\begin{proposition}[One complex unit, not three]\label{prop:one-complex-unit}
Suppose the generators satisfy
\[
  T^{2}=S^{2}=R^{2}=1,\qquad \varepsilon_i\varepsilon_j+\varepsilon_j\varepsilon_i=0
  \ \ (i\neq j),
\]
and set $\Omega=TSR$.  Then $\Omega$ is central and
\[
  \Omega^{2}=-1 .
\]
Consequently $J=\Omega\cdot$ preserves $W$, swaps $\Lambda^{1}E$ with
$\Lambda^{2}E$ by cyclic pairing
\[
  J(T)=SR,\quad J(SR)=-T;\qquad J(S)=RT,\quad J(RT)=-S;\qquad
  J(R)=TS,\quad J(TS)=-R,
\]
and on each plane $W_T,W_S,W_R$ restricts to
$\begin{pmatrix}0&-1\\ 1&0\end{pmatrix}$.  Hence
$\chi_J(\lambda)=(\lambda^{2}+1)^{3}$.
\end{proposition}

\begin{proof}
\emph{Centrality.}  A generator $T$ commutes with $SR$ up to $(-1)^{2}=+1$, so
$\Omega$ lies in the centre.

\emph{The square.}  Moving the fourth factor of $\Omega^{2}=T\,S\,R\,T\,S\,R$ to the
front costs the two transpositions past $R$ and $S$, hence $(-1)^{2}=+1$:
\[
  T\,S\,R\,T\,S\,R = T^{2}\,S\,R\,S\,R = S\,R\,S\,R ,
\]
using $T^{2}=1$.  Moving the third factor past the fourth costs one transposition:
\[
  S\,R\,S\,R = -\,S^{2}R^{2} = -1 .
\]
Hence $\Omega^{2}=-1$.

\emph{Pairing.}  For the first plane, $RT=-TR$ gives
\[
  \Omega T=T\,S\,R\,T=T\,S\,(-T\,R)=-(T\,S\,T)R=(T\,T\,S)R=T^{2}SR=SR,
\]
and then $\Omega(SR)=T\,S\,R\,S\,R=-T\,S^{2}R^{2}=-T$.  The other two planes follow by
cyclic symmetry.  On the plane spanned by $T$ and $SR$ the action is therefore
$T\mapsto SR\mapsto -T$, which is the displayed matrix.

\emph{Characteristic polynomial.}  $J$ acts as the same $2\times2$ rotation on each of
the three planes, so
$\chi_J(\lambda)=\prod_{i}\det\bigl(\lambda I-J|_{W_i}\bigr)=(\lambda^{2}+1)^{3}$.
\end{proof}

\begin{calibration}[The algebraic form is an assumption]\label{cal:algebraic-form}
Proposition~\ref{prop:one-complex-unit} takes the anticommutation relations as
\emph{hypotheses}.  It is a bounded algebraic calibration: it shows what the carrier
does once such generators exist.  \S\ref{sec:six-state-derivation} supplies the
derivation.
\end{calibration}

\begin{proposition}[Chirality sign law]\label{prop:sign-law}
For every permutation $\pi$ of the three directions,
\[
  P_\pi\,\Omega\,P_\pi^{-1}=\sgn(\pi)\,\Omega .
\]
In particular exchanging temporal and spatial directions gives
$\Omega\mapsto-\Omega=\Omega^{-1}$, so cyclic permutations preserve the orientation
expression and transpositions reverse it.
\end{proposition}

\begin{proof}
$P_\pi$ permutes the generators, hence multiplies the word $TSR$ by the sign of the
permutation.  Invariantly: a transposition reverses the induced orientation of $E$,
and $\Omega$ is the orientation class of $E$ followed by the volume of
$\Lambda^{2}E$ under $\Omega$'s own duality.
\end{proof}

\begin{remark}
Proposition~\ref{prop:sign-law} is the carrier-level form of
Definition~\ref{def:chirality}: the central sign is exactly what a transposition
produces, and it is invisible to any projective reading.
\end{remark}

\subsection{The signature trichotomy}

If the three directions are granted arbitrary signs, the complex unit becomes
conditional.

\begin{proposition}[Signature trichotomy]\label{prop:signature-trichotomy}
With $T^{2}=\varepsilon_T$, $S^{2}=\varepsilon_S$, $R^{2}=\varepsilon_R$ and
$\varepsilon_i\in\{-1,0,+1\}$,
\[
  \Omega^{2}=-\varepsilon_T\varepsilon_S\varepsilon_R ,
\]
so the carrier is elliptic when the product is positive, hyperbolic when negative, and
parabolic when the product vanishes.
\end{proposition}

\begin{proof}
Each $T^{2}$ contributes a factor $\varepsilon_T$ to the reduction performed in the
proof of Proposition~\ref{prop:one-complex-unit}; the signs multiply.
\end{proof}

\begin{warning}[Temporality is not a metric sign]\label{warn:temporality}
A Lorentzian reading with exactly one negative square forces the split case.  Hence
semantic temporality must \emph{not} be identified with a negative metric direction;
doing so silently converts the trichotomy of \S\ref{sec:scalar-trichotomy} into an
unannounced choice of $\kappa=1$.
\end{warning}

\subsection{Deriving the generators from cuts}\label{sec:six-state-derivation}

The material above is algebraic.  The geometric content of this section is that on a
declared finite cut interval the two operations needed are already present, so the
Clifford relation is \emph{derived}.

\begin{construction}[Local cut interval]\label{con:cut-interval}
Let $U$ be a finite set of events carrying the completion order of a checked acyclic
program, and let $g_1,g_2,g_3$ be three pairwise independent enabled events
generating a local Boolean cut interval.  Two operations are available on cut data:
\begin{enumerate}
\item \emph{oriented interchange}, which exchanges two independent enabled events and
  records the sign of the exchange, giving endomorphisms $\varepsilon_i$ on the cut
  cochains;
\item \emph{cut contraction}, which reads whether an event belongs to a cut, giving
  $\iota_i$ determined by the membership pairing
  \[
    dx_i(g_j)=\delta_{ij}.
  \]
\end{enumerate}
\end{construction}

\begin{theorem}[Endogenous Clifford relation]\label{thm:endogenous-clifford}
On the local cut interval of Construction~\ref{con:cut-interval} the following hold.
\begin{enumerate}
\item Oriented interchange yields the exterior relations
  $\varepsilon_i\varepsilon_j+\varepsilon_j\varepsilon_i=0$ for $i\neq j$.
\item With membership pairing $dx_i(g_j)=\delta_{ij}$, the operators
  $c_i=\varepsilon_i+\iota_i$ satisfy
  \[
    c_ic_j+c_jc_i=2\delta_{ij}I .
  \]
\item Hence $\Omega=c_1c_2c_3$ satisfies $\Omega^{2}=-I$, and on
  $W=\Lambda^{1}E\oplus\Lambda^{2}E$ the orientation action has
  $\chi_J(\lambda)=(\lambda^{2}+1)^{3}$ with three aspect--opposite-face planes.
\end{enumerate}
\end{theorem}

\begin{proof}[Proof sketch and status]
Items (i) and (ii) are the Chevalley construction of a Clifford algebra from an
exterior algebra together with a contraction satisfying the duality
$dx_i(g_j)=\delta_{ij}$; what the construction adds is that both operations are
\emph{supplied by the cut data} rather than assumed.  Item (iii) is then
Proposition~\ref{prop:one-complex-unit} applied to the derived generators.
The derivation is checked exhaustively on one local three-cube and is reported as a
bounded endogenous derivation, not as a theorem about arbitrary cut intervals.
\end{proof}

\begin{theorem}[Interchange alone is insufficient]\label{thm:interchange-insufficient}
Oriented interchange by itself yields only a nilpotent exterior volume: it produces
$\varepsilon_1\varepsilon_2\varepsilon_3$ satisfying
$(\varepsilon_1\varepsilon_2\varepsilon_3)^{2}=0$, and no complex structure.  The
contraction $\iota_i$ is indispensable.
\end{theorem}

\begin{proof}
Wedge products of degree three in a three-dimensional exterior algebra have square
zero because any repeated factor annihilates the product.  Without a contraction
there is no degree-lowering operation, so the Chevalley combination $c_i=\varepsilon_i
+\iota_i$ cannot be formed.
\end{proof}

\begin{remark}[Why this matters]
Theorems~\ref{thm:endogenous-clifford} and~\ref{thm:interchange-insufficient}
together answer the natural objection to \S\ref{sec:scalar-trichotomy}: the complex
structure is not smuggled in with the carrier.  It is produced by two operations that
the cut interpretation already needs, and the second of them is doing essential work.
\end{remark}

\subsection{Filtrations from a designated causal line}

One further structure is available, and it is important not to over-read it.

\begin{proposition}[Six-state filtration]\label{prop:filtration}
Designating one causal line refines the six-state carrier as
\[
  6=1+2+2+1 .
\]
The two readings $(1,2,3)$ and $(1,3,2)$ obtained by regrouping the transverse pairs
differ by an orientation-unit transfer and are \emph{not} a chiral pair.
\end{proposition}

\begin{proof}[Status]
The count is elementary from
Definition~\ref{def:six-state-carrier}.  The absence of chirality is the content: see
Theorem~\ref{thm:regrouping-not-chirality}.
\end{proof}

\begin{theorem}[No-go: regrouping is not chirality reversal]\label{thm:regrouping-not-chirality}
Orientation-unit transfer changes block membership but not $\Omega$, whereas an actual
odd exchange of two transverse directions satisfies $P\Omega P^{-1}=-\Omega$.  Hence
the passage $(1,2,3)\leftrightarrow(1,3,2)$ is a change of filtered reading and not a
chiral pair.  Consequently an ordinary scalar spectrum cannot distinguish the two
readings, while the filtration can.
\end{theorem}

\begin{proof}
The two facts are verified separately: the transfer preserves the orientation word by
Proposition~\ref{prop:sign-law} applied to an even permutation, while an odd exchange
is a transposition and therefore reverses it.  A scalar spectrum depends only on
$\Omega^{2}$ and hence cannot see the sign at all.
\end{proof}

\begin{remark}
Theorem~\ref{thm:regrouping-not-chirality} is a small result with a large moral, which
recurs in \S\ref{sec:residual-catalogue}: quantities built from squares, norms, or
spectra are blind precisely to the central sign that
Definition~\ref{def:chirality} isolates.
\end{remark}

\section{Triadic characteristic inference}\label{sec:triadic-inference}
% Section 5: Triadic characteristic inference
\label{sec:triadic-inference-body}

\S\ref{sec:observers} fixed what an observer is and what a characteristic costs.
This section reads one finite characteristic algebra in three ways at once and asks
what the three readings share.  The answer has two parts: a variance signature that is
stable under every calibration performed so far, and an inference procedure that is
nothing more than the intersection of three fibres.  Both are developed on the
affine monoid of Lemma~\ref{lem:affine-product} and then calibrated on the square and
the cube in \S\ref{sec:multiplicity}.

\subsection{The three readings}

\begin{definition}[Typed readings]\label{def:typed-readings}
Let $\cA$ be a class of processes closed under a composition $\odot$ read in execution
order.  A \emph{typed reading family} on $\cA$ consists of
\begin{enumerate}
\item a \emph{temporal} reading $\rho_t$, covariant in execution, taking values in
  endomorphisms of the state space;
\item a \emph{spatial} reading $\rho_X$, contravariant, taking values in inverse-image
  maps on a space of admissible regions;
\item a \emph{constructive} reading $\rho_\cK$, covariant modulo a declared
  canonicalisation $\equiv_{\mathrm{can}}$, taking values in construction data.
\end{enumerate}
\end{definition}

The three readings have different variance by design: a process acts forward on states,
backward on regions, and forward again on constructions.  This asymmetry is not a
defect; it is what the phrase "three readings" is meant to encode.

\begin{proposition}[Variance signature]\label{prop:variance-signature}
Let $f_1$ be executed before $f_2$ and write $f_2\odot f_1$ for the composite.  Then
\[
  \rho_t(f_2\odot f_1)=\rho_t(f_2)\circ\rho_t(f_1),\qquad
  \rho_X(f_2\odot f_1)=\rho_X(f_1)\circ\rho_X(f_2),
\]
\[
  \rho_\cK(f_2\odot f_1)\equiv_{\mathrm{can}}\rho_\cK(f_2)\circ\rho_\cK(f_1).
\]
Schematically, writing $\mathfrak T$ for the triple,
\[
  \mathfrak T(f_2\odot f_1)=\bigl(T_2\circ T_1,\ X_1\circ X_2,\ K_2\circ K_1
  \bmod \mathrm{can}\bigr),
\]
so the signature is $({+},{-},{+})$ in the order $(t,X,\cK)$.
\end{proposition}

\begin{proof}
The temporal and constructive statements are functoriality of $\rho_t$ and of
$\rho_\cK$ modulo canonicalisation.  The spatial statement is contravariance of the
inverse-image construction: a region transformed by the composite $f_1$ then $f_2$ is
$\rho_X(f_2)^{-1}\bigl(\rho_X(f_1)^{-1}(R)\bigr)=\bigl(\rho_X(f_1)\circ\rho_X(f_2)
\bigr)^{-1}(R)$.
\end{proof}

\begin{calibration}[Scope of the signature]\label{cal:variance-scope}
The signature of Proposition~\ref{prop:variance-signature} is established by hand
inside the finite affine model, and re-verified in the calibrations of
\S\ref{sec:multiplicity} and \S\ref{sec:classical}.  It is \emph{not} claimed as a
theorem about every conceivable three-domain model.  What changes from model to model
is the structure each reading must carry; the variance law has so far not changed.
\end{calibration}

\subsection{Inference as intersection of fibres}

\begin{proposition}[Triadic inference]\label{prop:triadic-inference}
For observations $o_t,o_X,o_\cK$ presented to the three readings, the inferred set of
candidates is
\[
  \cI(o_t,o_X,o_\cK)
  =\Fib_{t,Q}(o_t)\cap\Fib_{X,Q}(o_X)\cap\Fib_{\cK,Q}(o_\cK).
\]
Each reading contributes constraints and the constraints are intersected.
\end{proposition}

\begin{proof}
By Definition~\ref{def:observation-fibre}, a process lies in $\Fib_{D,Q}(o_D)$ exactly
when its $D$-reading is $Q$-indistinguishable from $o_D$.  A process satisfying all
three observations therefore lies in the intersection, and conversely.
\end{proof}

\begin{remark}[What "one algorithm" can and cannot mean]\label{rem:one-algorithm}
Proposition~\ref{prop:triadic-inference} says that the three readings share one
\emph{operation} in characteristic space.  It does not say that they share an untyped
decoder; the fibres live in different codomains and the intersection is taken after
transport.  Reading the proposition as a claim about a universal interpreter is a
mistake, and it is the mistake that Part~II is designed to prevent.
\end{remark}

The outcome of the intersection is reported through
Definition~\ref{def:outcome-discipline}.  In particular
Warning~\ref{warn:unknown-not-nonexistence} applies with force here: an empty
intersection is \emph{inconsistent} only if all three searches were exhausted within
$B_Q$.

\subsection{The deformed crossing}

Inference is only half the story; the other half is how readings move past one
another.  The affine model gives an exact law, and it is not a swap.

\begin{proposition}[Order reduction is not a swap]\label{prop:deformed-crossing}
Write $S$ for successor and $D$ for doubling.  Then
\[
  S;D\equiv D;S;S,\qquad\text{while}\qquad S;D\not\equiv D;S,
\]
where $\equiv$ denotes extensional equality of the induced value.  Moving $D$ leftward
across $S$ \emph{duplicates} the translation generator, and every pair has a canonical
representative $D^{k};S^{b}$.
\end{proposition}

\begin{proof}
By Lemma~\ref{lem:affine-product}, $S;D$ has parameters $(1,2)$ while $D;S$ has
$(1,1)$; and $D;S;S$ has $(1,2)$ as well.  Canonicality of $D^{k};S^{b}$ is the normal
form of the monoid.
\end{proof}

\begin{warning}[Extensional, not confluent]\label{warn:not-confluent}
Proposition~\ref{prop:deformed-crossing} is one exact extensional crossing law.  It is
\emph{not} a terminating and confluent rewriting theorem for the abstract
presentation, and this paper does not claim one.  The distinction matters because
Part~II proves genuine normalization theorems under linearity hypotheses, and
Proposition~\ref{prop:deformed-crossing} must not be mistaken for one of them.
\end{warning}

\subsection{The square map: loss, branching, and reverse transport}

The clearest model of what a reading can and cannot retain is the real square.

\begin{proposition}[Temporal loss is an orbit quotient]\label{prop:square-orbit}
Let $q(x)=x^{2}$ and $\sigma(x)=-x$.  Then $q\circ\sigma=q$, so $q$ factors as
\[
  \R\xrightarrow{\ \mathrm{quot}\ } \R/C_2\cong[0,\infty)
  \xrightarrow{\ \hookrightarrow\ } \R .
\]
The generic fibre is a single two-element $C_2$-orbit; $\{0\}$ is a fixed point, so the
quotient is non-free exactly at the origin.
\end{proposition}

\begin{proposition}[Spatial branching is the same quotient, read backwards]
\label{prop:square-covering}
On $\R^{\times}\to\R_{>0}$, $q$ is a two-sheeted covering and $\sigma$ exchanges the
sheets.  For $0<a<b$,
\[
  q^{-1}(a,b)=(-\sqrt b,-\sqrt a)\cup(\sqrt a,\sqrt b),
\]
while $q^{-1}(-1,1)=(-1,1)$ exhibits the sheets merging through the critical source
point.
\end{proposition}

\begin{proposition}[Reverse transport is relation-valued]\label{prop:relation-valued}
The backward reading $q^{\dagger}$ is $\varnothing$, $\{0\}$, or
$\{-\sqrt y,\sqrt y\}$ according as $y<0$, $y=0$, or $y>0$.  Hence a plain global
inverse function does not exist, although forward execution remains functional.
\end{proposition}

\begin{remark}
Propositions~\ref{prop:square-orbit}--\ref{prop:relation-valued} are the smallest
complete instance of the warning that contravariance forces relation-valued transports.
The spatial reading of a critical source is not a function, and any formalism that
demands one will either lose the singularity or invent a branch convention that is not
intrinsic to the data.
\end{remark}

\subsection{Scaling crosses; translation does not}

\begin{proposition}[Scaling crosses by deformation]\label{prop:scaling-crosses}
For $D_a(x)=ax$ and $Q(x)=x^{2}$,
\[
  D_a;Q\equiv Q;D_{a^{2}},
\]
with equal extensional value and different raw histories.
\end{proposition}

\begin{proof}
$q(D_a x)=a^{2}x^{2}$ and $D_{a^{2}}(q(x))=a^{2}x^{2}$.
\end{proof}

\begin{theorem}[No-go: translation has no affine crossing]\label{thm:translation-obstruction}
Let $S_b(x)=x+b$ with $b\neq0$.  There is no affine $A(y)=\alpha y+\beta$ with
\[
  S_b;Q\equiv Q;A .
\]
\end{theorem}

\begin{proof}
Every post-composite of $Q$ with an affine map has the form
$A(q(x))=\alpha x^{2}+\beta$, whose linear coefficient vanishes identically.  But
$q(S_b x)=x^{2}+2bx+b^{2}$ has linear coefficient $2b\neq0$ since the field has
characteristic $0$ and $b\neq0$.
\end{proof}

\begin{remark}[Scope of Theorem~\ref{thm:translation-obstruction}]
\label{rem:obstruction-scope}
The theorem forbids an \emph{affine} crossing.  It does not forbid a non-affine
expansion, and the next result exhibits exactly that.  Reading the obstruction as a
claim about all rewrites would overstate it; this is the sense in which the statement
is a boundary statement about the affine post-normal-form language, not about
reduction in general.
\end{remark}

\begin{proposition}[The word-to-graph boundary]\label{prop:word-to-graph}
Crossing a translation through a square forces a parallel decomposition with explicit
copy:
\[
  x\xrightarrow{\ \Delta_3\ } (x_1,x_2,x_3)
  \xrightarrow{\ Q\otimes D_{2b}\otimes C_{b^{2}}\ } (x_1^{2},\,2bx_2,\,b^{2})
  \xrightarrow{\ \mathrm{add}_3\ } x^{2}+2bx+b^{2}.
\]
Consequently general reduction in this setting is not a global ordering of two word
alphabets, but typed graph rewriting with declared copy and merge.
\end{proposition}

\begin{proof}
The right-hand side of the obstruction is a sum of three terms whose coefficients come
from three distinct uses of one source.  A word over the one-input alphabets $\{S_b\}$
and $\{Q\}$ cannot express three simultaneous uses of its single input; a graph with an
explicit fan-out can.
\end{proof}

\begin{calibration}[Constructive fibre is not a singleton]\label{cal:constructive-fibre}
$\chi_\cK(Q)=\chi_\cK(N;Q)=q$ while the raw histories and the construction data of $Q$
and $N;Q$ remain distinct.  Hence the constructive fibre is not a singleton even when
the temporal and spatial fibres are, and by
Proposition~\ref{prop:residual-principle} a usable characteristic must retain the
residual.
\end{calibration}

\subsection{Summary}

Three results in this section are load-bearing for what follows.

\begin{enumerate}
\item The variance signature $(+,-,+)$ is stable across the models tested
  (Proposition~\ref{prop:variance-signature}, Calibration~\ref{cal:variance-scope}).
\item Crossings are deformations, and translation does not cross an affine normal form
  at all (Propositions~\ref{prop:deformed-crossing}, \ref{prop:scaling-crosses},
  Theorem~\ref{thm:translation-obstruction}).
\item The consequence is a change of carrier: from words to typed graphs with explicit
  copy (Proposition~\ref{prop:word-to-graph}).  Part~II takes up precisely this
  carrier, where the copy discipline becomes a linearity discipline.
\end{enumerate}

\section{Typed multiplicity, equivariance, and the word-to-graph boundary}\label{sec:multiplicity}
% Section 6: Typed multiplicity, equivariance, and the word-to-graph boundary
\label{sec:multiplicity-body}

The temptation produced by \S\ref{sec:triadic-inference} is to suppose that a reading
has a single "multiplicity" --- a degree, a number of preimages, a count of sheets ---
and that these agree.  They do not.  This section separates them, shows that the
correct general form of the stability a characteristic enjoys is equivariance rather
than invariance, and exhibits the coaction responsible for the failure of translation
to cross a power inside a word normal form.

\subsection{Four invariants, not one}

\begin{definition}[Observer-relative degree profile]\label{def:degree-profile}
For a process $P$ and a contract $Q$, the \emph{degree profile} is the quadruple
\[
  \mathbf d_Q(P)=\bigl(d_{\mathrm{alg}},\,d_{\mathrm{fib}},\,d_{\mathrm{sheet}},\,
  d_{\mathrm{occ}}\bigr),
\]
where $d_{\mathrm{alg}}$ is the algebraic degree of the induced characteristic,
$d_{\mathrm{fib}}$ the cardinality of a generic fibre of the reverse reading over the
declared scalar domain, $d_{\mathrm{sheet}}$ the covering degree of the associated
spatial correspondence, and $d_{\mathrm{occ}}$ the constructive occurrence
multiplicity.
\end{definition}

The point of Definition~\ref{def:degree-profile} is that these four quantities are
\emph{typed}: they live in different codomains, are computed by different readings, and
may take different values for one process.

\subsection{The cube counter-calibration}

The square of \S\ref{sec:triadic-inference} is too symmetric to make the point, because
over $\R$ it already has a genuine branch.  The cube is decisive: it is algebraically
of degree three, it is a real homeomorphism, and it has three constructive occurrences.

\begin{calibration}[Power maps separate the invariants]\label{cal:cube-degrees}
For $P_{\R}$ the real cube $x\mapsto x^{3}$ read through the contract $Q_{\R}$, and for
$P_{\C}$ the complex cube read through $Q_{\C}$ whose scalar domain is $\C$,
\[
  \mathbf d_{Q_{\R}}(P_{\R})=(3,1,1,3),\qquad
  \mathbf d_{Q_{\C}}(P_{\C})=(3,3,3,3).
\]
The total multiplicity $3$ appears in both profiles but in different coordinates, so
the two cannot be reconciled by a single number.
\end{calibration}

\begin{warning}[Criticality does not imply branching]\label{warn:criticality}
$C(x)=x^{3}$ has $C'(0)=0$, so the map is critically singular at the origin, yet it is
a real homeomorphism: crucially, criticality does not split a real open neighbourhood
into several branches.  Consequently no argument of the form "the derivative vanishes,
therefore there are several sheets over the image" is admissible.  A branch count is a
statement about the covering over the declared scalar domain and must be computed
there.
\end{warning}

\begin{remark}
Warning~\ref{warn:criticality} is the exact obstruction to the naive reading of
\S\ref{sec:six-state}: the complex structure $J$ detects the presence of a
second direction, not the multiplicity of preimages.  The two data are independent.
\end{remark}

\subsection{Invariance is the trivial case of equivariance}

A characteristic need not be invariant under a symmetry of its carrier; it may
transform.

\begin{proposition}[Equivariance law]\label{prop:equivariance}
Let $\sigma(x)=-x$ and let $P_n(x)=x^{n}$.  Then
\[
  P_n(\sigma x)=\sigma^{n}\,P_n(x),\qquad\text{i.e.}\qquad P_n(-x)=(-1)^{n}P_n(x).
\]
Equivalently, $P_n$ is equivariant for the sign representation of $C_2$ when $n$ is odd,
and invariant --- the trivial representation --- when $n$ is even.
\end{proposition}

\begin{proof}
$(-x)^{n}=(-1)^{n}x^{n}$.
\end{proof}

\begin{proposal}[Equivariant characteristics]\label{prop:equivariant-characteristics}
The correct general notion of "a characteristic that a symmetry preserves" is
equivariance $f(gx)=\pi(g)f(x)$ for a declared representation $\pi$; invariance
$f(gx)=f(x)$ is the case $\pi$ trivial.  Characteristics should be recorded together
with $\pi$.
\end{proposal}

\begin{remark}
Proposition~\ref{prop:equivariance} is elementary, and it is included because the
distinction it names is the one that fails in practice: characters are frequently
assumed invariant, and the resulting theory then appears to lose odd-degree data.  The
data are not lost; the variance was mis-declared.
\end{remark}

\subsection{The binomial coaction}

The final piece explains the obstruction of Theorem~\ref{thm:translation-obstruction}
in structural terms.

\begin{proposition}[Binomial coaction]\label{prop:binomial-coaction}
On $\Q[x]$ define $\Delta(x)=x\otimes 1+1\otimes x$, extended as an algebra map.  Then
for every $b$,
\[
  (\id\otimes \mathrm{ev}_b)\Delta(x^{n})=(x+b)^{n},
\]
where $\mathrm{ev}_b$ evaluates the right tensor factor at $b$.
\end{proposition}

\begin{proof}
$\Delta(x^{n})=(x\otimes1+1\otimes x)^{n}=\sum_{k}\binom{n}{k}x^{n-k}\otimes x^{k}$ by
the binomial theorem in the commutative algebra $\Q[x]\otimes\Q[x]$.  Applying
$\id\otimes\mathrm{ev}_b$ gives $\sum_k\binom{n}{k}b^{k}x^{n-k}=(x+b)^{n}$.
\end{proof}

\begin{corollary}[Why translation expands a word]\label{cor:translation-expands}
The passage of a translation through the $n$-th power is governed by a coaction with
$n+1$ terms; hence it is a multi-term expansion for $n\ge 2$ and cannot be a single
post-composite.  This is the structural reason behind
Theorem~\ref{thm:translation-obstruction}.
\end{corollary}

\begin{proof}
The expansion has $\binom{n}{k}$ nonvanishing summands for $k=0,\dots,n$; for $n\ge2$
there are at least three terms with distinct powers of $x$, whereas a post-composite of
$x^{n}$ with an affine map produces exactly one term of degree $n$ and one constant.
\end{proof}

\begin{remark}[Vocabulary closed under translation]
Proposition~\ref{prop:binomial-coaction} also says which vocabularies are stable under
translation: a power vocabulary is translation-closed exactly when it is closed under
the coaction, which for a finite generating set forces the vocabulary to be
translation-invariant as a set.  \S\ref{sec:boundaries} records the corresponding open
problem for finite vocabularies.
\end{remark}

\subsection{Consequences carried forward}

\begin{enumerate}
\item A characteristic is typed by the reading that produces it, and its multiplicity
  coordinates are not interchangeable
  (Calibration~\ref{cal:cube-degrees}, Warning~\ref{warn:criticality}).
\item Stability under a symmetry is equivariance for a declared representation, not
  invariance (Proposition~\ref{prop:equivariance},
  Proposal~\ref{prop:equivariant-characteristics}).
\item Translation cannot cross a power inside a word normal form, for a structural
  reason given by a coaction (Proposition~\ref{prop:binomial-coaction},
  Corollary~\ref{cor:translation-expands}).  Combined with
  Proposition~\ref{prop:word-to-graph}, this is the paper's justification for replacing
  words by typed graphs with explicit copy.
\end{enumerate}

\section{Classical calibrations: isogenies and the Legendre crossing}\label{sec:classical}
% Section 7: Classical calibrations: isogenies and the Legendre crossing
\label{sec:classical-body}

A programme that claims a genuinely new triadic structure must be tested against the
places where mathematics already has one.  This section reports two calibrations, and
the first is not flattering to the programme: a complex elliptic isogeny already
carries three exactly-bridged representations, so the bare existence of a bridge is
prior art.  What survives as novel is narrower and is stated precisely.  The second
calibration is constructive: in the Legendre family one cusp really does support three
readings at once, and the residual that obstructs their identification is exact.

\subsection{The isogeny calibration: the bridge is not new}

\begin{theorem}[Three readings of an isogeny]\label{thm:isogeny-three-readings}
Let $\phi\colon E_1\to E_2$ be a complex elliptic isogeny with kernel $\ker\phi$, and let
$M_\alpha$ be the matrix of the induced map on the period lattice attached to a complex
multiplication $\alpha$.  Then
\[
  N(\alpha)=|\det M_\alpha|=\deg\phi=|\ker\phi| ,
\]
so the norm, the lattice determinant, the covering degree, and the kernel cardinality
are four descriptions of one number, computed in three different domains.
\end{theorem}

\begin{proof}
For a complex isogeny the degree equals the size of the kernel and the index of the
image lattice, and the lattice index is the absolute determinant of the matrix of the
map in a lattice basis.  For a CM endomorphism the degree is the algebraic norm.
\end{proof}

\begin{calibration}[The novelty boundary]\label{cal:isogeny-novelty}
Theorem~\ref{thm:isogeny-three-readings} shows that a triadic bridge of the kind
developed in \S\ref{sec:triadic-inference} is already realised classically in the
theory of complex elliptic isogenies.  The programme's claimable novelty is therefore
\emph{not} the existence of three bridged representations.  It is, at most:
\begin{enumerate}
\item observer-relativity: the bridge is stated relative to a declared finite contract
  $Q$ and a hypothesis class;
\item finite inference: the reconstruction is an intersection of fibres subject to a
  budget, with a four-valued outcome;
\item accountable residual: where the classical equalities hold, the residual is
  provably empty, and where they fail, the residual is characterized.
\end{enumerate}
\end{calibration}

\begin{warning}[Two traps in the classical case]\label{warn:isogeny-traps}
Two errors are easy to make and both are avoided in the statement above.
First, coordinate branching of a composite $x\circ\phi$ is \emph{not} ramification of
$\phi$; the two are different structures attached to different maps.
Second, the equalities of Theorem~\ref{thm:isogeny-three-readings} are equalities of
\emph{numbers} obtained from three carriers; they do not identify the carriers, and by
Definition~\ref{def:characteristic} an identification of carriers would require a
residual argument that the theorem does not supply.
\end{warning}

\begin{proposal}[Characteristics as a groupoid]\label{prop:characteristic-groupoid}
A characteristic is better modelled as an object of a representation groupoid --- a
carrier together with declared readings and the isomorphisms between them --- than as a
single normal-form value.  Under this reading, the value of a characteristic is the
isomorphism class, and the residual is the stabilizer datum.
\end{proposal}

\subsection{The marked rank-two carrier and its three cusps}

The geometry that the triadic programme wants is supplied by the level-two modular
curve.

\begin{proposition}[Three cusps are three slope classes]\label{prop:three-cusps}
The level-two modular curve
\[
  Y(2)=\Gamma(2)\backslash\bH\cong\PP^{1}\setminus\{0,1,\infty\}
\]
has exactly three cusps, corresponding to the three points of $\PP^{1}(\F_2)$, that is,
to the three slope classes $0/1$, $1/1$, $1/0$.  Equivalently, they are the three
oriented slope classes of one marked rank-two carrier.
\end{proposition}

\begin{proof}[Status]
The isomorphism is the classical description of $Y(2)$ as the Legendre $\lambda$-line;
the three cusps are $\lambda\in\{0,1,\infty\}$.  The slope-class description is obtained
by reading $(\omega_1,\omega_2)$ modulo $\Gamma(2)$ and reducing the ratio modulo $2$.
The corresponding mod-two relations are
$\delta_T=\delta_S+\delta_C$ and its cyclic variants.
\end{proof}

\begin{remark}[Mod-two erases chirality]\label{rem:mod2-erases}
The reduction above is mod-two, and \S\ref{sec:six-state} showed exactly what
mod-two forgets: the sign of the orientation word.  Consequently the identification of
Proposition~\ref{prop:three-cusps} describes the three readings up to a sign that the
mod-two picture cannot see, and an integer lift is required.  The following
construction supplies it.
\end{remark}

\begin{construction}[Oriented integer lift]\label{con:oriented-lift}
Let $\delta_T=a$, $\delta_S=b$, $\delta_C=-a-b$ in a rank-two lattice with symplectic
form $\langle\cdot,\cdot\rangle$, chosen so that all three cyclic pairings equal $+1$:
\[
  \langle\delta_T,\delta_S\rangle=\langle\delta_S,\delta_C\rangle
  =\langle\delta_C,\delta_T\rangle=+1 .
\]
The order-three chart rotation is
\[
  R=\begin{pmatrix}0&-1\\ 1&-1\end{pmatrix},\qquad R^{3}=I .
\]
\end{construction}

\begin{lemma}[Tri-chart coherence]\label{lem:tri-chart}
The three coordinate charts $(x,y)$, $(y-x,-x)$, $(-y,x-y)$ are permuted by $R$, and
$R^{3}=I$.
\end{lemma}

\begin{proof}
Applying $R$ to a pair yields the next chart, by direct substitution.  The
characteristic polynomial of $R$ is $\lambda^{2}+\lambda+1$, so $R^{3}=I$ by
Cayley--Hamilton.
\end{proof}

\begin{remark}
Construction~\ref{con:oriented-lift} and Lemma~\ref{lem:tri-chart} are the smallest exact
repair of Remark~\ref{rem:mod2-erases}: the sign lost mod two is retained by working in
an oriented integer lattice, at the cost of a nontrivial chart rotation that must be
tracked.
\end{remark}

\subsection{The Legendre crossing}

\begin{proposition}[One cusp, three readings]\label{prop:legendre-three-readings}
Consider the Legendre family
\[
  E_\lambda:\ y^{2}=x(x-1)(x-\lambda),\qquad \lambda\in\PP^{1}\setminus\{0,1,\infty\},
\]
with its nonzero two-torsion satisfying $P_0+P_1+P_\lambda=0$.  At each cusp
$\lambda\in\{0,1,\infty\}$, the same local singularity is simultaneously
\begin{enumerate}
\item a collision of branch factors of the defining equation;
\item a vanishing class $\delta$ with Dehn transvection $U_\delta=L_\delta^{2}$;
\item a monodromy of the Picard--Fuchs system
  \[
    \lambda(1-\lambda)\Pi''+(1-2\lambda)\Pi'-\tfrac14\Pi=0 .
  \]
\end{enumerate}
\end{proposition}

\begin{proof}[Status]
Items (i)--(iii) are classical local descriptions of the Legendre degeneration.  The
content claimed here is only that they apply at one point to one object; the
identification of the three \emph{descriptions} with the three \emph{readings} of
\S\ref{sec:triadic-inference} is not proved and remains the open problem of
\S\ref{sec:boundaries}.
\end{proof}

\begin{nogo}[Bare monodromy is spectrally empty]\label{nogo:unipotent}
The cusp monodromy acts by
\[
  M_\delta(v)=v+\langle v,\delta\rangle\,\delta ,
\]
which is unipotent.  Hence a bare cusp monodromy supplies no rich spectrum, and any
useful spectral invariant must be attached to paths, transports, correspondences, or
gluing data rather than to a single local monodromy operator.
\end{nogo}

\begin{proof}
$M_\delta-\id$ has image contained in the line spanned by $\delta$ and has square zero,
so every eigenvalue of $M_\delta$ equals $1$.
\end{proof}

\begin{proposition}[Central sign and projective closure]\label{prop:central-sign-closure}
With positive unipotent cusp lifts $U_0,U_1,U_\infty$ one has
\[
  U_0U_1U_\infty=-I \quad\text{in the linear lift, while}\quad
  M_0M_1M_\infty=I \quad\text{after typing the sign into the infinity chart.}
\]
Thus the three local operators close projectively, and a central $-1$ survives in any
selected linear lift.
\end{proposition}

\begin{proof}[Status]
The projective closure is classical.  The central sign is the same phenomenon as
Proposition~\ref{prop:order-two-four}: projectivization forgets a scalar, and a lift
retains it.
\end{proof}

\begin{proposition}[Label transport, not universal swap]\label{prop:label-transport}
The general law governing how a vanishing cycle moves when the base chart is changed is
\[
  L_{g\delta}=g\,L_\delta\,g^{-1},
\]
a conjugation of labels.  There is no universal swap law for $\delta$ across charts.
\end{proposition}

\begin{proof}
Conjugation transports the cycle $\delta$ to $g\delta$, and the transvection is
determined by its vanishing cycle, so the conjugated operator is the transvection of the
transported cycle.
\end{proof}

\subsection{The crossing-coherence prism and its decomposable defect}

The three cards are now placed on one object, and the result is the paper's principal
negative calibration.

\begin{construction}[Crossing-coherence prism]\label{con:prism}
Take three typed carriers
\[
  \cK=(\text{construction data }q,p,h),\qquad
  \cX=(\text{charge }q\in\Z^{2}),\qquad
  \cT=(\text{covector }\ell\in(\Z^{2})^{\ast}),
\]
with crossing actions $J^{c}_{\cK}$, $J^{c}_{\cX}=M_c$, $J^{c}_{\cT}=M_c^{-T}$ for each
cusp $c$.  The interpretation maps are $F_{\cK\cX}$, $F_{\cX\cT}=J$, and $F_{\cT\cK}$,
with residuals carried by the last.
\end{construction}

\begin{theorem}[Crossing squares]\label{thm:crossing-squares}
In Construction~\ref{con:prism}:
\begin{enumerate}
\item the $\cK$--$\cX$ square commutes strictly on charge;
\item the $\cX$--$\cT$ square commutes strictly, by symplectic duality
  $JM=M^{-T}J$;
\item the $\cT$--$\cK$ square commutes only up to the exact history residual, recorded
  as $[c]$.
\end{enumerate}
\end{theorem}

\begin{proof}
Item (i) is immediate, since both actions act on the same charge coordinate.  Item (ii)
is the defining property of the symplectic form: for $M$ preserving
$\langle\cdot,\cdot\rangle$ the relation $JM=M^{-T}J$ holds.  Item (iii) fails by
construction, since a period covector does not determine construction history
(Proposition~\ref{prop:period-not-history}); the failure is recorded rather than
removed.
\end{proof}

\begin{proposition}[Two independent residual channels]\label{prop:two-residual-channels}
The prism of Construction~\ref{con:prism} carries two independent residuals:
\begin{enumerate}
\item a central lift sign, of which $-I$ versus $+I$ is the visible instance;
\item a history suffix $[c]$ recording the construction path.
\end{enumerate}
A coarse projective observer that identifies $v\sim-v$ and $M\sim-M$, and that erases the
history, sees neither.
\end{proposition}

\begin{proof}
The first is invisible under projectivization by
Proposition~\ref{prop:order-two-four}; the second is invisible because it is not a
function of the charge or of the covector.  Their independence is exhibited by varying
one while fixing the other.
\end{proof}

\begin{theorem}[No-go: the triadic defect decomposes]\label{thm:defect-decomposes}
The total finite defect of the prism is
\[
  D_Q=\bigl(-I,\,[\infty,1,0]\bigr),
\]
and it decomposes into the two channels of
Proposition~\ref{prop:two-residual-channels}.  Consequently no irreducible triadic
$3$-cocycle is exhibited by this model.
\end{theorem}

\begin{proof}
The central sign is a scalar and the history suffix is a word; the projections onto the
two summands are maps of defect data, and each vanishes on the kernel of the other.  A
defect admitting such a decomposition is by definition not irreducible.
\end{proof}

\begin{calibration}[Status of the negative calibration]\label{cal:defect-negative}
Theorem~\ref{thm:defect-decomposes} is reported as the paper's principal negative
calibration.  It establishes that the finite Legendre prism does \emph{not} produce a
genuinely triadic coherence defect; it does not establish that no such defect exists.
\end{calibration}

\begin{proposal}[What genuine triadicity would require]\label{prop:coherence-defect}
Genuine triadicity requires a coherence defect $\Theta_c$ that is not reconstructible
from any pair of the three crossings.  In the language of
\S\ref{sec:satisfaction-logic} this is the statement that the triple support strictly
exceeds what the three pairwise supports determine.
\end{proposal}

\begin{conjecture}[Triadic singular crossing]\label{conj:triadic-crossing}
There exists a carrier on which the three readings of \S\ref{sec:triadic-inference} are
related by crossing characteristics with a non-decomposable cyclic defect, functorial
under refinement of the observer contract $Q$.
\end{conjecture}

\subsection{What this section settles}

\begin{enumerate}
\item The bare triadic bridge is prior art
  (Theorem~\ref{thm:isogeny-three-readings}, Calibration~\ref{cal:isogeny-novelty}); the
  programme's novelty must be sought in observer-relativity and residual accounting.
\item The classical carrier for three readings is the level-two modular curve with its
  three cusps, and the mod-two picture requires an oriented integer lift to retain
  chirality (Proposition~\ref{prop:three-cusps},
  Construction~\ref{con:oriented-lift}, Lemma~\ref{lem:tri-chart}).
\item Crossing coherence is strict for two pairs and residual-laden for the third, and
  the resulting defect decomposes (Theorem~\ref{thm:crossing-squares},
  Proposition~\ref{prop:two-residual-channels}, Theorem~\ref{thm:defect-decomposes}).
\item Whether a non-decomposable defect exists is open
  (Proposal~\ref{prop:coherence-defect}, Conjecture~\ref{conj:triadic-crossing}).
\end{enumerate}

\section{Through-crossings, resolution, and monodromy}\label{sec:through}
% Section 8: Through-crossings, resolution, and monodromy
\label{sec:through-body}

\S\ref{sec:classical} produced local monodromy operators and then observed
(No-go~\ref{nogo:unipotent}) that a single one of them carries no spectrum.  This
section supplies the missing structure.  Its central claim is a change of genre: passing
through a singular fibre is not an arrow but a cospan together with a choice, and once
that is made explicit, monodromy becomes a \emph{difference} of the two choices rather
than a primitive.  The change pays immediately: the Legendre monodromy is derived from
a quadratic base change, and a primitive transvection is realised on a finite cell
complex.

\subsection{Crossing is a cospan with a resolution}

\begin{definition}[Through-crossing]\label{def:through-crossing}
Let $F_0$ be a singular fibre and $F_{\pm}$ the two smooth fibres on either side of the
critical value.  A \emph{through-crossing} is the cospan
\[
  F_-\xrightarrow{\ \pi_- \ } F_0 \xleftarrow{\ \pi_+ \ } F_+ ,
\]
whose \emph{through relation} is the fibre product
\[
  R_c=F_-\times_{F_0}F_+ .
\]
\end{definition}

\begin{proposal}[A crossing is not a transport arrow]\label{prop:crossing-not-arrow}
There is no canonical smooth transport arrow $F_-\to F_+$ across a singular fibre.
What exists is the through relation $R_c$ together with a further choice, and the choice
is the mathematical content.
\end{proposal}

\begin{definition}[Resolution torsor]\label{def:resolution-torsor}
A \emph{resolution} of the through-crossing is a selection, for each point of the
critical fibre, of which smooth point on the negative side is joined to which on the
positive side.  The set of resolutions is a torsor under the monodromy group generated
by the vanishing transvection:
\[
  \mathrm{Res}(R_c)\ \text{is a torsor under}\ \langle T_\delta\rangle\cong\Z .
\]
\end{definition}

\begin{remark}[Why the torsor language is not decoration]
A resolution has no privileged origin: any two resolutions differ by a power of the
transvection, and no resolution is singled out by the geometry of the cospan alone.  This
is the precise sense in which \S\ref{sec:classical}'s monodromy operators are not
intrinsic to a single fibre.
\end{remark}

\begin{remark}[A worked affine instance, imported from Paper 0]
\label{rem:affine-instance-of-the-torsor}
The resolution torsor of Definition~\ref{def:resolution-torsor} is an abstract
statement until an instance is produced.  Paper~0 supplies one in the affine sector,
and we cite it rather than re-derive it.  There,
\(\operatorname{Aff}^{+}(1,\mathbb R)\) acts on \(\mathbb H^2\); a dilation is
hyperbolic and its conjugacy class is determined by its multiplier, so its
translation length is a function of the class; a nontrivial translation is
parabolic, and by the two-layer proposition of Paper~0
\cite{Yuan2026AEGPaperZero} every nontrivial translation of a fixed sign is
conjugate to every other, the stabiliser being the translation subgroup.  The class of a translation therefore carries no
continuous invariant, and the set of its conjugates is a torsor under the dilations
with no privileged origin --- the same shape as
Definition~\ref{def:resolution-torsor}.  Two consequences are worth stating,
because they are the reason this paper does not identify the two torsors.
First, the affine instance is a torsor under \(\mathbb R_{>0}\), not under
\(\mathbb Z\), so it does not yet supply the \(\langle T_\delta\rangle\cong\mathbb Z\)
structure that Theorem~\ref{thm:monodromy-difference} uses.  Second, the passage
from the affine sector to the projective one needs the trace trichotomy
\(\operatorname{tr}^2-4\det\), which Paper~0 does not develop and this paper does
not develop either.  So the instance is a calibration of the \emph{shape}, not an
identification of the objects.
\end{remark}

\subsection{Monodromy as a difference of resolutions}

\begin{theorem}[Monodromy from two through-resolutions]\label{thm:monodromy-difference}
Let $P^{\uparrow}$ and $P^{\downarrow}$ be the two resolutions of
Definition~\ref{def:resolution-torsor} obtained by passing through the critical value
along the two admissible sides.  Then the around-the-critical-value monodromy is
\[
  \bigl(P^{\downarrow}\bigr)^{-1}P^{\uparrow}=T_\delta ,
\]
the vanishing transvection.
\end{theorem}

\begin{proof}
Both resolutions identify the same two smooth fibres, and their composite is the deck
transformation of the passage around the critical value.  For a single node that deck
transformation is the Picard--Lefschetz transvection of the vanishing cycle $\delta$, by
the local theory of nondegenerate critical points.
\end{proof}

\begin{remark}
Theorem~\ref{thm:monodromy-difference} is the section's organising result.  It replaces
"insert the monodromy matrix" by "compare two resolutions", and it explains why
No-go~\ref{nogo:unipotent} had to be true: one resolution alone determines nothing,
because the torsor has no origin.
\end{remark}

\subsection{The Legendre cusp as a quadratic base change}

\begin{proposition}[Local reduction]\label{prop:legendre-reduction}
The local equation of the Legendre family at a cusp reduces, after an analytic change
of coordinates, to
\[
  UV=\lambda^{2}.
\]
The map $\lambda\mapsto s=\lambda^{2}$ is a quadratic base change of the node
$UV=s$, and consequently
\[
  M_{\mathrm{Legendre}}=T_\delta^{2}.
\]
\end{proposition}

\begin{proof}
The Legendre family has an $A_1$ node at $\lambda=0$ (and, after the standard change of
chart, at the other two cusps as well).  Pulling a node back along $\lambda\mapsto\lambda^{2}$
doubles the local monodromy, since a full turn in $s$ is covered by half a turn in
$\lambda$.  By Theorem~\ref{thm:monodromy-difference} the node monodromy is $T_\delta$,
so the pullback has monodromy $T_\delta^{2}$.
\end{proof}

\begin{remark}[What has been gained]
Proposition~\ref{prop:legendre-reduction} derives a monodromy matrix that
\S\ref{sec:classical} could only state.  The derivation uses the \emph{shape} of the
base change, not a computation with periods.  This is the sense in which the through
form is a better carrier than the operator: the operator is its shadow.
\end{remark}

\subsection{Cellular realisation on a finite annulus}

The abstract torsor has an exact finite model.

\begin{construction}[Cubical annulus and shear]\label{con:cubical-annulus}
Let $A_n$ be the cubical annulus obtained from a strip of $n$ squares by identifying the
two boundary circles with a shift, so that
\[
  \chi(A_n)=0,\qquad (b_0,b_1,b_2)(A_n)=(1,1,0).
\]
Let $D$ be the cellular shear that displaces each vertex one step along the annulus, and
let $g\colon A_n\to T_n$ be the gluing to the torus obtained by identifying the two
boundaries.
\end{construction}

\begin{proposition}[Shear induces the transvection]\label{prop:shear-transvection}
The induced map on homology is the Picard--Lefschetz transvection
\[
  D_\ast(x)=x+\langle \delta,x\rangle\,\delta ,\qquad
  M_D=\begin{pmatrix}1&1\\ 0&1\end{pmatrix},\qquad
  M_D^{T}JM_D=J ,
\]
where $J$ is the symplectic form.  Hence the shear is symplectic and realises the
primitive Dehn twist.
\end{proposition}

\begin{proof}
The shear displaces a representative of $\delta$ once around the annulus while fixing
the transverse class, which is precisely the defining formula of the transvection.  The
matrix is read off on the basis $(\delta,\gamma)$ with
$\langle\delta,\gamma\rangle=1$, and symplecticity follows from
$\langle x+\langle\delta,x\rangle\delta,\ y+\langle\delta,y\rangle\delta\rangle
=\langle x,y\rangle$ by antisymmetry and $\langle\delta,\delta\rangle=0$.
\end{proof}

\begin{theorem}[No-go: vertex closure does not imply path closure]
\label{thm:vertex-not-path}
In Construction~\ref{con:cubical-annulus} the vertex-level shear satisfies $D_0^{\,n}=I$,
while the lifted one-chain level satisfies $D_1^{\,n}\neq I$ and
$M_D^{\,n}=\begin{pmatrix}1&n\\ 0&1\end{pmatrix}\neq I$ for $n\neq0$.  Thus closure of a
finite vertex permutation is not the closure of the lifted path.
\end{theorem}

\begin{proof}
The vertex permutation is a cyclic shift of a finite set, hence has finite order dividing
$n$.  The induced homology map is unipotent with nontrivial nilpotent part, so it has
infinite order.
\end{proof}

\begin{remark}[The pair of separations]
Theorem~\ref{thm:vertex-not-path} is the finite model of
Proposition~\ref{prop:crossing-not-arrow}: a closed finite quotient can look trivial
while the lifted object is not.  \S\ref{sec:residual-catalogue} records this as one
instance of a general pattern rather than as an isolated curiosity.
\end{remark}

\begin{proposition}[Marked identity and the nodal pinch]\label{prop:marked-identity}
Let $q\colon T_n\to$ (the nodal quotient) be the pinch.  Then $q_\ast D_\ast=q_\ast$, so
the shear is invisible to the nodal reading; but the marked outgoing condition gives a
strict identity $q_+D=q_-$, which upgrades $D$ to a genuine through-resolution.
\end{proposition}

\begin{proof}
The pinch kills precisely the vanishing direction, and the shear acts on the transverse
direction by the transvection, which is trivial modulo $\delta$.  The marked statement
tracks which boundary the identity is asserted on, whence the strict identity.
\end{proof}

\subsection{Through forms as fibre products over a singular middle}

\begin{definition}[Through form]\label{def:through-form}
For a typed singular middle object $N$ and two fibres $F_\pm$ over it, the \emph{through
form} is
\[
  \langle L\mid_N R\rangle \defeq F_-\times_N F_+ .
\]
\end{definition}

\begin{proposition}[Order and compatibility are independent]\label{prop:order-compat-independent}
Conway number-admissibility of a cut $\{L\mid R\}$ and through-compatibility of
$\langle L\mid_N R\rangle$ are logically independent: either may hold while the other
fails.
\end{proposition}

\begin{proof}
Admissibility is a condition on the order structure of $L$ and $R$ alone; compatibility
is a condition on the projections to $N$.  Choosing $N$ degenerate makes compatibility
fail with admissible sides, and choosing an incompatible order on $N$ makes compatibility
hold with inadmissible sides.
\end{proof}

\begin{proposition}[Resolution fibre count]\label{prop:resolution-count}
If the critical fibre consists of $r$ regular points together with $n$ phase points,
the through form carries $r+n^{2}$ resolutions, indexed by the relative offset
$P_k\colon i\mapsto i+k$ modulo $n$.
\end{proposition}

\begin{proof}
Each regular point has one resolution; each pair of phase points may be joined, and the
join is determined by a relative offset $k$ modulo $n$, giving $n$ classes for each of
$n$ points.
\end{proof}

\begin{remark}[Duality is relational converse]
The duality appropriate to through forms is not an inverse but the relational converse
$T^{\star}=T^{\mathrm{op}}$ with contravariant composition.  This is forced by
Proposition~\ref{prop:relation-valued}: a transport that is relation-valued cannot have a
functional dual.
\end{remark}

\subsection{One feature, three readings}

\begin{theorem}[Figure-eight calibration]\label{thm:figure-eight}
There is a specific figure-eight braid word $w$ for which
\[
  -D_w=-\phi\!\left(\frac{\partial w}{\partial b}\right)
  =\det(tI-M)=t^{2}-3t+1 ,
\]
where $D_w$ is the AEG affine defect, $\phi$ is the abelianised Fox derivative
\cite{Fox1953FreeDifferential}, and $M$ the fibre monodromy.  The common value is a
feature of the ordered history and does not reconstruct the knot.
\end{theorem}

\begin{proof}[Status]
An exact computation for one word, with all three quantities evaluated independently.
The equality is a calibration on a single word, not a general identity, and the
non-reconstruction clause is essential: the polynomial is a characteristic in the sense
of Definition~\ref{def:characteristic}, with the word as its residual.
\end{proof}

\begin{proposition}[Double periods]\label{prop:double-periods}
Two visible periods $p,q$ give a phase torus $\Z/p\times\Z/q$ under the successor
$\tau(i,j)=(i+1,j+1)$, with
\[
  \#\text{orbits}=\gcd(p,q),\qquad \text{orbit length}=\lcm(p,q),\qquad
  \text{primitive winding}=\bigl(q/g,\ p/g\bigr),\quad g=\gcd(p,q).
\]
The unit-tangent lift supplies the third direction, and the orbifold meridians
\[
  m_P=(p-1,-1),\qquad m_Q=(-1,q-1),\qquad \det = pq-p-q
\]
recover the arithmetic.  Equal numerals across the three types are not identified.
\end{proposition}

\begin{proof}
An orbit of $\tau$ has length $\lcm(p,q)$ because both coordinates must return, and the
number of orbits is $pq/\lcm(p,q)=\gcd(p,q)$.  The winding is read off by dividing the
periods by $g$.  The determinant computation is direct.
\end{proof}

\subsection{Return is not period}

\begin{nogo}[No intrinsic periodic dynamics]\label{nogo:return-not-period}
There is a finite diagram containing the source-free components
$\mathrm{constant};\mathrm{discard}$ and
$\mathrm{constant};\mathrm{neg};\mathrm{discard}$ whose least boundary-return lengths are
$2$ and $3$ respectively; but monotonicity of the completed-past rank forces the verdict
\emph{NotRepresentable} for intrinsic periodic dynamics.  Hence a boundary-return word is
not a periodic orbit.
\end{nogo}

\begin{proof}
The completed past of a state grows monotonically along a checked execution, so no state
can recur.  The return lengths are therefore properties of words, not periods of a
dynamical system; a genuine period requires a feedback operator that the syntax does not
supply.
\end{proof}

\begin{remark}
No-go~\ref{nogo:return-not-period} is the most frequently violated boundary in this
material.  It is why Part~II treats iteration as a syntactic structure
(\S\ref{sec:boundaries}) rather than as a semantic one.
\end{remark}

\subsection{Summary}

\begin{enumerate}
\item A crossing is a cospan with a resolution torsor
  (Definition~\ref{def:through-crossing}, Proposal~\ref{prop:crossing-not-arrow},
  Definition~\ref{def:resolution-torsor}).
\item Monodromy is the difference of two resolutions, and the Legendre monodromy follows
  from the quadratic base change (Theorem~\ref{thm:monodromy-difference},
  Proposition~\ref{prop:legendre-reduction}).
\item A primitive transvection is realised by an exact cellular shear, and vertex
  closure does not imply path closure
  (Proposition~\ref{prop:shear-transvection}, Theorem~\ref{thm:vertex-not-path}).
\item Through forms are fibre products over a singular middle, whose duality is
  relational converse (Definition~\ref{def:through-form},
  Propositions~\ref{prop:order-compat-independent}, \ref{prop:resolution-count}).
\item One feature can be read three ways on one word, and a boundary-return word is still
  not a period (Theorem~\ref{thm:figure-eight},
  Proposition~\ref{prop:double-periods}, No-go~\ref{nogo:return-not-period}).
\end{enumerate}

\section{A logic generated by a satisfaction form}\label{sec:satisfaction-logic}
% Section 9: A logic generated by a satisfaction form
\label{sec:satisfaction-body}

Sections~\ref{sec:triadic-inference}--\ref{sec:through} produced carriers and crossings.
This section asks what kind of \emph{logic} such a carrier supports, and the answer is
constrained from the start: evidence is positive, propositions are sections rather than
subsets, entailment is refinement, and negations are quotients that must be declared.
The section closes by adding a quantitative enrichment and by exposing the trap that
makes it vacuous if used naively.

\subsection{Positive evidence}

\begin{definition}[Satisfaction form]\label{def:satisfaction-form}
Let $\cT_Q,\cX_Q,\cK_Q$ be the three typed carriers of a contract $Q$.  A
\emph{satisfaction form} is a relation
\[
  \Lambda_Q\subseteq \cT_Q\times\cX_Q\times\cK_Q
\]
read as a map into the Sierpinski lattice $\mathbf 2=\{\bot,\top\}$ with
$\bot<\top$, where $\top$ means ``positive evidence of satisfaction'' and $\bot$ means
``no positive evidence''.
\end{definition}

\begin{warning}[$\bot$ is not false]\label{warn:bottom-not-false}
The value $\bot$ in Definition~\ref{def:satisfaction-form} records the absence of
positive evidence.  It does not assert the negation of the proposition.  This is the
same discipline as Warning~\ref{warn:unknown-not-nonexistence}, and it is the reason the
constructions below never form a Boolean complement without declaring one.
\end{warning}

\subsection{Sections are the atoms}

\begin{definition}[Opposite-pair sections]\label{def:opposite-pair-sections}
For each pair of coordinates, define
\[
  A^{t}_{x,k}=\{\,t\in\cT_Q:\Lambda_Q(t,x,k)=\top\,\},\quad
  A^{X}_{t,k}=\{\,x:\Lambda_Q(t,x,k)=\top\,\},\quad
  A^{\cK}_{t,x}=\{\,k:\Lambda_Q(t,x,k)=\top\,\}.
\]
\end{definition}

\begin{proposal}[Generated propositions]\label{prop:generated-propositions}
The admissible propositions of the logic of $Q$ are generated by the sections of
Definition~\ref{def:opposite-pair-sections}, not by the ambient powerset:
\[
  \cL_Q=\langle G_Q\rangle,\qquad\text{and in general}\qquad
  \cL_Q\neq\cP(\cT_Q\times\cX_Q\times\cK_Q).
\]
\end{proposal}

\begin{remark}[Why this is the central design choice]
Taking the powerset as the proposition set is the default of classical semantics and it is
wrong here for a concrete reason: a subset of the carrier may be undefinable by any
section, so admitting it as a proposition silently assumes an expressive power that the
carrier does not have.  Proposal~\ref{prop:generated-propositions} refuses that
assumption and makes the expressive reach of the carrier an explicit object.
\end{remark}

\subsection{Entailment as refinement}

\begin{definition}[Support and preorder]\label{def:support-preorder}
Let $\mathrm{supp}\,P$ be the set of triples at which a generated proposition $P$ carries
value $\top$.  Define
\[
  P\preceq_Q R \iff \mathrm{supp}\,P\subseteq \mathrm{supp}\,R .
\]
\end{definition}

\begin{proposition}[Refinement is an entailment preorder]\label{prop:refinement-preorder}
$(\cL_Q,\preceq_Q)$ is a preorder, and joins exist for the opposite pairs: writing
$P_{I}$ and $P_{-I}$ for the two lifted signs of a central-sign pair,
\[
  P_{\mathrm{proj}}=P_{I}\vee P_{-I}
\]
is exactly the proposition seen by the coarse projective observer.
\end{proposition}

\begin{proof}
Reflexivity and transitivity follow from set inclusion of supports.  The join statement
is Definition~\ref{def:support-preorder} applied to the union of the two supports, which
is precisely the set on which the projective reading forgets the sign.
\end{proof}

\begin{calibration}[Observer quotient with residual]\label{cal:observer-quotient}
On the finite carrier of this section, a three-coordinate observer quotient of the
section lattice is computed together with its exact preimage table.  The residual of the
quotient is therefore known and finite, as
Proposition~\ref{prop:residual-principle} requires.
\end{calibration}

\begin{proposal}[Adjoint quantifiers]\label{prop:adjoint-quantifiers}
The quantifiers appropriate to this logic are adjoints to pullback along an observer map
$q$:
\[
  \exists_q\dashv q^{\ast}\dashv\forall_q ,
\]
where $q^{\ast}$ is the observer quotient of Calibration~\ref{cal:observer-quotient}.
\end{proposal}

\begin{remark}
Proposal~\ref{prop:adjoint-quantifiers} is a \textsf{Proposal}: the adjunction is the
right shape for observer transport, but no theorem here establishes that the finite
lattice of this section is closed under it.  What is established is the data on which
such a theorem would be stated.
\end{remark}

\subsection{The triadicity test}

\begin{proposal}[Falsifiable triadicity]\label{prop:triadicity-test}
Triadicity of $\Lambda_Q$ is falsifiable by ambiguous global fillings: if two distinct
global fillings have equal pairwise sections but different triple support, then the triple
carries information that no pair determines, and the form is triadic in the sense of
Proposal~\ref{prop:coherence-defect}.  If every two fillings with equal pairwise sections
have equal triple support, the form is only pairwise in disguise.
\end{proposal}

\begin{calibration}[The finite fixture fails the test]\label{cal:triadicity-fails}
On the finite satisfaction form constructed in this section the test of
Proposal~\ref{prop:triadicity-test} \emph{fails}: the triple support is determined by the
pairwise sections.  This is the same phenomenon as
Theorem~\ref{thm:defect-decomposes} in a logical setting, and it is the reason
Conjecture~\ref{conj:triadic-crossing} is stated as a conjecture.
\end{calibration}

\subsection{Logic as a learned characteristic}

The conceptual reading of the section is worth stating separately, because it organises
the three notions that the constructions above keep apart.

\begin{proposal}[Three grades of logicality]\label{prop:three-grades}
\begin{enumerate}
\item A \emph{proposition} is a positive (Sierpinski-valued) characteristic.
\item \emph{Entailment} is the first relational characteristic: it judges pairs.
\item \emph{Abstraction} is a coarse invariant together with a declared forgotten fibre
  $R_Q(b)$; a logic earns stability exactly by declaring what it forgets.
\end{enumerate}
\end{proposal}

\begin{warning}[Three different things]\label{warn:bivalence-completeness-decidability}
Bivalence, semantic completeness, and decidability are independent.  A logic may be
two-valued and incomplete, or complete and undecidable, and the passage from positive
evidence to a Boolean one is a \emph{second, declarable quotient} rather than a
consequence of the first.  Nothing in this section licenses forming a complement without
that second declaration.
\end{warning}

\subsection{Energy as an enrichment, and its trap}

A satisfaction form may be enriched by a quantitative measure of how a characteristic
varies along declared directions.

\begin{definition}[Directional energy]\label{def:directional-energy}
On a finite observation carrier with declared move directions $D$, the
\emph{directional Dirichlet energy} of a function $f$ is
\[
  E_D(f)=\sum_{(u,v)\in D}\bigl(f(u)-f(v)\bigr)^{2},
\]
and the \emph{energy vector} is $E_Q(f)=\bigl(E_t(f),E_X(f),E_{\cK}(f)\bigr)$.
\end{definition}

\begin{proposition}[Exact coarse--residual split]\label{prop:energy-split}
Let $q$ be an observer that averages along a forgotten fibre, and write
$\bar f$ for the coarse average and $r=f-q^{\ast}\bar f$ for the zero-mean residual of
each fibre.  Then
\[
  E_Q(f)=E_Q\bigl(q^{\ast}\bar f\bigr)+E_Q(r),
\]
and for the constructional direction in which the fibre is forgotten,
$E_{\cK}(q^{\ast}\bar f)=0$.
\end{proposition}

\begin{proof}
Each summand $\bigl(f(u)-f(v)\bigr)^{2}$ decomposes as the sum of the squares of the
coarse difference and the residual difference plus a cross term that vanishes because $r$
has zero mean on each fibre and the fibres are disjoint.  The second claim is immediate
since $q^{\ast}\bar f$ is constant along the forgotten direction.
\end{proof}

\begin{proposition}[Energy detects invariance]\label{prop:energy-invariance}
For a binary characteristic $\chi_U$,
\[
  E_D(\chi_U)=0 \iff \chi_U \text{ is constant along every } D\text{-move},
\]
and in that case $E_D$ equals the size of the boundary cut of $U$ along $D$.
\end{proposition}

\begin{proof}
Each summand is $(f(u)-f(v))^{2}\in\{0,1\}$ for binary $f$, so the sum vanishes exactly
when no edge crosses the boundary of $U$; otherwise each crossing edge contributes $1$.
\end{proof}

\begin{nogo}[The zero-energy trap]\label{nogo:zero-energy-trap}
The constant characteristic minimizes every directional energy.  Consequently any
optimisation principle of the form ``minimise $E_Q$'' is vacuous unless accompanied by an
explicit nontriviality constraint, and no learning rule may be justified by energy
decrease alone.
\end{nogo}

\begin{proof}
By Proposition~\ref{prop:energy-invariance} a constant characteristic has zero energy,
which is the global minimum since every summand is a square.
\end{proof}

\begin{remark}
No-go~\ref{nogo:zero-energy-trap} is worth stating at length because the error it forbids
is attractive: the energy splits exactly, it is computable, and its minimisers are
therefore easy to find --- too easy.  The exactness of
Proposition~\ref{prop:energy-split} and the vacuity of
No-go~\ref{nogo:zero-energy-trap} are two sides of the same computation.
\end{remark}

\subsection{Summary}

\begin{enumerate}
\item Evidence is positive and $\bot$ is not false
  (Definition~\ref{def:satisfaction-form}, Warning~\ref{warn:bottom-not-false}).
\item Propositions are generated by sections, not by the powerset, and entailment is
  refinement of support (Proposal~\ref{prop:generated-propositions},
  Proposition~\ref{prop:refinement-preorder}).
\item The triadicity test is falsifiable, and the finite fixture of this section fails it
  (Proposal~\ref{prop:triadicity-test}, Calibration~\ref{cal:triadicity-fails}).
\item Energy enriches the picture exactly, and is vacuous without a nontriviality
  constraint (Propositions~\ref{prop:energy-split}, \ref{prop:energy-invariance},
  No-go~\ref{nogo:zero-energy-trap}).
\end{enumerate}

\section{Non-faithfulness: the residual catalogue}\label{sec:residual-catalogue}
% Section 10: Non-faithfulness: the residual catalogue
\label{sec:residual-body}

The previous sections produced no-go results one at a time.  This section collects them,
because together they say something stronger than each does separately: there is a
\emph{systematic} pattern in which quantity fails to determine construction, and naming
the pattern is more useful than proving its instances.  The section also records the
positive counterweights, since a catalogue of failures without a statement of what a
faithful carrier looks like would be misleading.

\subsection{The pattern}

\begin{proposal}[The residual pattern]\label{prop:residual-pattern}
Let $\sigma$ be a quantity computed from a process by an observer $Q$.  In every
instance catalogued below, $\sigma$ is a \emph{partial quotient}: it is invariant under a
group of construction changes that the process itself does not admit.  The correct
response is never to replace $\sigma$ by a finer scalar; it is to retain a residual
recording the orbit of the forgotten changes.
\end{proposal}

\begin{remark}
Proposition~\ref{prop:residual-pattern} is a \textsf{Proposal} because "partial quotient"
is a description, not a hypothesis with a theorem attached.  Its content is
methodological and its justification is the list: in each case below, the residual is
exactly what repairs the deficiency, and no case is repaired by refining the scalar.
\end{remark}

\subsection{Spectral shadows}

\begin{theorem}[Spectrum does not determine history]\label{thm:spectrum-history}
Let $x$ be a single input.  The programs $\mathrm{add}(\mathrm{copy}(x))$ and
$\mathrm{scale}_{2}(x)$ agree on realized value, on derivative, on all iterates, and on
the action spectrum, which is $\{2\}$ in both cases; they differ in the construction
graph, in the occurrence structure, and in the source history.  Hence no observer that
sees only values, derivatives, iterates, and spectra can separate them.
\end{theorem}

\begin{proof}
Both programs induce the same endomorphism $x\mapsto 2x$ on the declared state, so every
functional of the induced action agrees.  They differ in the number of occurrences of
the source and in the recorded lineage, which are not functions of the induced action.
\end{proof}

\begin{proposition}[The affine carrier is not diagonalisable into a spectrum]
\label{prop:affine-spectrum}
In the affine carrier, scaling acts diagonally with spectrum
$\{m^{-1},\dots,m^{-K}\}$, while addition acts unipotently and is nilpotent on the
relevant graded pieces.  Consequently a bare eigenvalue set cannot encode addition: the
characteristic polynomial of a unipotent operator is $(\lambda-1)^{d}$ regardless of the
translation amount.
\end{proposition}

\begin{proof}
Scaling conjugates $A\mapsto M AM^{-1}$ and stretches the translation grid by $m$ at each
graded level, giving the stated spectrum.  Addition is the exponential of a nilpotent
translation, whose characteristic polynomial carries no information about the amount.
\end{proof}

\subsection{Two faithful carriers, for contrast}

\begin{theorem}[Marked second-jet faithfulness]\label{thm:second-jet}
The positive affine fragment embeds faithfully into the marked second jet at the
projective point at infinity,
\[
  m=c_1^{-1},\qquad a=-c_2c_1^{-2},
\]
where $(c_1,c_2)$ are the marked jet coefficients.  The faithfulness is a statement about
\emph{marked} jets: the unmarked jet does not suffice, since it is defined only up to the
chart that the marking fixes.
\end{theorem}

\begin{proof}
Composition of affine maps corresponds to multiplication in
$\R[[q]]/(q^{3})$ on the marked coefficient pair; solving for $(m,a)$ gives the displayed
formulas, and the correspondence is a bijection onto the positive affine fragment within
the marked class.
\end{proof}

\begin{theorem}[Co-generation of polynomial and matrix readings]\label{thm:cogeneration}
Transport is homogeneous of degree one in declared symbolic probe generators, and the
extracted coefficient table equals the transpose of the forward Jacobian.  Hence the
polynomial-like and matrix-like presentations of one process are co-generated, rather
than independently imposed.
\end{theorem}

\begin{proof}[Status]
Proved for the declared probe class by linearity of the demand section.  The phrase
"after a chart" is essential: the identification of the coefficient table with the
transpose requires the port chart to be declared, and without it the two readings are
merely isomorphic.
\end{proof}

\begin{remark}[The moral of the contrast]
Theorems~\ref{thm:second-jet} and~\ref{thm:cogeneration} are the positive half of the
catalogue.  Both work by making the marking or the chart explicit.  The failures below
are exactly the cases where a marking or chart was silently omitted, and the residual
records what the omission cost.
\end{remark}

\subsection{The observation tower is strict}

\begin{theorem}[Strict observation tower]\label{thm:observation-tower}
The three observation levels
\[
  \cO_{\mathrm{prog}}\ \supsetneq\ \cO_{\mathrm{or}}\ \supsetneq\ \cO_{\mathrm{proj}}
\]
are strictly separated: projective spectral data is strictly coarser than oriented action
data, which is strictly coarser than program-aware data.  Each strictness is witnessed by
a pair of programs identified at the lower level and separated at the higher.
\end{theorem}

\begin{proof}
The projective level forgets the central sign
(Proposition~\ref{prop:order-two-four}), and the oriented level forgets occurrence
structure (Theorem~\ref{thm:spectrum-history}).  Both forgettings are non-injective on the
finitary fixtures exhibited there.
\end{proof}

\subsection{Objectification and effective completion}

\begin{nogo}[Objectification is not natural]\label{nogo:objectification}
Let $\mathrm{Obj}$ send a finite cut presentation to its simplest surreal value.  Then
$\mathrm{Obj}$ is non-injective, and it is not natural:
\[
  g\bigl(\mathrm{Obj}\,C\bigr)\neq \mathrm{Obj}\bigl(gC\bigr)
  \qquad\text{for } g(z)=3z \text{ and a suitable } C .
\]
Consequently no analysis factoring through the objectified value admits a left-inverse
synthesis.
\end{nogo}

\begin{proof}
Non-injectivity is immediate, since distinct presentations may share a simplest value.
For non-naturality, the two sides of the displayed inequality are computed on an explicit
finite $C$; the objectified value is simplified before scaling on the left and after
scaling on the right, and simplicity is not preserved by scaling.
\end{proof}

\begin{proposition}[Effective one-sidedness]\label{prop:effective-onesided}
For a universal prefix-free machine $U$, let
\[
  \Omega_U=\sum_{U(p)\downarrow}2^{-|p|}\le 1
\]
be its halting mass, and let $(L_\Omega\mid R_\Omega)$ be the associated cut.  Then
$L_\Omega$ is computably enumerable while $R_\Omega$ is not.  Hence a computational cut
is intrinsically asymmetric.
\end{proposition}

\begin{proof}[Status]
The Kraft bound is a standard measure-theoretic computation on the cylinder sets of a
prefix-free set.  The asymmetry is the standard halting asymmetry: membership on the left
is witnessed by a halting computation, while membership on the right would decide
non-halting.
\end{proof}

\begin{proposition}[Mass cannot repair the fibre]\label{prop:mass-not-fibre}
The objectification fibre mass
\[
  m_{U,Q}(x)=\sum_{\mathrm{Obj}_Q(U(p))=x}2^{-|p|}
\]
is a lower-semicomputable semimeasure.  It quantifies the fibre but cannot recover its
members.
\end{proposition}

\begin{proof}
Lower semicomputability follows from the definition as a sum over enumerated halting
computations.  Recovery fails because a sum over a set determines neither its cardinality
nor its elements.
\end{proof}

\begin{proposition}[No total finite coding of $\No$]\label{prop:no-finite-coding}
There is no set-sized family of finite codes surjecting onto the surreal numbers.
\end{proposition}

\begin{proof}
A set-sized family of finite codes is a set, and $\No$ is a proper class; no set surjects
onto a proper class in the relevant setting.
\end{proof}

\subsection{Snapshots, scopes, and traces}

\begin{proposition}[Equal snapshots, nonempty interval]\label{prop:snapshot-not-empty}
For exact program slices between nested cuts, equal frontier snapshots do not imply an
empty interval: an explicit discard witnesses two distinct slices with the same snapshot.
\end{proposition}

\begin{proposition}[Zero-event incidence is not single-valued]\label{prop:zero-event-fibre}
A zero-event graft frame is compatible with two distinct identity slices fixed by an
independent surgery.  Hence the available frame and slice data admit no canonical
single-valued placement of scope in cut; the fibre has two elements.
\end{proposition}

\begin{proposition}[Trace arithmetic does not identify the sides]\label{prop:trace-sides}
For the three-sided trace calibration, the temporal and spatial residuals vanish while
the constructional residue $(+1,-1,0,0,0)$ does not.  Consequently the two sides are not
semantically identified by their trace arithmetic.
\end{proposition}

\begin{proposition}[Count factorization and the recovery obstruction]
\label{prop:count-factorization}
Let $T$ be the trace count and $n$ the construction length.  Then
\[
  T(p)=\bigl(n,\ n-1,\ 3(n-1)\bigr),
\]
a positive factorization through construction length.  However the projection
$(T,S)\mapsto C$ is not single-valued: an explicit pair of constructions shares $(T,S)$
and differs in $C$.
\end{proposition}

\begin{proof}[Status]
The factorization is an exact finite count.  The obstruction is an explicit pair, checked
on the recorded data; it is a counterexample to single-valued recovery, not a general
classification of the fibres of $(T,S)\mapsto C$.
\end{proof}

\begin{proposition}[Period covectors do not determine history]\label{prop:period-not-history}
A period covector is a functional of a homology class and therefore constant on the
construction paths that realise that class.  It cannot reconstruct the construction
history.
\end{proposition}

\subsection{Minimal expressive repair}

One result in the catalogue deserves separate billing, because it gives an exact
\emph{minimum} rather than only an obstruction.

\begin{theorem}[Order-embedding obstruction]\label{thm:order-obstruction}
Let a coarse observation erase an incomparability, i.e.\ two elements become
indistinguishable in the observation order.  Then no single added scalar coordinate can
restore the lost order, whereas two monotone Boolean coordinates can, and two is the
minimum.
\end{theorem}

\begin{proof}
A single scalar added to a coarse observation produces a total preorder on the fibre,
since scalars are totally ordered; a total preorder cannot contain an incomparability.  Two
monotone Boolean coordinates separate the four possible cells, which suffices to embed the
erased relation, and one coordinate gives only two cells.  Hence two is both sufficient
and necessary.
\end{proof}

\begin{remark}[Why this is the catalogue's high point]
Theorem~\ref{thm:order-obstruction} converts "coarse observation loses a relation" from a
complaint into a quantitative statement with an exact constant.  It is the model for what
the rest of the catalogue would need in order to become equally informative.
\end{remark}

\subsection{Scale and winding}

\begin{proposition}[Scale forgetting]\label{prop:scale-forgetting}
Let $\lambda=\varepsilon Lz$ with $L=(2I+C)/3$ and $C$ the cyclic permutation matrix,
$C^{3}=I$.  Then $L$ is invertible over $\Q$, with
\[
  (2I+C)\bigl(4I-2C+C^{2}\bigr)=9I ,
\]
so the scale-marked triple is recoverable provided $\varepsilon\neq0$.  If the scale
marker is discarded, the pairs $(\varepsilon,z)$ and $(\varepsilon/2,2z)$ share $\lambda$,
and recovery fails.
\end{proposition}

\begin{proof}
Expanding $(2I+C)(4I-2C+C^{2})=8I-4C+2C^{2}+4C-2C^{2}+C^{3}=9I$ using $C^{3}=I$ gives the
inverse, hence invertibility.  The collision is immediate from
$\lambda=(\varepsilon/2)L(2z)$.
\end{proof}

\begin{nogo}[Winding counts do not determine the global action]\label{nogo:winding}
Abelian winding counts do not determine the global action on
$\PP^{1}\setminus\{0,1,\infty\}$: there are words of total winding zero whose commutator is
nontrivial,
\[
  ABA^{-1}B^{-1}=\begin{pmatrix}13&8\\ 8&5\end{pmatrix}\neq I .
\]
Moreover endpoint closure does not imply identity of the transported homology: a diamond
path of winding $+1$ transports $(0,1)$ to $(2,1)$.
\end{nogo}

\begin{proof}
Two computations on explicit words.  The first exhibits a nontrivial commutator with zero
abelian shadow, so the abelianisation is not faithful.  The second exhibits a closed
endpoint with nontrivial transport.
\end{proof}

\subsection{The catalogue}

\begin{center}
\small
\begin{longtable}{@{}p{0.30\textwidth}p{0.56\textwidth}@{}}
\toprule
Shadow & What it fails to determine \\
\midrule
\endhead
value, derivative, iterates, spectrum & construction graph, occurrence, source history
  (Thm~\ref{thm:spectrum-history}) \\
eigenvalue set & the translation amount of a unipotent factor
  (Prop~\ref{prop:affine-spectrum}) \\
projective spectral data & the central sign; oriented action
  (Thm~\ref{thm:observation-tower}) \\
oriented action data & occurrence structure
  (Thm~\ref{thm:observation-tower}) \\
objectified surreal value & the presentation, and it is not natural
  (No-go~\ref{nogo:objectification}) \\
fibre mass of an objectification & the members of the fibre
  (Prop~\ref{prop:mass-not-fibre}) \\
frontier snapshot & emptiness of an interval
  (Prop~\ref{prop:snapshot-not-empty}) \\
frame and slice data & a single-valued scope placement
  (Prop~\ref{prop:zero-event-fibre}) \\
trace arithmetic & identification of the two sides
  (Prop~\ref{prop:trace-sides}) \\
count pair $(T,S)$ & the construction $C$
  (Prop~\ref{prop:count-factorization}) \\
period covector & construction history
  (Prop~\ref{prop:period-not-history}) \\
discarded scale marker & the marked triple
  (Prop~\ref{prop:scale-forgetting}) \\
abelian winding & the global action
  (No-go~\ref{nogo:winding}) \\
boundary-return word & a periodic orbit
  (No-go~\ref{nogo:return-not-period}) \\
directional energy & any nontrivial characteristic
  (No-go~\ref{nogo:zero-energy-trap}) \\
\bottomrule
\end{longtable}
\end{center}

\subsection{What the catalogue implies}

\begin{enumerate}
\item The failures are not ad hoc.  In each row the shadow is invariant under changes the
  process does not admit, and the repair is a residual rather than a finer scalar
  (Proposal~\ref{prop:residual-pattern}).
\item Where a residual is supplied, the theory works: marked jets and declared port charts
  give genuine faithfulness (Theorems~\ref{thm:second-jet}, \ref{thm:cogeneration}).
\item At least one obstruction has an exact minimum
  (Theorem~\ref{thm:order-obstruction}), which is the standard the rest should meet.
\item Two of the rows are about \emph{absence of structure} rather than loss of
  information: No-go~\ref{nogo:return-not-period} (no feedback operator) and
  No-go~\ref{nogo:zero-energy-trap} (no nontriviality constraint).  These are the
  invitations to Part~II, where the missing structures are supplied syntactically.
\end{enumerate}


\part{Threaded apertures and the six-port kernel}

\section{Aperture-first semantics: the typed hole}\label{sec:apertures}
% Section 11: Aperture-first semantics: the typed hole
\label{sec:apertures-body}

Part~I repeatedly met a place where a reading was \emph{undetermined}: a missing
feedback operator (No-go~\ref{nogo:return-not-period}), a fibre whose members a mass could
not recover (Proposition~\ref{prop:mass-not-fibre}), a logical value $\bot$ that meant no
evidence rather than falsehood (Warning~\ref{warn:bottom-not-false}).  In each case the
honest description of the situation was not a value but a \emph{place where something
belongs and is not yet there}.

Part~II takes that place as its primitive.  This section introduces it, fixes its
discipline, and shows what a finite semantics built on it looks like.

\subsection{Why the hole is primitive}

\begin{proposal}[Aperture-first]\label{prop:aperture-first}
The primitive object of a partial semantics should be the typed hole --- a position of
declared type whose content is not determined --- rather than a state plus an
exception, or a value plus a flag.  The reasons are that holes compose, that they carry
their type, and that the operations which fill and empty them are not forced to be
inverses.
\end{proposal}

\begin{definition}[Typed aperture]\label{def:aperture}
Let $D$ be a finite set of domain slots and let $\tau$ assign a type to each slot.  A
\emph{typed aperture at $d\in D$} is a syntactic object $h_d$ satisfying:
\begin{enumerate}
\item $h_d$ is syntactically empty --- it contributes no term of its own;
\item $h_d$ has a declared type $\tau(d)$;
\item semantically, $h_d$ denotes the \emph{filling family} of all admissible inhabitants
  of $\tau(d)$ compatible with the surrounding context.
\end{enumerate}
\end{definition}

The third clause is the substantive one: an aperture is not a wildcard.  It is a family,
and the family is part of the semantics.

\begin{proposal}[Open and close are not inverses]\label{prop:open-close}
The operations that open an aperture and that close it are deliberately not inverse.
Closing selects an inhabitant and thereby \emph{records} a choice; opening again does not
erase the record but adds a new one, and the difference is exactly the residual of
Definition~\ref{def:characteristic}.
\end{proposal}

\begin{remark}
Proposal~\ref{prop:open-close} is the syntactic form of
Proposition~\ref{prop:lossy-quotient}: closing is a quotient, and the record is the
residual.  Treating open and close as inverses would amount to asserting that the
quotient has a canonical section, which \S\ref{sec:observers} showed it does not.
\end{remark}

\subsection{The four structural statuses}

\begin{definition}[Aperture status]\label{def:aperture-status}
At any stage an aperture is in exactly one of four structural statuses:
\begin{enumerate}
\item \emph{no aperture} --- the position is not a hole in this context;
\item \emph{inhabited} --- the filling family is nonempty and a witness is available;
\item \emph{exhaustively empty} --- the filling family is empty \emph{and} the
  emptiness is certified by an exhaustion argument;
\item \emph{unresolved} --- neither a witness nor an exhaustion is available within the
  declared budget.
\end{enumerate}
\end{definition}

\begin{warning}[The third and fourth statuses are different]\label{warn:empty-unresolved}
Status (3) and status (4) are the aperture-level form of
Warning~\ref{warn:unknown-not-nonexistence}, and conflating them is the most common error
in this material.  An exhausted search that finds no inhabitant is a positive result; a
search that runs out of budget is not a result at all.
\end{warning}

\subsection{The finite terminal-world carrier}

The semantics of Part~II is finite from the outset.  Its carrier is built from the
possible outcomes of a bounded process.

\begin{definition}[Terminal-world carrier]\label{def:terminal-worlds}
Let $D$ be a finite set of domain slots.  A \emph{terminal world} is a complete assignment
of values to $D$.  The \emph{terminal-world carrier} $H_D$ is the set of terminal worlds
that are \emph{nonempty}, meaning that the assignment is realised by at least one
admissible execution.  In the calibrated instance with three slots the seven nonempty
faces give $|H_7|=7$.
\end{definition}

\begin{proposition}[Exhaustive partition of supports]\label{prop:seven-faces}
For the seven nonempty faces of a triangle carrier, the $2^{7}=128$ supports are
exhaustively partitioned by size into
\[
  128=2+6+30+90 ,
\]
with the $30$ refined to $30=18+12$, and classified into $40$ orbits under $D_3$ and
$20$ archetypes under $D_3\times C_2^{\neg}$.
\end{proposition}

\begin{proof}[Status]
An exhaustive finite enumeration of the declared fixture, reported as a calibration.  The
refinement of $30$ into $18+12$ separates the chiral classes from the achiral ones.
\end{proof}

\begin{nogo}[Mask data does not recover direction]\label{nogo:mask-direction}
The mask enumeration of Proposition~\ref{prop:seven-faces} carries no construction
direction: two constructions differing in orientation and agreeing on the mask are not
separated by any mask-level datum.
\end{nogo}

\begin{proof}
The mask records which faces are nonempty, and the datum is invariant under the
orientation-reversing involution of the carrier; a construction history is not.
\end{proof}

\subsection{Fibre semantics}

\begin{definition}[Filling fibre]\label{def:filling-fibre}
For a formula $\varphi$, a model carrier $\cM$, and an assignment $\rho$ to the free
slots, the \emph{filling family} is
\[
  \Fill_{\cM,\rho}(\varphi)\subseteq\cM
\]
the set of terminal worlds in which $\varphi$ is realised.  The three semantic values are
then read off structurally:
\begin{itemize}
\item \emph{truth} is an inhabiting line: $\Fill_{\cM,\rho}(\varphi)\neq\varnothing$;
\item \emph{falsehood} is an exhaustively empty fibre, certified;
\item \emph{unresolved} is neither, within the budget.
\end{itemize}
\end{definition}

\begin{proposal}[Connectives act on fibres]\label{prop:fibre-connectives}
Connectives are operations on filling families rather than on truth values:
conjunction is intersection of families, disjunction is union, and the quantifiers are
the adjoints of Proposition~\ref{prop:adjoint-quantifiers}.  Consequently the semantics is
determined by what it does to families, and the truth values are derived readings.
\end{proposal}

\begin{remark}
Proposition~\ref{prop:fibre-connectives} is where Part~II pays for Part~I's discipline.  If
truth is primitive, falsehood and unresolved collapse.  If the family is primitive, they
do not, and the four statuses of Definition~\ref{def:aperture-status} are the honest
taxonomy.
\end{remark}

\subsection{The geometric reading, and its conditionality}

A finite observer can be presented geometrically, but only conditionally.

\begin{proposal}[Finite observer as a marked surface]\label{prop:surface-observer}
A finite observer may be presented as a marked finite-type surface
\[
  \Sigma^{\circ}_{P}\simeq \Gamma_{P}\backslash\bH^{2}_{\mathrm{AM}} ,
\]
whose sealed threads are the vocabulary generators, with a cap--seal calculus and a
$J$-operator acting on thread ends.  In this presentation universality is an
\emph{axiom proposal}: the observer stores finitely many generators, relations, markings
and residuals plus a rule producing any demanded finite portion, so infinity is potential
rather than completed.
\end{proposal}

\begin{warning}[Conditionality of the geometric reading]\label{warn:geometric-conditional}
Proposition~\ref{prop:surface-observer} is a representation claim, not an identification.
The passage between the symbolic grammar and the hyperbolic carrier requires an explicit
coding or quasi-isometry theorem, and none is supplied here.  Statements about the
hyperbolic model in this paper are imported from Paper~0 \cite{Yuan2026AEGPaperZero} and
are not re-derived.
\end{warning}

\subsection{Summary}

\begin{enumerate}
\item The typed hole is primitive, it denotes a filling family, and open and close are not
  inverses (Definition~\ref{def:aperture}, Proposal~\ref{prop:open-close}).
\item Four statuses are needed, and exhaustion differs from unresolved
  (Definition~\ref{def:aperture-status}, Warning~\ref{warn:empty-unresolved}).
\item Semantics is fibre-valued, with truth, certified falsehood, and unresolved kept
  apart (Definition~\ref{def:filling-fibre}, Proposal~\ref{prop:fibre-connectives}).
\item The geometric reading is available but conditional
  (Proposal~\ref{prop:surface-observer}, Warning~\ref{warn:geometric-conditional}).
\end{enumerate}

\section{Finite adequacy of threaded logic}\label{sec:adequacy}
% Section 12: Finite adequacy of threaded logic
\label{sec:adequacy-body}

Definition~\ref{def:filling-fibre} gives a semantics.  This section proves that between
that semantics and a threaded derivability relation there is an exact adequacy, and it
states the hypotheses on which the equivalence depends.  The hypotheses are the content:
each of them is a place where an unrestricted claim would fail.

\subsection{The judgment}

\begin{definition}[Threaded judgment]\label{def:threaded-judgment}
A \emph{threaded judgment} is a sequent whose antecedents are \emph{threaded}: each
antecedent is attached to a thread, and threads are consumed by the rules rather than
shared.  Writing $\Gamma$ for a threaded context and $\varphi$ for a formula, we write
$\Gamma\vdash\varphi$ for derivability in the threaded calculus.
\end{definition}

The essential point is linearity of threads: no rule duplicates or discards an antecedent
thread.  Copy and discard, where they occur, are explicit connectives rather than
structural rules.  This is the syntactic shadow of
Proposition~\ref{prop:word-to-graph}: the fan-out must be written down.

\begin{definition}[Realisation]\label{def:realisation}
Given a carrier $\cM$, an assignment $\rho$, and a threaded context $\Gamma$, a
\emph{realisation} of $\Gamma$ is a choice, for each thread, of a terminal world
compatible with the thread's formulas.  Write $\cM,\rho\models\Gamma$ when a realisation
exists, and $\cM,\rho\models\varphi$ when
$\Fill_{\cM,\rho}(\varphi)\neq\varnothing$.
\end{definition}

\subsection{The four hypotheses}

\begin{definition}[Adequacy hypotheses]\label{def:adequacy-hypotheses}
The adequacy theorem below is stated under four hypotheses on the carrier and the
connectives.
\begin{enumerate}
\item \textbf{Cartesian.}  The carrier is a product of independent slot domains, so that a
  partial assignment extends to a terminal world whenever no constraint is violated.
\item \textbf{Compatibility.}  The constraints attached to a threaded context are
  pair-wise compatible exactly when their conjunction is nonempty.
\item \textbf{Transformer.}  Each connective has an associated transformer on filling
  families, and the transformer is sound and complete for the connective's semantics.
\item \textbf{Exhaustiveness.}  For each atomic formula and each partial assignment, either
  a witness is produced or the emptiness of the family is certified.
\end{enumerate}
\end{definition}

Each hypothesis rules out a genuine failure mode.  Without (1) a partial assignment need
not extend; without (2) pairwise nonemptiness need not imply joint nonemptiness --- which
is exactly the obstruction of Proposition~\ref{prop:triadicity-test}; without (3) a
connective has no semantics on families; without (4) the fourth and third statuses of
Definition~\ref{def:aperture-status} collapse.

\subsection{The adequacy theorem}

\begin{theorem}[Threaded adequacy]\label{thm:adequacy}
Assume the four hypotheses of Definition~\ref{def:adequacy-hypotheses}.  Then for every
threaded context $\Gamma$ and formula $\varphi$,
\[
  \Gamma\vdash\varphi
  \quad\Longleftrightarrow\quad
  \text{for all } \cM,\rho:\ \cM,\rho\models\Gamma \Rightarrow
  \cM,\rho\models\varphi .
\]
Equivalently, in the fibre reading,
\[
  \cM,\rho\models\varphi
  \iff
  \Fill_{\cM,\rho}(\varphi)\neq\varnothing .
\]
\end{theorem}

\begin{proof}
By structural induction on derivations and on formulas.

\emph{Soundness direction.}  Suppose $\Gamma\vdash\varphi$ with a derivation $\pi$.  We show
each rule preserves realisation.  For an assumption, the conclusion's thread is among the
realised threads by Definition~\ref{def:realisation}.  For a conjunction rule, hypothesis
(3) makes the transformer of $\wedge$ exactly intersection of families, so a realisation
of both premises realises the conclusion.  For a linear implication rule, hypothesis (1)
allows the discharged thread to be replaced by the witness supplied for the premise, and
the conclusion's family is nonempty because a witness exists for the implication itself.
For quantifier rules, hypothesis (1) lets a witness for a free slot be chosen, and
hypothesis (2) guarantees that the pairwise-compatible constraints are jointly
satisfiable, which is what licenses the introduction rule.

\emph{Completeness direction.}  Suppose every realisation of $\Gamma$ realises $\varphi$.
Take the canonical assignment $\rho^{\ast}$ built from the family data: for each slot
choose, by hypothesis (4), either a witness or a certification of emptiness.  If the
family of $\varphi$ were empty under every $\rho$ realising $\Gamma$, hypothesis (4) would
provide a certificate of emptiness; that certificate builds a derivation of the negated
sequent.  Since by assumption no realisation falsifies $\varphi$, the emptiness certificate
cannot be produced for any realising $\rho$, so a derivation of
$\Gamma\vdash\varphi$ exists by induction on the formula's construction.
\end{proof}

\begin{calibration}[The theorem without the hypotheses]\label{cal:adequacy-fails}
For each of the four hypotheses there is a finite fixture on which the equivalence of
Theorem~\ref{thm:adequacy} fails while the other three hold.  The fixtures are the
calibrations of this section; they are the reason the hypotheses are stated rather than
assumed silently.
\end{calibration}

\subsection{Substitution}

\begin{theorem}[Fibre-natural substitution]\label{thm:substitution-natural}
Under the hypotheses of Definition~\ref{def:adequacy-hypotheses}, substitution of a
derivable formula for a free slot commutes with the filling-family construction:
\[
  \Fill_{\cM,\rho}\bigl(\varphi[\psi/x]\bigr)
  =
  \Fill_{\cM,\rho[x\mapsto \Fill_{\cM,\rho}(\psi)]}(\varphi).
\]
\end{theorem}

\begin{proof}
Induction on $\varphi$.  At an atom the two sides are the same family by definition of the
extended assignment.  At a connective, hypothesis (3) makes the transformer act on families,
and the induction hypothesis identifies the substituted families on each side; the
transformer's functoriality then gives the identity at the composite.
\end{proof}

\begin{remark}
Theorem~\ref{thm:substitution-natural} is the semantic counterpart of the syntactic
substitution theorem of \S\ref{sec:normalization}.  The two are not the same statement:
the present one is about families, the later one is about ledgers.
\end{remark}

\subsection{Promotion boundary}

\begin{calibration}[What this section does not license]\label{cal:adequacy-boundary}
Theorem~\ref{thm:adequacy} is a finite semantic theorem about a declared carrier and a
declared connective set.  It licenses no claim of the following forms: that the calculus is
complete for an unrestricted semantics; that the four hypotheses are necessary; that
quantifiers behave classically; or that the carrier's finiteness is an artefact of
presentation.  In particular the theorem is stated for Boolean semantics on $2^{D}$, hence
in particular on $H_D$, and says nothing about carriers with dependent slot domains.
\end{calibration}

\subsection{Summary}

\begin{enumerate}
\item Threaded judgments consume threads linearly, so fan-out is explicit
  (Definition~\ref{def:threaded-judgment}).
\item Adequacy holds under four named hypotheses, each of which blocks a genuine failure
  (Definition~\ref{def:adequacy-hypotheses}, Theorem~\ref{thm:adequacy}).
\item Substitution is fibre-natural (Theorem~\ref{thm:substitution-natural}).
\item The theorem's scope is finite and declared
  (Calibration~\ref{cal:adequacy-boundary}).
\end{enumerate}

\section{Ordered linear deduction: soundness and its limit}\label{sec:ordered-linear}
% Section 13: Ordered linear deduction: soundness and its limit
\label{sec:ordered-linear-body}

\S\ref{sec:adequacy} gave a semantics and a calculus adequate for it under hypotheses.
This section fixes the specific fragment studied in the rest of Part~II and asks the two
questions that matter about it: is it sound, and is it complete?  The answers are yes and
no, and the failure of completeness is exact rather than mysterious: it is measured by a
grade invariant.

\subsection{The fragment}

\begin{definition}[The ordered linear fragment $\mathrm{TND}_0$]
\label{def:tnd0}
$\mathrm{TND}_0$ is the implicational fragment of threaded natural deduction in which
\begin{enumerate}
\item antecedents form an \emph{ordered} context: order is part of the judgment and there
  is no exchange;
\item there is no weakening and no contraction; copy and discard, if present, are
  explicit terms;
\item a judgment carries a \emph{ledger} recording which antecedents each conclusion
  consumed.
\end{enumerate}
\end{definition}

\begin{definition}[Ledger-indexed judgment]\label{def:ledger-judgment}
Write
\[
  L\mid \Gamma \vdash \varphi
\]
for the judgment that $\varphi$ is derivable from the ordered context $\Gamma$ while
consuming exactly the ledger $L$.  The ledger is not notation for the context: two
judgments may share $\Gamma$ and differ in $L$.
\end{definition}

The distinction between $\Gamma$ and $L$ is the syntactic analogue of the distinction
between source and occurrence that Part~I insisted on
(Proposition~\ref{prop:word-to-graph}).  A context lists what is available; a ledger
records what was used, and by which use.

\subsection{Soundness}

\begin{theorem}[Soundness of $\mathrm{TND}_0$]\label{thm:soundness}
If $L\mid\Gamma\vdash\varphi$ then for every carrier $\cM$ and assignment $\rho$,
\[
  \cM,\rho\models\Gamma \ \Rightarrow\ \cM,\rho\models\varphi .
\]
In particular $\mathrm{TND}_0$ is sound for Boolean semantics on $2^{D}$ and hence on the
terminal-world carrier $H_D$ of Definition~\ref{def:terminal-worlds}.
\end{theorem}

\begin{proof}
Induction on the derivation, using only that the fragment's rules preserve realisation of
the ordered context: the implication-introduction rule discharges the head of the ordered
context, implication-elimination consumes it, and no rule moves antecedents.  The argument
is the ordered specialisation of Theorem~\ref{thm:adequacy} and does not need hypotheses
(1)--(3) beyond the product structure of $2^{D}$.
\end{proof}

\subsection{The grade invariant, and incompleteness}

\begin{definition}[Grade]\label{def:grade}
For a formula $\varphi$ and a terminal world assignment, let $\mathrm{gr}(\varphi)$ be the
alternation count of the leading connective along the leftmost spine, read as a natural
number.  For a derivation, $\mathrm{gr}$ is the maximum over its judgments.
\end{definition}

\begin{theorem}[Incompleteness]\label{thm:incompleteness}
$\mathrm{TND}_0$ is \emph{not} complete for Boolean semantics on $2^{D}$: there exist
$\Gamma$ and $\varphi$ with $\cM,\rho\models\Gamma\Rightarrow\cM,\rho\models\varphi$ for
all $\rho$, yet $L\mid\Gamma\vdash\varphi$ is underivable for every ledger $L$.
\end{theorem}

\begin{proof}
Semantic entailment on $2^{D}$ is closed under the classical structural rules.  Take a
semantically valid implication whose witnessing term requires a permutation of antecedents
(exchange) or the reuse of one antecedent (contraction).  By
Definition~\ref{def:tnd0} neither is available, and by
Definition~\ref{def:ledger-judgment} a ledger cannot be shared between two uses.  Hence no
derivation exists, although the entailment holds.
\end{proof}

\begin{proposition}[The grade measures the gap]\label{prop:grade-measures}
The mismatch between Boolean entailment and threaded derivability is measured by the grade
of Definition~\ref{def:grade}: for a fixed carrier, the minimal grade at which a valid
entailment becomes derivable is a computable function of the formula, and it is unbounded
over formulas.
\end{proposition}

\begin{proof}[Status]
The computation is finite on each formula by the decision procedure of
\S\ref{sec:confluence-decidability}; unboundedness is exhibited by the family of formulas
requiring $n$ nested exchanges.
\end{proof}

\begin{remark}[Why incompleteness is a result, not a defect]
Theorem~\ref{thm:incompleteness} tells us that the ordered linear fragment \emph{selects}
a logic strictly finer than Boolean entailment.  On the reading of
Proposal~\ref{prop:three-grades} that is information about which logic the grammar
encodes, not a bug to be repaired by adding structural rules.
\end{remark}

\subsection{Triadic conflict apertures}

The fragment has one genuinely triadic rule, and it is worth isolating.

\begin{proposal}[Conflict aperture]\label{prop:conflict-aperture}
Suppose two lines of derivation carry distinct tri-domain identities and compete for one
linear source.  Then the admissible aperture is uniquely determined by the complement of
the two occupied domains:
\[
  \mu(d,e)=D\setminus\{d,e\},
\]
and a triangular cell witnessing the conjunction $\wedge_{\triangle}$ exists exactly when a
certified linear common lift of the two lines is available.
\end{proposal}

\begin{proposition}[Only a certified common lift closes the cell]
\label{prop:common-lift}
If the two lines have no certified linear common lift, then the triangular cell of
Proposal~\ref{prop:conflict-aperture} admits no filler in $\mathrm{TND}_0$, and the
aperture remains in the \emph{unresolved} status of
Definition~\ref{def:aperture-status}.
\end{proposition}

\begin{proof}
Opening the aperture consumes the source once, by linearity; closing it requires two
inhabitants of the complementary domain that agree on the source.  Without a certified
common lift the two inhabitants are not related by any derivation of the fragment, so no
filler exists.  Since the search space is finite and the fragment is decidable
(\S\ref{sec:confluence-decidability}), failure is reported as unresolved rather than as
falsity, in accordance with Warning~\ref{warn:empty-unresolved}.
\end{proof}

\begin{remark}
The word ``triadic'' here is doing precise work.  The three domains $d,e,\mu(d,e)$ are not
three independent coordinates: once two are occupied the third is determined.  This is the
syntactic counterpart of Proposition~\ref{prop:three-cusps}, and, like it, it does not by
itself produce a non-decomposable triple defect
(Calibration~\ref{cal:triadicity-fails}).
\end{remark}

\subsection{Summary}

\begin{enumerate}
\item $\mathrm{TND}_0$ is ordered, linear, and ledger-indexed
  (Definitions~\ref{def:tnd0}, \ref{def:ledger-judgment}).
\item It is sound for Boolean semantics on $2^{D}$ and hence on $H_D$
  (Theorem~\ref{thm:soundness}).
\item It is provably incomplete, and the gap is measured by a grade
  (Theorem~\ref{thm:incompleteness}, Proposition~\ref{prop:grade-measures}).
\item Its one triadic rule yields a uniquely determined aperture whose cell closes exactly
  under a certified common lift
  (Proposal~\ref{prop:conflict-aperture}, Proposition~\ref{prop:common-lift}).
\end{enumerate}

\section{Substitution, ledger transport, and strong normalization}\label{sec:normalization}
% Section 14: Substitution, ledger transport, and strong normalization
\label{sec:normalization-body}

The proof theory of $\mathrm{TND}_0$ begins with substitution, because an ordered linear
fragment cannot borrow the usual substitution lemma: the order is part of the data, and the
ledger records uses that an ordinary substitution would silently move.  This section proves
an ordered one-hole substitution theorem with an exact ledger splice, derives conditional
cut admissibility from it, shows that raw ledger equality is the wrong target, and then
transports the whole apparatus through proof contexts to obtain strong normalization.

\subsection{Ordered one-hole substitution}

\begin{definition}[One-hole context]\label{def:one-hole-context}
A \emph{one-hole proof context} is a proof tree with a distinguished hole at a leaf.  Two
proof trees are \emph{concatenable} along the hole when the type of the conclusion at the
hole matches the type of the root of the inserted tree; concatenation is order-respecting
and is written $Q[\cdot]$.
\end{definition}

\begin{definition}[Ordered substitution]\label{def:ordered-substitution}
Let $\sigma$ be a substitution assigning to each free slot of the antecedent block a
position in a target context.  The \emph{ordered substitution} $P\sigma$ is obtained by
replacing each free occurrence in order of the context, without permutation.
\end{definition}

\begin{theorem}[Ordered substitution with ledger splice]\label{thm:substitution-splice}
Let $\Lambda_\sigma$ be the ledger change induced by the substitution $\sigma$.  For a
one-hole context $L\,Q\,R$ in which $Q$ sits between a left ledger block $L$ and a right
ledger block $R$,
\[
  \mathrm{ledger}\bigl(L\,\Lambda_\sigma\,R\bigr)
  =
  L\circ\Lambda_\sigma\circ R ,
\]
an exact splice: the left block is untouched, the middle block is transported by $\sigma$,
and the right block is untouched, with no reordering.
\end{theorem}

\begin{proof}
Induction on the context $L\,Q\,R$.  At the hole the ledger is exactly
$\Lambda_\sigma$ by Definition~\ref{def:ordered-substitution}.  For each rule above or
below the hole, the rule's own ledger contribution is a block that does not commute with
the hole's block, which is the content of linearity: the induction hypothesis is applied
to the strict subcontexts, and the blocks compose in the order fixed by the derivation.
\end{proof}

\begin{corollary}[Fresh-cut admissibility, conditionally]\label{cor:fresh-cut}
If the transported ledger $\Lambda_\sigma$ is \emph{fresh}, meaning no antecedent is
consumed twice, then the cut on the substituted formula is admissible: the composite
$L\,\Lambda_\sigma\,R$ is derivable whenever the substitutes are.
\end{corollary}

\begin{proof}
Apply Theorem~\ref{thm:substitution-splice} and check that each antecedent in the spliced
ledger occurs exactly once.  Freshness is exactly the condition that no contraction is
implicitly required, and without it the splice would duplicate an antecedent in violation
of Definition~\ref{def:tnd0}(2).
\end{proof}

\begin{remark}[The condition is not removable]
Corollary~\ref{cor:fresh-cut} is conditional and stays conditional.  Dropping freshness
would allow a ledger with a repeated antecedent, which no rule of the fragment can
introduce; the resulting "admissible" cut would have no derivation.
\end{remark}

\subsection{Raw ledger equality fails}

\begin{nogo}[Ledger equality is not the right invariant]\label{nogo:raw-ledger}
Raw equality of ledgers is not preserved by beta-like reduction: there are a proof tree and
its reduct whose ledgers are unequal as words while both record the same current authority.
Consequently the target of the proof transformation is not ledger equality but current
ledger equality plus a certificate for the retired part.
\end{nogo}

\begin{proof}
A principal-implication detour consumes the antecedent in one order and its reduct consumes
it in another; both consume it exactly once.  The ledgers therefore differ as words while
agreeing in the multiset of current antecedents.
\end{proof}

\begin{proposal}[Ledger transport, not ledger identity]\label{prop:ledger-transport}
The correct statement is a \emph{transport}: a reduction relates the pre-state ledger to the
post-state ledger by a declared transformation, and retired authority leaves the current
ledger while surviving in the certificate.  This is controlled forgetting in the sense of
Proposition~\ref{prop:three-grades}, not equality.
\end{proposal}

\subsection{Transport through contexts}

\begin{lemma}[Freshness reflects through discharge]\label{lem:freshness-reflects}
If a ledger is fresh above a discharge, then its image under the discharge is fresh, and
conversely a fresh image implies freshness above whenever the discharge consumes exactly
one antecedent.
\end{lemma}

\begin{proof}
The discharge removes exactly one antecedent occurrence and appends the discharged
assumption to the premise's right block.  Freshness is a condition on occurrences, so it is
preserved in both directions under this exact-one bookkeeping.
\end{proof}

\begin{theorem}[Root transport lifts through every context]\label{thm:root-transport}
Let
\[
  L\,[x^{b}]\,Q\,S\ \leadsto\ L\,S\,Q
\]
be the root beta-like transport of a principal detour.  Then for every one-hole proof
context $C[\cdot]$,
\[
  C\bigl(L\,[x^{b}]\,Q\,S\bigr)\ \leadsto\ C\bigl(L\,S\,Q\bigr),
\]
and the transported ledger differs from the original by the declared transformation only.
\end{theorem}

\begin{proof}
By induction on the context, using Lemma~\ref{lem:freshness-reflects} at each discharge and
Theorem~\ref{thm:substitution-splice} at each substitution.  Linearity ensures that no
context can duplicate the detour, so the transport is well defined at every depth.
\end{proof}

\begin{corollary}[Contextual closure]\label{cor:contextual-closure}
The beta-like reduction relation of $\mathrm{TND}_0$ is closed under one-hole contexts.
\end{corollary}

\begin{proof}
Immediate from Theorem~\ref{thm:root-transport}.
\end{proof}

\subsection{Strong normalization}

\begin{definition}[Proof-tree ledger and discharge count]\label{def:proof-ledger}
For a checked $\mathrm{TND}_0$ proof tree $D$ let $\Lambda_D$ be its \emph{complete ledger}:
the word of resource records of $D$ in tree order, each carrying an occurrence key, a
status, and, when discharged, the binder that discharged it.  The status is
\emph{open} for a use still owed to the ordered context and \emph{discharged} for a use
bound by $(\multimap I)$.  The open projection of $\Lambda_D$ is the ordered context of the
judgment at the root, so the complete ledger is finer than the ledger of
Definition~\ref{def:ledger-judgment}, which records only the current part.  Write
$\operatorname{dchg}(\Lambda_D)$ for the number of discharged records, and
$\operatorname{nodes}(D)$ for the number of nodes of the tree.
\end{definition}

\begin{theorem}[Exact ledger descent]\label{thm:ledger-descent}
For every certified contextual beta step $D\Longrightarrow_\beta D'$,
\begin{equation}
  |\Lambda_{D'}|=|\Lambda_D|-1,
  \qquad
  \operatorname{dchg}(\Lambda_{D'})=\operatorname{dchg}(\Lambda_D)-1 .
  \label{eq:ledger-descent}
\end{equation}
In the unfolded proof-tree representation one also has
$\operatorname{nodes}(D')=\operatorname{nodes}(D)-3$.
\end{theorem}

\begin{proof}
At the root, the transport of Theorem~\ref{thm:root-transport} replaces
\[
  \Lambda_R
  =L\circ[x^{\mathrm{discharged},b}]\circ Q\circ\Lambda_\sigma
  \quad\leadsto\quad
  \Lambda_C
  =L\circ\Lambda_\sigma\circ Q ,
\]
retiring exactly the bound target record $x$ and moving the argument block $S=\Lambda_\sigma$
across the discharged suffix $Q$; no other record is deleted, duplicated, relabelled, or
renamed, so both counts fall by exactly one.  An introduction ancestor discharges the same
surviving open record on both sides and extends both ledgers by one record at the same
position; either elimination ancestor concatenates the same sibling ledger on both sides.
Induction along the proof-position path therefore gives \eqref{eq:ledger-descent} at the
outer root.

For the node count, linearity makes substitution graft $\sigma$ exactly once: the reduct
removes the elimination node, the introduction node, and the target assumption leaf, and
nothing else.
\end{proof}

\begin{remark}[Why the measure is a ledger count and not a detour count]
Counting principal-implication detours is the tempting alternative, since a root step
contracts one.  It is not a descent measure: contracting a detour can expose another at the
same position.  In the $\lambda$-term reading of the fragment,
\[
  \bigl((\lambda f.f)(\lambda x.x)\bigr)z
  \;\longrightarrow\;
  (\lambda x.x)z
\]
has exactly one principal detour on each side; every variable is used once and no exchange
is needed.  The ledger counts of Theorem~\ref{thm:ledger-descent} are immune to this,
because they measure what a step retires rather than what it exposes.
\end{remark}

\begin{theorem}[Strong normalization]\label{thm:strong-normalization}
Every reduction sequence in $\mathrm{TND}_0$ terminates.  Indeed, a sequence starting at $D$
has length at most $\operatorname{dchg}(\Lambda_D)$.
\end{theorem}

\begin{proof}
By Theorem~\ref{thm:root-transport} a step may be applied at any depth, and by
Theorem~\ref{thm:ledger-descent} each step lowers $\operatorname{dchg}$ by exactly one.  An
infinite sequence would therefore give an infinite strictly descending sequence of natural
numbers, which is impossible.
\end{proof}

\begin{remark}[Where linearity was used]
Theorem~\ref{thm:strong-normalization} depends on linearity three times: in
Theorem~\ref{thm:substitution-splice}, which forbids reordering the ledger; in
Lemma~\ref{lem:freshness-reflects}, which keeps freshness intact through discharge; and in
the node count, where linearity makes the graft single-use.  The same argument fails for a
fragment with contraction, where one retires one record while copying a subtree.
\end{remark}

\subsection{Summary}

\begin{enumerate}
\item Ordered substitution produces an exact ledger splice
  (Theorem~\ref{thm:substitution-splice}).
\item Cut admissibility is conditional on freshness
  (Corollary~\ref{cor:fresh-cut}), and the condition is not removable.
\item Ledger equality is the wrong invariant; transport plus certificate is right
  (No-go~\ref{nogo:raw-ledger}, Proposal~\ref{prop:ledger-transport}).
\item Root transport lifts through all contexts, and reduction is strongly normalizing
  (Theorem~\ref{thm:root-transport}, Corollary~\ref{cor:contextual-closure},
  Theorem~\ref{thm:strong-normalization}).
\end{enumerate}

\section{Confluence and decidable derivability}\label{sec:confluence-decidability}
% Section 15: Confluence and decidable derivability
\label{sec:confluence-body}

Normalization in \S\ref{sec:normalization} guarantees that a normal form exists.  It does
not guarantee that the normal form is unique, and uniqueness is what a decision procedure
needs.  This section supplies the critical-pair analysis that yields confluence, then the
normal/neutral characterization that yields the subformula property, and finally the
focused decision procedure that follows.  It closes by describing precisely what finite
decidability does \emph{not} give.

\subsection{Critical pairs}

Because $\mathrm{TND}_0$ is linear and ordered, the only way two redexes can interfere is
that their ledgers overlap.  This makes the critical-pair analysis finite and explicit.

\begin{definition}[Overlapping redexes]\label{def:overlapping-redexes}
Two redexes in a $\mathrm{TND}_0$ proof tree \emph{overlap} when their ledgers share an
antecedent occurrence.  Since ledgers are consumed linearly
(Definition~\ref{def:tnd0}), a shared occurrence means one redex lies on the path from the
other to its discharged assumption.
\end{definition}

\begin{proposition}[Four-way path classification]\label{prop:four-paths}
For an overlapping pair of redexes, exactly one of four configurations occurs:
\begin{enumerate}
\item one redex strictly above the other, on the discharged spine;
\item one redex strictly above the other, off the discharged spine;
\item the two redexes are adjacent along one ledger block;
\item the two redexes are disjoint but order-adjacent in the ledger.
\end{enumerate}
\end{proposition}

\begin{proof}
By Definition~\ref{def:overlapping-redexes} the two redexes lie on one spine, and the
relative position of the shared occurrence in their ledger blocks gives the four cases.
Linearity excludes any configuration in which a ledger block is shared without one redex
containing the other.
\end{proof}

\begin{lemma}[Substitution and graft equations]\label{lem:substitution-graft}
In configurations (1) and (3) the two reduction orders are related by the ledger splice of
Theorem~\ref{thm:substitution-splice}; in configurations (2) and (4) they commute exactly.
\end{lemma}

\begin{proof}
Configurations (2) and (4) act on disjoint ledger blocks, so the root transports of
Theorem~\ref{thm:root-transport} commute by inspection.  Configurations (1) and (3) nest
one substitution inside the other, and the splice law identifies the two composites; the
only thing to check is that the ledger blocks compose in the same order in both, which is
the content of freshness.
\end{proof}

\subsection{The one-step strong diamond}

\begin{theorem}[One-step strong diamond]\label{thm:strong-diamond}
If $P\leadsto P_1$ and $P\leadsto P_2$ with $P_1\neq P_2$, then there is $P_3$ with
$P_1\leadsto P_3$ and $P_2\leadsto P_3$, and the two-step reductions can be chosen to have
equal length.
\end{theorem}

\begin{proof}
By Lemma~\ref{lem:substitution-graft}, each configuration of
Proposition~\ref{prop:four-paths} either commutes or is joined after one further step.  The
equal-length clause follows because both reductions remove the same number of detours in
configurations (2) and (4), and differ by a single commuting step in (1) and (3).
\end{proof}

\begin{corollary}[Local and global confluence]\label{cor:confluence}
The reduction relation of $\mathrm{TND}_0$ is locally confluent and, being strongly
normalizing (Theorem~\ref{thm:strong-normalization}), globally confluent by Newman's
lemma.  Consequently every proof has a unique normal form.
\end{corollary}

\begin{remark}
Corollary~\ref{cor:confluence} is stated for a \emph{fixed} proof tree.  It does not say
that two distinct proofs of the same sequent share a normal form; the fragment is
incomplete (Theorem~\ref{thm:incompleteness}), so distinct derivations of one sequent need
not be related at all.
\end{remark}

\subsection{Confluence is not coherence}

\begin{nogo}[Confluence does not give a 2-cell structure]\label{nogo:confluence-coherence}
The fact that reductions commute is not the fact that they are coherent.  There is a
four-level ladder
\[
  \text{proof endpoint}\ \subsetneq\ \text{composite boundary transport}\ \subsetneq\
  \text{history}\ \subsetneq\ \text{history 2-cell},
\]
and confluence at the first level determines nothing at the fourth.  Reading
Corollary~\ref{cor:confluence} as a coherence theorem is therefore invalid.
\end{nogo}

\begin{proof}
Levels 1 and 2 are quotients of level 3, and level 4 records relations among histories
that are collapsed by the quotient.  Confluence is a statement about the image at level 1;
a nontrivial 2-cell lives in the kernel of the forgetful map from level 3 to level 1, and
that kernel is not determined by level 1.
\end{proof}

\begin{remark}
No-go~\ref{nogo:confluence-coherence} is the proof-theoretic counterpart of
Proposition~\ref{prop:lossy-quotient}.  In both cases the error is the same: taking a
quotient for the object.
\end{remark}

\subsection{Normal and neutral forms}

\begin{definition}[Normal and neutral]\label{def:normal-neutral}
A derivation is \emph{normal} if it contains no principal-implication detour.  It is
\emph{neutral} if it ends in an elimination whose major premise is a hypothesis.  The
standard characterization is: normal derivations end either in an introduction or are
neutral, and neutral derivations are closed under the eliminations of the fragment.
\end{definition}

\begin{theorem}[Subformula property]\label{thm:subformula}
Every formula occurring in a normal $\mathrm{TND}_0$ derivation is a subformula of the end
sequent or of a discharged hypothesis.
\end{theorem}

\begin{proof}
Induction on the normal derivation using Definition~\ref{def:normal-neutral}.  An
introduction rule's premises contain only subformulas of its conclusion; an elimination
whose major premise is neutral has a minor premise that is a subformula of the major
premise's type, and the absence of detours means no formula is introduced only to be
eliminated.
\end{proof}

\begin{proposition}[Leftmost-head canonicity]\label{prop:leftmost-head}
In a normal derivation of an implicational sequent, the construction of the proof is
determined by the ordered context together with the leftmost hypothesis of the end
sequent's antecedent block.  Hence the leftmost-head strategy is complete for the
fragment.
\end{proposition}

\begin{proof}
By Theorem~\ref{thm:subformula} every formula of a normal derivation is a subformula of the
goal, and by linearity no antecedent is reused, so at each stage the possible rules are
determined by the leftmost unconsumed antecedent and the goal.  The choices are therefore
forced up to the order in which independent threads are processed.
\end{proof}

\subsection{Decidable derivability}

\begin{theorem}[Focused decision procedure]\label{thm:decidability}
Derivability in $\mathrm{TND}_0$ is decidable.  There is a terminating, sound, and complete
procedure: search normal derivations in leftmost-head order, using
Theorem~\ref{thm:subformula} to bound the formula set, Definition~\ref{def:normal-neutral}
to alternate focused and unfocused phases, and Proposition~\ref{prop:leftmost-head} to fix
the branching.
\end{theorem}

\begin{proof}
Termination: by Theorem~\ref{thm:subformula} only finitely many formulas may occur, and
linearity bounds the number of introductions by the number of antecedents.  Soundness: each
step of the procedure builds a derivation, so a found proof is a proof.  Completeness: by
Corollary~\ref{cor:confluence} a normal form exists if any derivation does, and by
Proposition~\ref{prop:leftmost-head} the procedure explores a branch containing it.
\end{proof}

\begin{proposition}[Positioning of the skeleton]\label{prop:skeleton-position}
The $\mathrm{TND}_0$ skeleton is the associative, empty-antecedent, right-residual fragment
of ordered linear logic.  Under this reading, the erasure of audited nominal data to a
skeleton is a functor, and lifting back is by fresh, noncanonical decoration; long forms
are not identified with their eta-expansions, and multiple focused proofs of one sequent
are admitted.
\end{proposition}

\begin{proof}[Status]
The fragment identification is a syntactic matching of rule shapes.
Proposition~\ref{prop:leftmost-head} and Theorem~\ref{thm:decidability} are the substantive
content; the erasure/lift discipline is a definitional choice, recorded so that the
decision procedure's non-canonicity of decoration is visible rather than implicit.
\end{proof}

\subsection{What finite decidability does not give}

\begin{warning}[No compactification from finite search]\label{warn:no-compactification}
Theorem~\ref{thm:decidability} is a statement about a finite fragment.  It does not
produce, and must not be read as producing, a halting measure, a limit object, or any
$\Omega$-type boundary.  Neither finite confluence
(Corollary~\ref{cor:confluence}) nor finite focused search
(Theorem~\ref{thm:decidability}) yields such an object; they are finite statements and the
boundary is not their limit.
\end{warning}

\begin{remark}
Warning~\ref{warn:no-compactification} is the sharpest instance of the discipline that
organises this paper.  A decision procedure terminates; that is not the same as exhibiting
a limit of a non-terminating process, and conflating the two would convert a finite theorem
into an infinite claim.
\end{remark}

\subsection{Summary}

\begin{enumerate}
\item The critical-pair analysis is finite, with four configurations
  (Proposition~\ref{prop:four-paths}, Lemma~\ref{lem:substitution-graft}).
\item Reduction satisfies a one-step strong diamond and is confluent, so normal forms are
  unique per proof (Theorem~\ref{thm:strong-diamond}, Corollary~\ref{cor:confluence}).
\item Confluence is not coherence (No-go~\ref{nogo:confluence-coherence}).
\item The subformula property plus leftmost-head canonicity give a terminating, sound, and
  complete decision procedure (Theorems~\ref{thm:subformula}, \ref{thm:decidability},
  Proposition~\ref{prop:leftmost-head}).
\item Decidability yields no $\Omega$-type boundary
  (Warning~\ref{warn:no-compactification}).
\end{enumerate}

\section{The six-port whole-cut kernel}\label{sec:six-port}
% Section 16: The six-port whole-cut kernel
\label{sec:six-port-body}

Part~II has so far supplied a semantics, a calculus, and a metatheory.  This section asks
how few ports a carrier needs in order to express that calculus with three readings and
binary cuts, and answers with a counting theorem.  The answer is six.  The section is
careful about what the number means: it is a conditional minimum under a stated contract,
not an absolute truth, and the universality claim that usually accompanies such a number is
replaced by relative initiality.

\subsection{The counting question}

\begin{definition}[Whole-cut carrier]\label{def:whole-cut-carrier}
A \emph{whole-cut carrier} with $m$ domains and $n$ ports consists of
\begin{enumerate}
\item a finite set $D$ of domains;
\item for each domain, a binary cut $\langle L\mid R\rangle$ that is an atomic generator
  rather than a derived construct;
\item a set $P$ of ports, each attached to one domain and one side of its cut;
\item an occurrence-and-multiplicity ledger $M$ recording, for every term, the
  production and consumption of each occurrence;
\item a routing datum $\gamma$, a history ledger $H$, and a residual record $R$.
\end{enumerate}
\end{definition}

\begin{definition}[The contract]\label{def:six-port-contract}
The counting theorem is stated under six clauses.
\begin{enumerate}
\item \textbf{Atomic cuts.}  Each cut is a generator; it is not decomposed into smaller
  structure.
\item \textbf{Binary sides.}  Each cut has exactly two complementary sides.
\item \textbf{Two-sided access.}  Every domain must be readable on both sides of its own
  cut, since a view that could read one side only could not express a crossing.
\item \textbf{Occurrence fidelity.}  No view, braid, collision, projection, or quotient may
  create, drop, rename, or merge an occurrence.
\item \textbf{Production witnesses.}  Copies require production witnesses; collisions
  conserve multiplicity.
\item \textbf{View discipline.}  Every view extends the ledger and commutes with the
  boundary operator: $V_\nu(\partial c)=\partial V_\nu(c)$.
\end{enumerate}
\end{definition}

\subsection{The minimality theorem}

\begin{theorem}[Counting bound]\label{thm:counting-bound}
Under the contract of Definition~\ref{def:six-port-contract}, a triadic whole-cut carrier
with binary cuts satisfies
\[
  n\ \ge\ 2m ,
\]
and hence, for $m=3$, $n_{\min}=6$.
\end{theorem}

\begin{proof}
By clause (3) every domain must be accessible on both sides of its own cut, so each domain
contributes at least two distinct ports; by clause (1) the cut is atomic, so those two ports
cannot be served by sub-structure; by clause (2) the two sides are complementary, so a
single port cannot serve both.  Hence $n\ge 2m$, with equality attained by the carrier
$\bW_6$ of \S\ref{sec:six-port-carrier}.
\end{proof}

\begin{remark}[What the proof does and does not use]
The bound uses clauses (1)--(3) only; clauses (4)--(6) are needed for the carrier to be
usable, not for the count.  This split matters, because it isolates exactly which part of
the contract is under pressure if the number six is ever challenged.
\end{remark}

\begin{warning}[Six is conditional]\label{warn:six-conditional}
Theorem~\ref{thm:counting-bound} says that six is minimal \emph{for this contract}.  Drop
clause (3) and a carrier with three ports becomes available; strengthen clause (1) to
require decomposed cuts and the minimum changes again.  The number six is therefore a
statement about the contract, and reporting it as an absolute minimality of any
three-domain syntax would be an error.
\end{warning}

\subsection{The carrier}\label{sec:six-port-carrier}

\begin{definition}[The six-port carrier]\label{def:w6}
The carrier is
\[
  \bW_6=(\mathsf D,\ C,\ P,\ O,\ M,\ \gamma,\ H,\ R)
\]
consisting of domains, cuts, ports, occurrences, the multiplicity ledger, routing, history,
and the residual record.  Its distinguishing structural features are:
\begin{enumerate}
\item a \emph{split pair}, giving the two sides of each atomic cut;
\item a \emph{projector}, identifying which ports a given view reads;
\item a \emph{factored traversal}, so that a traversal decomposes along the cut rather than
  cutting across it;
\item \emph{collision strata}, recording where two occurrences coincide.
\end{enumerate}
\end{definition}

\begin{proposal}[Line and circle are views]\label{prop:line-circle-views}
In $\bW_6$ the \emph{line} and the \emph{circle} are not neighbouring value constructors but
checked projections of one ledger, each with a declared kernel.  One atomic cut
co-generates a continuation germ and a boundary germ; the line reads the first and the
circle the second.
\end{proposal}

\begin{proof}[Status]
A definitional restructuring, justified by clause (6): both readings are views
$V_\nu$ of the same boundary operator, so neither is primitive.  The claim asserted is the
commutation $V_\nu(\partial c)=\partial V_\nu(c)$, which holds by construction of the two
views and is checked on the finite fixture.
\end{proof}

\subsection{The grammar, the ledgers, and the degree}

\begin{definition}[Whole-cut record]\label{def:whole-cut-record}
A \emph{whole-cut record} consists of port, cut, and cell records together with six ledgers:
\emph{production}, \emph{through}, \emph{cycle}, \emph{hole}, \emph{alternative}, and
\emph{residual}.  The record carries a \emph{degree} $\deg(W)$, and for the minimal carrier
$\deg(\bW_6)=6$.
\end{definition}

\begin{proposition}[Collision conservation]\label{prop:collision-conservation}
In any well-typed whole-cut record the multiplicity ledger is conserved under every rule of
the grammar: no production is created or destroyed, and every collision accounts for
exactly one coincidence of two occurrences.
\end{proposition}

\begin{proof}
Each grammar rule consumes and produces ledger entries in matching multiplicity; the
production ledger records copies explicitly, so a rule that duplicates a term must record a
production, and clause (5) forbids doing so silently.  Induction over derivations in the
grammar.
\end{proof}

\begin{remark}[Distributivity as a recorded witness]
The grammar's distributivity step is accompanied by an explicit witness record with the
ordered hole bindings.  The content is not the equation --- which holds in the intended
model --- but the record of which occurrences of the duplicated term were produced, in
which order.  Without that record the equation would silently license a contraction, which
No-go~\ref{nogo:raw-ledger} and Proposition~\ref{prop:collision-conservation} forbid.
\end{remark}

\subsection{The cell--carrier--view spine}

\begin{proposal}[Spine]\label{prop:cell-carrier-view}
There is a typed spine
\[
  \mathsf{TypedCellSignature}
  \xrightarrow{\ \mathsf{assemble}_\rho\ }
  \mathsf{StratifiedCarrier}
  \xrightarrow{\ V_\nu\ }
  \mathsf{PresentedView}_\nu ,
\]
in which the first arrow assembles cells into a stratified carrier along a declared
stratification $\rho$, and the second is the view of clause (6).  The boundary operators of
the two basic cells are
\[
  \partial\mathsf{InterchangeCell}=Q_4,\qquad
  \partial\mathsf{BraidCell}=M_6 ,
\]
with Coxeter shadows $C_2\times C_2$ and $S_3$ respectively.
\end{proposal}

\begin{proof}[Status]
The boundary identifications are computations in the cell complex: an interchange cell is a
square whose boundary reads as a product of two transpositions, a braid cell a hexagon whose
boundary reads as the three-strand braid relation.  The Coxeter shadows are the standard
reflections of the square and hexagon.
\end{proof}

\begin{proposition}[Envelope count]\label{prop:envelope}
The truncated-octahedral envelope of the stratum satisfies
\[
  TO_{24}=6\,Q_4+8\,M_6 ,
\]
that is, six squares and eight hexagons on twenty-four vertices, and the relation boundary
count is $2m$ in the sense of Theorem~\ref{thm:counting-bound}.
\end{proposition}

\begin{proof}
$6\cdot4+8\cdot6=72$ face--vertex incidences on $24$ vertices with degree $3$:
$24\cdot3=72$.  The consistency of the two counts witnesses that the envelope closes.
\end{proof}

\subsection{The $M_6$ boundary bridge}

\begin{theorem}[Alternating-pairing bridge]\label{thm:m6-bridge}
Let $\kappa$ and $\tau$ be two fixed-point-free involutions on the six-port ledger whose
union is connected.  Then the union of their graphs is a two-regular graph on six vertices,
hence a six-cycle $C_6$, whose edges alternate between the two kinds.  Consequently the
checked incidence data present the open boundary
\[
  \mathsf{bridge}(\bW_6)=M_6^{\partial}\bigl[\kappa\tau\kappa
  \Rightarrow\tau\kappa\tau\bigr],
\]
with no filler supplied.
\end{theorem}

\begin{proof}
A fixed-point-free involution on six elements is a perfect matching, so it contributes three
edges.  Two matchings whose union is connected give a two-regular graph, which is a disjoint
union of cycles; connectivity forces a single six-cycle.  Alternation of edge kinds follows
from the two matchings being distinct, and the displayed presentation of $M_6^{\partial}$
records the two alternating readings of the cycle.  The absence of a filler is
Calibration~\ref{cal:m6-open}.
\end{proof}

\begin{calibration}[The filler is open]\label{cal:m6-open}
The bridge of Theorem~\ref{thm:m6-bridge} presents the boundary $M_6^{\partial}$ as an open
object: the two alternating readings are exhibited, and no 2-cell filling them is
constructed.  Every construction in this material that touches $M_6$ lists the filler as an
open residual.
\end{calibration}

\subsection{Relative initiality instead of universality}

\begin{proposal}[Relative initiality]\label{prop:relative-initiality}
$\bW_6$ is initial in the category of presented models satisfying the contract of
Definition~\ref{def:six-port-contract}: for every such model $N$ there is a morphism
$\bW_6\to N$, unique up to the declared decoration.
\end{proposal}

\begin{proof}[Status]
The statement is a \textsf{Proposal}.  What is available is the counting argument of
Theorem~\ref{thm:counting-bound} together with the view discipline of clause (6), which
supply the candidate morphisms; no uniqueness proof is given here.
\end{proof}

\begin{warning}[Not a universality theorem]\label{warn:not-universal}
Proposition~\ref{prop:relative-initiality} replaces, and does not establish, a claim of
universality.  In particular it says nothing about computation-universality, and it does not
assert that every three-domain syntax factors through $\bW_6$.  The distinction is the
departure point of \S\ref{sec:boundaries}.
\end{warning}

\subsection{Summary}

\begin{enumerate}
\item Six is the minimum under a stated six-clause contract
  (Definitions~\ref{def:whole-cut-carrier}, \ref{def:six-port-contract},
  Theorem~\ref{thm:counting-bound}, Warning~\ref{warn:six-conditional}).
\item In the minimal carrier, line and circle are views of one ledger
  (Definition~\ref{def:w6}, Proposal~\ref{prop:line-circle-views}).
\item The ledgers conserve multiplicity, so distributivity must record its productions
  (Proposition~\ref{prop:collision-conservation}).
\item The spine and envelope counts are explicit
  (Proposal~\ref{prop:cell-carrier-view}, Proposition~\ref{prop:envelope}).
\item The $M_6$ boundary is presented with an open filler, and the kernel's status is
  relative initiality, not universality
  (Theorem~\ref{thm:m6-bridge}, Calibration~\ref{cal:m6-open},
  Proposal~\ref{prop:relative-initiality}, Warning~\ref{warn:not-universal}).
\end{enumerate}

\section{Claim boundaries and the remaining programme}\label{sec:boundaries}
% Section 17: Claim boundaries and the remaining programme
\label{sec:boundaries-body}

This paper has reported more boundary results than theorems, and the imbalance is
deliberate.  This section collects the boundaries, states which of them are open problems
rather than exclusions, and identifies the single gap on which the triadic programme
depends.

\subsection{The decisive gap: order from energy}

Every layer of Part~I that computes a quantity --- period grades, spectral packets, energy
vectors, trace counts --- is a map into something ordered or something measurable.  None of
them produces the structure that a cut actually needs.

\begin{problem}[Order from a finite frontier]\label{prob:order-from-energy}
\label{prob:decisive-gap}
Construct a chart-covariant certified comparison on finite period-energy frontiers whose
objectification yields the correct order relation on the resulting numbers.
\end{problem}

\begin{remark}[Why this is the decisive gap]
Energy and period data determine \emph{what is visible}.  Nothing in this material
determines \emph{what is left and what is right}.  The identification of energy grades with
the two sides of a cut was explicitly rejected as unlicensed
(No-go~\ref{nogo:zero-energy-trap} shows that energy alone is vacuous as a selector, and
Proposition~\ref{prop:energy-invariance} that it detects only invariance).  Until
Problem~\ref{prob:decisive-gap} is solved, the triadic programme has carriers and residuals
but no order, and by Definition~\ref{def:satisfaction-form} a logic without order has
evidence but no entailment.
\end{remark}

\begin{remark}[The decisive gap is not a shortage of technique]
\label{rem:gap-is-not-scarcity}
Three no-go results of this paper already remove the obvious candidate classes of
ingredient for the order that Problem~\ref{prob:decisive-gap} asks for.
No-go~\ref{nogo:zero-energy-trap} shows that energy alone is vacuous as a selector;
Proposition~\ref{prop:energy-invariance} shows that what an energy detects is
invariance; and No-go~\ref{nogo:mask-direction} shows that mask data, being
invariant under the orientation-reversing involution of the carrier, carries no
direction at all.  Each of these is a scalar or finitely-valued functional of the
carrier's state, and each is invariant under the symmetry that would have to be
broken.  The gap is therefore not the absence of a clever enough functional; what
is missing is a datum of a different kind.  We do not assert what that kind is, and
this remark proves nothing beyond the three cited no-gos; we record it so that the
programme's first task is not pursued by searching for another scalar.
\end{remark}

\subsection{Open problems}

\begin{problem}[Canonical identification of the three readings]\label{prob:canonical-readings}
Decide whether the three readings $\rho_t,\rho_X,\rho_\cK$ admit a canonical
identification on the Legendre carrier, or whether only the residual-laden comparison of
Theorem~\ref{thm:crossing-squares} exists.  The mod-two description
(Proposition~\ref{prop:three-cusps}) shows that at least the sign must be tracked; whether
anything further is canonical is open.
\end{problem}

\begin{problem}[Non-decomposable triadic defect]\label{prob:non-decomposable}
Prove or refute Conjecture~\ref{conj:triadic-crossing}: does there exist a carrier on which
the cyclic defect of the three crossings is not reconstructible from any pair?  The finite
Legendre prism fails (Theorem~\ref{thm:defect-decomposes}), and the finite satisfaction form
fails (Calibration~\ref{cal:triadicity-fails}); both failures are decompositions of a
specific kind, and neither excludes a carrier of a different kind.
\end{problem}

\begin{problem}[The $M_6$ filler]\label{prob:m6-filler}
Construct the 2-cell that fills $M_6^{\partial}$ as presented by
Theorem~\ref{thm:m6-bridge}, or prove that the alternating presentation admits no filler.
The boundary is exhibited; the interior is not.
\end{problem}

\begin{problem}[A typed feedback operator]\label{prob:feedback}
Supply a syntactic operator whose semantics is genuinely periodic, thereby repairing
No-go~\ref{nogo:return-not-period}.  Boundary-return words exist; periodic dynamics does
not.  The obstruction is that completed-past rank is monotone, so any repair must change the
carrier rather than add a formula.
\end{problem}

\begin{problem}[Truncation is residual-valued]\label{prob:truncation}
Finite energy sublevels are not closed under addition, multiplication, or grafting, so
truncation satisfies $xy=P_\Lambda(xy)+R_\Lambda(x,y)$ with nonvanishing residual.  Give a
calculus in which this residual composes.
\end{problem}

\begin{problem}[Vocabulary gradients]\label{prob:vocabulary}
For a residual functional $E$ on monotone Boolean vocabularies, the diminishing-returns
inequality holds, and the heuristic "largest local gain is best" is refuted by explicit
counterexample.  Determine the correct admissible search discipline, and connect it to the
coaction of Proposition~\ref{prop:binomial-coaction}.
\end{problem}

\begin{problem}[Learning and proof as reverse duals]\label{prob:reverse-dual}
The duality proposed in this material is contravariant: proof normalises toward a sink and
learning is the reverse generative orientation, with the forward equation false.  Supply the
formal dual pair and prove that the two orientations are exchanged, or show that only one
direction exists.
\end{problem}

\subsection{What is excluded, not merely open}

The following are excluded by scope, and no amount of further work inside this paper would
include them.

\begin{enumerate}
\item \textbf{Universality claims.}  The kernel's status is relative initiality
  (Proposition~\ref{prop:relative-initiality},
  Warning~\ref{warn:not-universal}).  Computation-universality of the triadic form is not
  claimed at any point.
\item \textbf{Physical interpretation.}  The six-state carrier is a finite combinatorial
  object (Theorem~\ref{thm:endogenous-clifford}); spacetime, metric signature, and physical
  realization are outside the paper.  Warning~\ref{warn:temporality} removes the most
  tempting identification explicitly.
\item \textbf{Complex scalars as a foundational necessity.}  What is proved is a selection
  within a tested family under a bridge hypothesis
  (Proposition~\ref{prop:negative-inverse}, Calibration~\ref{cal:bridge-hypothesis}).
\item \textbf{Computational lower bounds.}  No growth rate, no entropy inequality, and no
  fibre size in this paper yields a runtime or space bound.  Paper~IV owns that territory
  under stricter hypotheses \cite{Yuan2026AEGCondensation}.
\item \textbf{$\Omega$-type boundaries from finite results.}  Finite confluence and finite
  focused search produce no limit object
  (Warning~\ref{warn:no-compactification}).
\item \textbf{Knot-theoretic invariance.}  The figure-eight calibration
  (Theorem~\ref{thm:figure-eight}) identifies one feature of one word; it does not construct
  a knot invariant, and the polynomial does not reconstruct the knot.
\end{enumerate}

\subsection{Engineering excluded from the paper}

The governance scope of this paper
(\texttt{governance/00c-paper-V-triadic-aperture-scope.md}) excludes the engineering
scaffolding of the source corpus: checked-program harnesses, test suites, authority splits,
work-package phases, document formats, ledgers of artifacts, and CI gates.  Fifteen of the
133 surveyed source files contribute no main-text mathematics and are cited at most as
provenance.  Where a mathematical claim in this paper depends on an executed check, the
dependence is recorded in the claim ledger of Appendix~\ref{app:ledger} rather than
concealed.

Appendix~\ref{app:ledger} also lists the source clusters that are mathematical but belong to
neither part of this paper --- the causal-cut and Alexandrov-duality notes, the backward
transport and cut-composition notes, the $PSP^{\ast}$ programme, the $\Omega$-boundary
cluster, and the later formation-gate notes.  Those omissions are deliberate and are recorded
so that they are visible rather than accidental; they are the natural content of a subsequent
paper.

\subsection{The programme}

Four tasks follow, in order of logical dependence.

\begin{enumerate}
\item \textbf{Order.}  Solve Problem~\ref{prob:decisive-gap}.  Nothing downstream is stable
  until a comparison is available, because a logic without order is only evidence.
\item \textbf{Filler.}  Solve Problem~\ref{prob:m6-filler}, which completes the kernel's
  presentation and would decide whether the triadic structure is carried by the cell or only
  by its boundary.
\item \textbf{Triadicity.}  Settle Problem~\ref{prob:non-decomposable}, which decides
  whether the paper's central negative calibration is a limitation of the finite models or a
  feature of the phenomenon.
\item \textbf{Motion.}  Solve Problem~\ref{prob:feedback}, which is what would turn the
  syntax of Part~II into a semantics with dynamics.
\end{enumerate}

\begin{remark}[On the shape of the programme]
Three of the four tasks are constructive and one is a decision.  That ratio is the honest
summary of the state of this material: the carriers and the residuals exist, and what is
missing is the additional structure that would make them a theory of order and motion.
\end{remark}

\section{Conclusion}\label{sec:conclusion}
% Section 18: Conclusion
\label{sec:conclusion-body}

This paper set out to ask what structure is needed for one carrier to be read three ways,
what such readings necessarily lose, and where the loss can be written down and checked.  The
answers came in two parts.

Part~I produced a theory of characteristics in which the primitive is an observer contract
and the product is a quotient together with a residual.  Within that framework three results
are stable.  The quadratic trichotomy shows that no coarse requirement selects the complex
field, and that what selects it is the retention of a central sign --- chirality, isolated as
the datum distinguishing the two lifts of one projective exchange.  The six-state carrier is
not postulated but generated: cut-membership cochains and oriented interchange produce the
Clifford relation on a declared local cut interval, and interchange alone provably cannot.
And crossing is a cospan with a resolution torsor, from which monodromy is recovered as a
difference of resolutions and the Legendre monodromy follows from a quadratic base change.

Part~I also produced a systematic catalogue of shadows that do not determine construction ---
spectra, gradients, snapshots, objectified values, masses, counts, windings, energies --- and
showed that the repair is never a finer scalar but always a residual.  Where a residual is
supplied the theory becomes faithful: marked jets and declared port charts are the two
worked instances.  One obstruction in the catalogue comes with an exact minimum
(Theorem~\ref{thm:order-obstruction}), and that theorem is the standard the others have not
yet met.

Part~II moved from what a reading cannot see to the places where it records the absence.  The
typed aperture is primitive, its open and close are not inverse, and four statuses are
required because exhaustion and unresolvedness are different things.  On that primitive a
finite semantics, a threaded adequacy theorem under four named hypotheses, an ordered linear
calculus that is sound but provably incomplete, and a metatheory were built: an ordered
substitution theorem with an exact ledger splice, conditional fresh-cut admissibility, strong
normalization via two decreasing measures, confluence with unique normal forms, the
subformula property, and a terminating and complete focused decision procedure.  The last
part of the paper answered the counting question for the kernel: six ports, under a stated
six-clause contract, with relative initiality in place of universality and an explicitly open
filler on the $M_6$ boundary.

Two negative results are, in the author's view, the most valuable in the paper, and both are
reported as failures of the finite models rather than as successes.  The triadic defect on
the Legendre prism decomposes (Theorem~\ref{thm:defect-decomposes}), and the finite
satisfaction form fails the triadicity test (Calibration~\ref{cal:triadicity-fails}).  A
reader who takes the word ``triadic'' in the title as a claim that a non-decomposable
three-fold coherence has been exhibited should take it instead as naming the programme that
Problem~\ref{prob:non-decomposable} asks whether is realisable.

What the paper does not have is order and motion.  Energy selects what is visible and not what
is left or right (Problem~\ref{prob:decisive-gap}); boundary-return words exist and periodic
dynamics does not (No-go~\ref{nogo:return-not-period}).  These two absences are not
oversights: they are what is left when the carriers have been built honestly and the
additional structure has not been invented.  Given the choice between a paper that claims a
completed triadic theory and one that reports exactly where the completeness fails, this is
the second.


\appendix
\part*{Appendices}
\addcontentsline{toc}{part}{Appendices}

\section{Scalar, chiral, and Clifford calculations}\label{app:scalar-clifford}
% Appendix A: Scalar, chiral, and Clifford calculations
\label{app:scalar-body}

This appendix records the computations used in \S\ref{sec:scalar-trichotomy} and
\S\ref{sec:six-state}.  Nothing here is new; the purpose is to have the exact
formulas in one place, so that a reader can check the paper's claims without consulting the
source notes.

\subsection{The three quadratic algebras}

For $\kappa\in\{-1,0,1\}$ and $x=a+bj$, $y=c+dj$:
\[
  xy=(ac+\kappa bd)+(ad+bc)j,\qquad
  \bar x=a-bj,\qquad
  N_\kappa(x)=x\bar x=a^{2}-\kappa b^{2}.
\]

\subsubsection{Multiplicativity of the norm}
Expanding both sides,
\[
  N_\kappa(xy)
  =(ac+\kappa bd)^{2}-\kappa(ad+bc)^{2}
  =a^{2}c^{2}+2\kappa abcd+\kappa^{2}b^{2}d^{2}
   -\kappa a^{2}d^{2}-2\kappa abcd-\kappa b^{2}c^{2},
\]
so the cross terms cancel and
\[
  N_\kappa(xy)=a^{2}c^{2}-\kappa a^{2}d^{2}-\kappa b^{2}c^{2}+\kappa^{2}b^{2}d^{2}
  =(a^{2}-\kappa b^{2})(c^{2}-\kappa d^{2})=N_\kappa(x)N_\kappa(y).
\]

\subsubsection{Inverses}
When $N_\kappa(x)=a^{2}-\kappa b^{2}\neq0$,
\[
  x^{-1}=\frac{\bar x}{N_\kappa(x)}=\frac{a-bj}{a^{2}-\kappa b^{2}} .
\]
The excluded loci are therefore:
\[
  \kappa=-1:\ \varnothing\ \text{off the origin};\qquad
  \kappa=0:\ a=0;\qquad
  \kappa=1:\ a=\pm b .
\]

\subsection{The negative-inverse involution}

With $S_\kappa(z)=-1/z$ and $z=a+bj$,
\[
  S_\kappa(a,b)=\Bigl(\frac{-a}{a^{2}-\kappa b^{2}},\ \frac{b}{a^{2}-\kappa b^{2}}\Bigr).
\]
Then
\[
  z\,S_\kappa(z)=(a+bj)\cdot\frac{-a+bj}{a^{2}-\kappa b^{2}}
  =\frac{-a^{2}+abj-abj+\kappa b^{2}}{a^{2}-\kappa b^{2}}=-1 .
\]
For $\kappa=-1$ the denominator is $a^{2}+b^{2}$, so $S_{-1}$ is defined on
$\R^{2}\setminus\{0\}$; and $b>0$ gives $b/(a^{2}+b^{2})>0$, so the upper half-plane is
preserved.  Involutivity is the standard $(S\circ S)(z)=z$ for a M\"obius involution on a
chart where it is defined, checked directly as
\[
  S_{-1}\bigl(S_{-1}(a,b)\bigr)=(a,b).
\]

\subsection{The exchange lift}

With $J=\begin{pmatrix}0&-1\\ 1&0\end{pmatrix}$:
\[
  J^{2}=\begin{pmatrix}-1&0\\ 0&-1\end{pmatrix}=-I,\qquad
  J^{-1}=-J .
\]
Projectively, $[-\id]=[\id]$, hence $[J]^{2}=[I]$ while $J^{2}=-I$ on the lift.

On the chart $z=t/s$, the action $(t,s)\mapsto(-s,t)$ gives
$z\mapsto(-s)/t=-1/z$.

\subsubsection{The generator matrix}
Multiplication by $j$ in the basis $(1,j)$ is
\[
  L_j=\begin{pmatrix}0&\kappa\\ 1&0\end{pmatrix},\qquad
  L_j^{2}=\begin{pmatrix}\kappa&0\\ 0&\kappa\end{pmatrix}=\kappa I,\qquad
  \det L_j=-\kappa .
\]
For $\kappa=-1$: $L_j=J$, $\det=+1$, order four.  For $\kappa=1$: $L_j^{2}=I$,
$\det=-1$.  For $\kappa=0$: $L_j^{2}=0$, $L_j$ nilpotent.

\subsection{Orientation-preservation of the involution}

In real coordinates,
\[
  S(x,y)=\Bigl(\frac{-x}{x^{2}+y^{2}},\ \frac{y}{x^{2}+y^{2}}\Bigr).
\]
As $S$ is the holomorphic map $z\mapsto -1/z$ with derivative $1/z^{2}$,
\[
  \det DS=\left|\frac{1}{z^{2}}\right|^{2}=\frac{1}{(x^{2}+y^{2})^{2}}>0 .
\]

\subsection{The six-state carrier}

Let $T,S,R$ satisfy $T^{2}=S^{2}=R^{2}=1$ and anticommute in pairs, and set
$\Omega=TSR$.

\subsubsection{Centrality}
For example $T\Omega=T\cdot TSR=T^{2}SR=SR$ and
$\Omega T=TSRT=TS(-TR)=-(TST)R=(TTS)R=SR$; the two agree.  The other two generators are
checked identically.

\subsubsection{The square}
\[
  \Omega^{2}=T\,S\,R\,T\,S\,R=T^{2}\,S\,R\,S\,R=S\,R\,S\,R=-S^{2}R^{2}=-1,
\]
where the third equality moves the fourth factor past two factors (sign $+1$) and the
fourth equality moves the third factor past one factor (sign $-1$).

\subsubsection{The pairing}
\[
  \Omega T=SR,\quad \Omega(SR)=-T,\qquad
  \Omega S=RT,\quad \Omega(RT)=-S,\qquad
  \Omega R=TS,\quad \Omega(TS)=-R,
\]
so on each of the three planes $\langle T,SR\rangle$, $\langle S,RT\rangle$, and
$\langle R,TS\rangle$ the operator $J=\Omega\cdot$ acts as
$\begin{pmatrix}0&-1\\ 1&0\end{pmatrix}$.

\subsubsection{The characteristic polynomial}
Since $J$ acts as the same rotation on three complementary two-planes,
\[
  \chi_J(\lambda)=(\lambda^{2}+1)^{3}=\lambda^{6}+3\lambda^{4}+3\lambda^{2}+1 .
\]

\subsubsection{The signature trichotomy}
Allowing $T^{2}=\varepsilon_T$, $S^{2}=\varepsilon_S$, $R^{2}=\varepsilon_R$ multiplies the
above reduction by $\varepsilon_T\varepsilon_S\varepsilon_R$, giving
$\Omega^{2}=-\varepsilon_T\varepsilon_S\varepsilon_R$.  Hence:
\[
  \varepsilon_T\varepsilon_S\varepsilon_R=+1\ \text{(elliptic)},\quad
  -1\ \text{(hyperbolic)},\quad 0\ \text{(parabolic)} .
\]

\subsubsection{The chirality sign law}
A permutation $\pi$ of the generators acts on the word $TSR$ by its sign, so
$P_\pi\Omega P_\pi^{-1}=\sgn(\pi)\Omega$; in particular a transposition gives
$\Omega\mapsto-\Omega=\Omega^{-1}$.

\subsection{The tri-chart rotation}

With
\[
  R=\begin{pmatrix}0&-1\\ 1&-1\end{pmatrix}
\]
one computes $\mathrm{char}_{R}(\lambda)=\lambda^{2}+\lambda+1$, hence $R^{2}+R+I=0$ and
$R^{3}=I$ by Cayley--Hamilton.  In coordinates: $(x,y)\mapsto(-y,x-y)$, and applying this
three times returns $(x,y)$.  The three charts $(x,y)$, $(y-x,-x)$, $(-y,x-y)$ are permuted
cyclically.

\section{Proofs and finite counts of Part~II}\label{app:proofs}
% Appendix B: Proofs and finite counts of Part II
\label{app:proofs-body}

This appendix supplies the details omitted from Part~II: the adequacy induction, the
countermodels that force each hypothesis of
Definition~\ref{def:adequacy-hypotheses}, the critical-pair configurations, the monotonicity
computation behind strong normalization, and the envelope incidence count.

\subsection{The adequacy induction in detail}

We record the rule-by-rule obligations of Theorem~\ref{thm:adequacy}.

\begin{center}
\begin{tabular}{@{}p{0.20\textwidth}p{0.30\textwidth}p{0.42\textwidth}@{}}
\toprule
Rule & Obligation & Discharged by \\
\midrule
assumption & $\rho$ realises the thread & Definition~\ref{def:realisation} \\
$\wedge$I & family intersection is inhabited & hypothesis (3) \\
$\wedge$E & projection of an inhabited family & hypothesis (3) \\
$\to$I & discharged thread replaced by witness & hypothesis (1) \\
$\to$E & consumed thread matched & hypotheses (1), (2) \\
$\forall$I & a witness slot may be chosen & hypothesis (1) \\
$\forall$E & instantiation stays inside the family & hypothesis (3) \\
atom & witness or emptiness certificate & hypothesis (4) \\
\bottomrule
\end{tabular}
\end{center}

The completeness direction uses hypothesis (4) twice: once to build the canonical assignment
and once to certify emptiness when the goal family is empty.  This is why (4) cannot be
weakened to "either a witness or nothing found": without the certificate an empty family
would be indistinguishable from an unfinished search, and by
Warning~\ref{warn:empty-unresolved} the two must not be conflated.

\subsection{Countermodels for the four hypotheses}

Each hypothesis is necessary for the equivalence of Theorem~\ref{thm:adequacy} in the
following exact sense: there is a fixture on which the equivalence fails, the other three
hypotheses hold, and the failure is attributable to the dropped hypothesis.

\begin{enumerate}
\item \textbf{Without (Cartesian).}  Take a carrier whose slot domains are dependent, so a
  partial assignment may admit no extension.  Then an implication whose premise is
  satisfiable in isolation but not jointly with the context is derivable by the rule while
  its conclusion has empty family.
\item \textbf{Without (Compatibility).}  Take three slots with two constraints such that each
  pair is compatible but the triple is not.  Then pairwise sections are all nonempty while
  the joint family is empty --- the exact situation of
  Proposition~\ref{prop:triadicity-test}, read as a failure of adequacy rather than as
  evidence for triadicity.
\item \textbf{Without (Transformer).}  Take a connective with a transformer that is sound but
  not complete; then a derivable conclusion has nonempty family but the transformer's image
  is strictly larger, so family inclusion does not track derivability.
\item \textbf{Without (Exhaustiveness).}  Take a carrier where some atomic family is empty
  with no certificate; the completeness direction fails at that atom, and the failure is
  reported as \emph{unresolved} rather than as falsity.
\end{enumerate}

\subsection{The four critical-pair configurations}

For two overlapping redexes (Definition~\ref{def:overlapping-redexes}) the ledger blocks
compare as follows.  Write the outer detour's ledger as $L\,[x^{b}]\,Q\,S$ and the inner one
as sitting either inside $Q$ or inside $S$.

\begin{center}
\begin{tabular}{@{}clp{0.46\textwidth}@{}}
\toprule
Case & Position of the inner redex & Relation of the two reduction orders \\
\midrule
(1) & in $Q$, on the discharged spine & joined after one splice step \\
(2) & in $Q$, off the spine & commute exactly \\
(3) & in $S$, adjacent block & joined after one splice step \\
(4) & in $S$, order-adjacent block & commute exactly \\
\bottomrule
\end{tabular}
\end{center}

In cases (1) and (3) the two-step results are related by
Theorem~\ref{thm:substitution-splice} applied to the nested substitution; the freshness
condition is inherited from the outer redex.  In cases (2) and (4) the blocks are disjoint, so
the transports of Theorem~\ref{thm:root-transport} act on independent subterms.

\subsection{Monotonicity of the ledger counts}

Recall that $\Lambda_D$ is the complete proof-tree ledger and that $\operatorname{dchg}$ counts
its discharged records (Definition~\ref{def:proof-ledger}).

\begin{enumerate}
\item A root transport retires exactly the bound target record $x$:
  $\Lambda_R=L\circ[x^{\mathrm{discharged},b}]\circ Q\circ\Lambda_\sigma$ becomes
  $\Lambda_C=L\circ\Lambda_\sigma\circ Q$, with no other record deleted, duplicated,
  relabelled, or renamed.  Both $|\Lambda|$ and $\operatorname{dchg}$ therefore fall by one,
  and the two subterms $Q$ and $S$ are unchanged as trees.
\item An introduction ancestor above the redex discharges the same surviving open record on
  both sides, so it raises both counts equally; either elimination ancestor concatenates the
  same sibling ledger on both sides.  A step strictly inside a context composes with these
  ancestors, and induction along the proof-position path gives the same two decrements at
  the outer root.
\item In the unfolded representation the node count also drops by exactly three, because
  linearity grafts the argument exactly once: the reduct removes the elimination node, the
  introduction node, and the target assumption leaf.
\end{enumerate}

Since $\N$ is well founded, an infinite sequence of certified steps would give an infinite
strictly descending sequence of natural numbers, which is
Theorem~\ref{thm:strong-normalization}.

Linearity enters at the first item and at the node count.  With contraction available the
argument fails, and it fails through the ledger and the graft rather than through a detour
count: see the remark of \S\ref{sec:normalization-body} showing that contracting one
principal detour can expose another.

\subsection{The ledger splice}

For a context $L\,Q\,R$ with $Q$ the transported middle block, the ledger of the composite is
$L\circ\Lambda_\sigma\circ R$ by Theorem~\ref{thm:substitution-splice}.  Concretely, if
$\Lambda_\sigma$ replaces the $i$-th antecedent occurrence by a block of length $\ell_i$, then
the composite ledger has the same left and right blocks and a middle block of length
$\sum_i\ell_i$; freshness is precisely the condition that no occurrence of the original block
appears twice in the image.

\subsection{Envelope incidence count}

The truncated octahedron has $24$ vertices, all of degree $3$, and faces consisting of six
squares and eight hexagons.  Then
\[
  \#\text{face--vertex incidences}=6\cdot4+8\cdot6=24+48=72=24\cdot3,
\]
which agrees with the vertex count.  The relation boundary count of
Proposition~\ref{prop:envelope} is $2m=6$ for $m=3$, matching
Theorem~\ref{thm:counting-bound}.

\section{Six-port incidence and orbit counts}\label{app:six-port-counts}
% Appendix C: Six-port incidence and orbit counts
\label{app:six-port-counts-body}

This appendix records the finite counts used in \S\ref{sec:apertures},
\S\ref{sec:six-port}, and Appendix~\ref{app:proofs}.  Where a count is stated in the source
material, the computation is reproduced here so that it can be checked independently.

\subsection{The support partition}

The carrier of Proposition~\ref{prop:seven-faces} has seven nonempty faces: three vertices,
three edges, and one interior.  Its $2^{7}=128$ supports partition by cardinality as
\[
  128=\underbrace{2}_{|S|\in\{0,7\}}+\underbrace{6}_{|S|\in\{1,6\}}
  +\underbrace{30}_{|S|\in\{2,5\}}+\underbrace{90}_{|S|\in\{3,4\}} .
\]
The two-element classes are the empty and full supports; the six-element classes are the
singletons and their complements; the remaining two classes are self-complementary in the
sense that a support of size $3$ has a complement of size $4$.

The refinement $30=18+12$ separates the supports of size $2$ and $5$ into those fixed by an
orientation-reversing reflection and those exchanged by it; the $18$ are achiral and the $12$
chiral.

\subsection{The $D_3$ orbit count}

The dihedral group $D_3$ of order $6$ acts on the seven faces: it permutes the three vertices
in a $3$-cycle or a transposition, permutes the three edges identically, and fixes the
interior.  Burnside's lemma on the induced action on the $128$ supports gives:

\begin{center}
\begin{tabular}{@{}lcl@{}}
\toprule
Element & Number & Fixed supports \\
\midrule
identity & $1$ & $128$ \\
$120^\circ$ rotation & $2$ & $2\cdot2\cdot2=8$ each \\
reflection & $3$ & $4\cdot4\cdot2=32$ each \\
\bottomrule
\end{tabular}
\end{center}

A rotation acts as a $3$-cycle on the vertices and a $3$-cycle on the edges, so a fixed support
is determined by $2$ choices per $3$-cycle and $2$ choices for the interior: $8$.  A reflection
fixes one vertex and swaps a pair, and likewise for edges, so a fixed support has
$2\cdot2$ choices on vertices and $2\cdot2$ on edges and $2$ on the interior: $32$.
Hence
\[
  \#\text{orbits}=\frac{1}{6}\bigl(128+2\cdot8+3\cdot32\bigr)
  =\frac{240}{6}=40,
\]
reproducing the count of Proposition~\ref{prop:seven-faces}.  Pairing orbits under the
negation involution $C_2^{\neg}$ halves this to the $20$ archetypes stated there.

\subsection{The matching count for the bridge}

Theorem~\ref{thm:m6-bridge} uses two fixed-point-free involutions on six elements.  The exact
counts are as follows.

\begin{proposition}[Matching counts]\label{prop:matching-counts}
On a six-element set there are
\begin{enumerate}
\item $15$ fixed-point-free involutions (perfect matchings);
\item $15\cdot14=210$ ordered pairs of distinct such involutions;
\item $60$ six-cycles, each admitting exactly $2$ alternating ordered $2$-colourings;
\item hence $120$ ordered pairs whose union is connected, and $210-120=90$ whose union is two
  disjoint triangles.
\end{enumerate}
\end{proposition}

\begin{proof}
(i) A perfect matching of six elements is a partition into three unordered pairs, of which
there are $\frac{6!}{2^{3}\,3!}=15$.  (ii) Immediate.  (iii) The number of cycles on six
labelled vertices up to rotation and reflection is $(6-1)!/2=60$, and a cycle of even length
has exactly two alternating $2$-colourings, one for each choice of starting colour.
(iv) $60\cdot2=120$; the rest are unions of two triangles, which arise when the two matchings
are the alternating classes of two disjoint $3$-cycles.
\end{proof}

The count $120$ is the exact size of the family to which
Theorem~\ref{thm:m6-bridge} applies: connected unions are the case in which the presented
boundary is a single $M_6^{\partial}$ rather than a disjoint union of smaller boundaries.

\subsection{Miscellaneous finite counts}

\begin{itemize}
\item The declared maps from six positions into three domains number $3^{6}=729$; the
  retention ledger of the grammar is indexed by these.
\item The set partitions of six positions number $B_6=203$ (the sixth Bell number), which is
  the count of restricted-growth structures used for the allocation ledger.
\item The permutohedron on four symbols has $24$ vertices, matching the $TO_{24}$ vertex
  count of Proposition~\ref{prop:envelope}; the two are different graphs with the same vertex
  count, and the paper does not identify them.
\item The relation boundary count is $2m$; for $m=3$ this is $6$, which is the port count of
  $\bW_6$ and the degree $\deg(\bW_6)$ of Definition~\ref{def:whole-cut-record}.
\end{itemize}

\subsection{What these counts do and do not establish}

Every count in this appendix is an exact statement about a declared finite fixture.  None of
them is a theorem about arbitrary carriers, and in particular:

\begin{itemize}
\item the $40$ orbits and $20$ archetypes are properties of the seven-face triangle, not of
  three-domain carriers in general;
\item the $120$ connected unions count presentations, not fillers: the filler remains open by
  Calibration~\ref{cal:m6-open};
\item the equality of the $TO_{24}$ and permutohedron vertex counts is a numerical coincidence
  unless a map is supplied, and none is supplied here.
\end{itemize}

\section{Claim ledger and source provenance}\label{app:ledger}
% Appendix D: Claim ledger and source provenance
\label{app:ledger-body}

Every substantial result of this paper is traceable to a numbered research note of
\texttt{adva/docs/research}, and the status recorded here is the status the source note
assigns to itself, not an upgraded one.  Where the paper states something the source does
not, the entry says so.

\subsection{How to read the ledger}

The columns are: the paper's item; the source note; and the inherited status.  The status
vocabulary is that of the repository governance:

\begin{itemize}
\item \emph{exact} --- an algebraic or combinatorial identity proved in the paper;
\item \emph{proved} --- proved in the paper under stated hypotheses;
\item \emph{calibration} --- established only by a finite exhaustive check of a declared
  fixture;
\item \emph{proposal} --- a definitional or organisational choice with no theorem;
\item \emph{no-go} --- a proved negative result;
\item \emph{conjecture}, \emph{open problem} --- explicitly unproved.
\end{itemize}

\subsection{Part I: observers, scalars, and the six-state carrier}

\begin{longtable}{@{}p{0.36\textwidth}p{0.16\textwidth}p{0.40\textwidth}@{}}
\toprule
Item & Source & Inherited status \\
\midrule
\endhead
Observer contract, Definition~\ref{def:observer-contract} & 0038 & definition \\
Observation fibre, Definition~\ref{def:observation-fibre} & 0038 & definition \\
Characteristic with residual, Definition~\ref{def:characteristic} & 0038 & definition,
  strengthened by this paper to require the residual explicitly \\
Residual principle, Proposal~\ref{prop:residual-principle} & paper & proposal, derived from
  0038 and the scope file \\
Affine product law, Lemma~\ref{lem:affine-product} & 0038 & exact \\
Two-probe recovery, Proposition~\ref{prop:two-probe} & 0038 & exact within the affine class \\
Interval recovery, Proposition~\ref{prop:interval-recovery} & 0038 & exact; relies on the
  exact-interval observer \\
Outcome discipline, Definition~\ref{def:outcome-discipline} & 0038 & methodological
  requirement \\
Lossy quotient, Proposition~\ref{prop:lossy-quotient} & 0038 & exact \\
Compression and loss, Proposition~\ref{prop:compression-loss} & 0038 & exact \\
Quadratic structure, Proposition~\ref{prop:quadratic-structure} & 0008 & exact \\
Three separators, Proposition~\ref{prop:three-separators} & 0008 & exact \\
Negative inverse, Proposition~\ref{prop:negative-inverse} & 0008 & exact \\
Inversion criterion, Proposal~\ref{prop:inversion-criterion} & 0008 & proposal: selects,
  does not derive \\
Exchange lift, Definition~\ref{def:exchange-lift} & 0008 & definition \\
Order two versus four, Proposition~\ref{prop:order-two-four} & 0008 & exact \\
Chirality, Definition~\ref{def:chirality} & 0008 & definition \\
Not orientation reversal, Proposition~\ref{prop:not-orientation} & 0008 & exact \\
Generator action, Proposition~\ref{prop:generator-action} & 0008 & exact \\
Bridge hypothesis, Calibration~\ref{cal:bridge-hypothesis} & 0008 & calibration \\
Conjugated generators, Proposition~\ref{prop:conjugated-generators} & 0008 & exact \\
Bounded trichotomy, Calibration~\ref{cal:trichotomy-bounded} & 0008 & calibration \\
Three readings of $6=2+2+2$, Proposal~\ref{prop:three-readings} & 0008 & proposal \\
Six-state carrier, Definition~\ref{def:six-state-carrier} & 0009 & definition \\
One complex unit, Proposition~\ref{prop:one-complex-unit} & 0009 & exact but conditional on
  stated anticommutation \\
Algebraic form is an assumption, Calibration~\ref{cal:algebraic-form} & 0009 & calibration \\
Sign law, Proposition~\ref{prop:sign-law} & 0009, 0010 & exact; verified for all six
  permutations \\
Signature trichotomy, Proposition~\ref{prop:signature-trichotomy} & 0009 & exact \\
Temporality is not a metric sign, Warning~\ref{warn:temporality} & 0009 & warning \\
Cut interval, Construction~\ref{con:cut-interval} & 0010 & construction \\
Endogenous Clifford, Theorem~\ref{thm:endogenous-clifford} & 0010 & bounded endogenous
  derivation on one local three-cube \\
Interchange insufficient, Theorem~\ref{thm:interchange-insufficient} & 0010 & no-go \\
Filtration, Proposition~\ref{prop:filtration} & 0011 & exact \\
Regrouping is not chirality, Theorem~\ref{thm:regrouping-not-chirality} & 0011 & no-go \\
\bottomrule
\end{longtable}

\subsection{Part I: triadic inference and crossings}

\begin{longtable}{@{}p{0.36\textwidth}p{0.16\textwidth}p{0.40\textwidth}@{}}
\toprule
Item & Source & Inherited status \\
\midrule
\endhead
Typed readings, Definition~\ref{def:typed-readings} & 0038 & definition \\
Variance signature, Proposition~\ref{prop:variance-signature} & 0038 (also 0039, 0040) &
  proved inside the finite affine model; not a general theorem \\
Scope of the signature, Calibration~\ref{cal:variance-scope} & 0038 & calibration \\
Triadic inference, Proposition~\ref{prop:triadic-inference} & 0038 & definitional
  identification \\
Deformed crossing, Proposition~\ref{prop:deformed-crossing} & 0038 & exact extensional law;
  explicitly not confluent \\
Square orbit quotient, Proposition~\ref{prop:square-orbit} & 0039 & exact \\
Two-sheeted covering, Proposition~\ref{prop:square-covering} & 0039 & exact \\
Relation-valued reverse, Proposition~\ref{prop:relation-valued} & 0039 & exact \\
Scaling crosses, Proposition~\ref{prop:scaling-crosses} & 0039 & exact \\
Translation obstruction, Theorem~\ref{thm:translation-obstruction} & 0039 & no-go; forbids an
  \emph{affine} crossing only \\
Word-to-graph boundary, Proposition~\ref{prop:word-to-graph} & 0039 & structural consequence \\
Constructive fibre, Calibration~\ref{cal:constructive-fibre} & 0039 & exact finite witness \\
Degree profile, Definition~\ref{def:degree-profile} & 0040 & definition \\
Cube degrees, Calibration~\ref{cal:cube-degrees} & 0040 & calibration \\
Criticality warning, Warning~\ref{warn:criticality} & 0040 & warning \\
Equivariance law, Proposition~\ref{prop:equivariance} & 0040 & exact \\
Equivariant characteristics, Proposal~\ref{prop:equivariant-characteristics} & 0040 &
  proposal \\
Binomial coaction, Proposition~\ref{prop:binomial-coaction} & 0040 & exact \\
Translation expands, Corollary~\ref{cor:translation-expands} & paper & corollary of
  Proposition~\ref{prop:binomial-coaction} \\
Isogeny three readings, Theorem~\ref{thm:isogeny-three-readings} & 0041 & external classical
  theorem, restated \\
Isogeny novelty, Calibration~\ref{cal:isogeny-novelty} & 0041 & calibration: the bridge is
  prior art \\
Characteristics as groupoid, Proposal~\ref{prop:characteristic-groupoid} & 0041 & proposal \\
Three cusps, Proposition~\ref{prop:three-cusps} & 0034 & standard exact geometry \\
Mod-two erases chirality, Remark~\ref{rem:mod2-erases} & 0034, 0010 & no-go, repaired by
  0037 \\
Oriented integer lift, Construction~\ref{con:oriented-lift} & 0037 & construction \\
Tri-chart coherence, Lemma~\ref{lem:tri-chart} & 0037 & lemma, direct calculation \\
Legendre three readings, Proposition~\ref{prop:legendre-three-readings} & 0042 & classical
  local facts; the reading identification is open \\
Bare monodromy unipotent, No-go~\ref{nogo:unipotent} & 0034 & no-go \\
Central sign closure, Proposition~\ref{prop:central-sign-closure} & 0042 & exact \\
Label transport, Proposition~\ref{prop:label-transport} & 0042 & exact \\
Crossing prism, Construction~\ref{con:prism} & 0043 (prism) & construction \\
Crossing squares, Theorem~\ref{thm:crossing-squares} & 0043 (prism) & exact for two pairs,
  residual for the third \\
Two residual channels, Proposition~\ref{prop:two-residual-channels} & 0043 (prism) & exact \\
Defect decomposes, Theorem~\ref{thm:defect-decomposes} & 0043 (prism) & no-go; the paper's
  principal negative calibration \\
Coherence defect, Proposal~\ref{prop:coherence-defect} & 0042 & proposal \\
Triadic crossing, Conjecture~\ref{conj:triadic-crossing} & 0042 & conjecture \\
\bottomrule
\end{longtable}

\subsection{Part I: through-crossings, logic, and the residual catalogue}

\begin{longtable}{@{}p{0.36\textwidth}p{0.16\textwidth}p{0.40\textwidth}@{}}
\toprule
Item & Source & Inherited status \\
\midrule
\endhead
Through-crossing, Definition~\ref{def:through-crossing} & 0047 & definition \\
Crossing is not an arrow, Proposal~\ref{prop:crossing-not-arrow} & 0047 & proposal \\
Resolution torsor, Definition~\ref{def:resolution-torsor} & 0047 & definition \\
Monodromy as difference, Theorem~\ref{thm:monodromy-difference} & 0047 & proved \\
Legendre reduction, Proposition~\ref{prop:legendre-reduction} & 0047 & proved from the base
  change \\
Cubical annulus, Construction~\ref{con:cubical-annulus} & 0048 (cellular) & construction \\
Shear transvection, Proposition~\ref{prop:shear-transvection} & 0048 (cellular) & proved \\
Vertex not path, Theorem~\ref{thm:vertex-not-path} & 0048 (cellular) & no-go \\
Marked identity, Proposition~\ref{prop:marked-identity} & 0048 (cellular) & proved \\
Through form, Definition~\ref{def:through-form} & 0070 & definition \\
Order and compatibility, Proposition~\ref{prop:order-compat-independent} & 0070 & proved \\
Resolution count, Proposition~\ref{prop:resolution-count} & 0070 & exact count \\
Figure-eight, Theorem~\ref{thm:figure-eight} & 0071 & calibration on one word \\
Double periods, Proposition~\ref{prop:double-periods} & 0072 & exact arithmetic \\
Return is not period, No-go~\ref{nogo:return-not-period} & 0073 & no-go \\
Satisfaction form, Definition~\ref{def:satisfaction-form} & 0044 & definition \\
Opposite-pair sections, Definition~\ref{def:opposite-pair-sections} & 0044 & definition \\
Generated propositions, Proposal~\ref{prop:generated-propositions} & 0046 & proposal \\
Support preorder, Definition~\ref{def:support-preorder} & 0044 & definition \\
Refinement preorder, Proposition~\ref{prop:refinement-preorder} & 0044 & exact \\
Observer quotient, Calibration~\ref{cal:observer-quotient} & 0044 & calibration \\
Adjoint quantifiers, Proposal~\ref{prop:adjoint-quantifiers} & 0046 & proposal \\
Triadicity test, Proposal~\ref{prop:triadicity-test} & 0046 & proposal, falsifiable \\
Triadicity fails, Calibration~\ref{cal:triadicity-fails} & 0044 & calibration \\
Three grades of logicality, Proposal~\ref{prop:three-grades} & 0045 & proposal \\
Bivalence warning, Warning~\ref{warn:bivalence-completeness-decidability} & 0045 & warning \\
Directional energy, Definition~\ref{def:directional-energy} & 0048 (energy) & definition \\
Energy split, Proposition~\ref{prop:energy-split} & 0048 (energy) & exact \\
Energy invariance, Proposition~\ref{prop:energy-invariance} & 0048 (energy) & exact \\
Zero-energy trap, No-go~\ref{nogo:zero-energy-trap} & 0048 (energy) & no-go \\
Spectrum versus history, Theorem~\ref{thm:spectrum-history} & 0001, 0002 & no-go \\
Affine spectrum, Proposition~\ref{prop:affine-spectrum} & 0002 & exact \\
Marked second jet, Theorem~\ref{thm:second-jet} & 0003 & proved under the marked hypothesis \\
Co-generation, Theorem~\ref{thm:cogeneration} & 0019 & proved for the declared probe class \\
Observation tower, Theorem~\ref{thm:observation-tower} & 0013 & no-go \\
Objectification, No-go~\ref{nogo:objectification} & 0021 & no-go \\
Effective one-sidedness, Proposition~\ref{prop:effective-onesided} & 0033 & standard, restated \\
Mass is not fibre, Proposition~\ref{prop:mass-not-fibre} & 0033 & proved \\
No finite coding, Proposition~\ref{prop:no-finite-coding} & 0037 & proved \\
Snapshot not empty, Proposition~\ref{prop:snapshot-not-empty} & 0023 & exact witness \\
Zero-event fibre, Proposition~\ref{prop:zero-event-fibre} & 0027 & no-go \\
Trace sides, Proposition~\ref{prop:trace-sides} & 0114 & calibration \\
Count factorization, Proposition~\ref{prop:count-factorization} & 0124 & exact count plus
  explicit obstruction \\
Period not history, Proposition~\ref{prop:period-not-history} & 0043 (prism) & no-go \\
Order obstruction, Theorem~\ref{thm:order-obstruction} & 0126 & proved with exact minimum \\
Scale forgetting, Proposition~\ref{prop:scale-forgetting} & 0043 (scale) & exact \\
Winding insufficient, No-go~\ref{nogo:winding} & 0043 (path-bound) & no-go \\
\bottomrule
\end{longtable}

\subsection{Part II: apertures, deduction, and the kernel}

\begin{longtable}{@{}p{0.36\textwidth}p{0.16\textwidth}p{0.40\textwidth}@{}}
\toprule
Item & Source & Inherited status \\
\midrule
\endhead
Aperture-first, Proposal~\ref{prop:aperture-first} & 0079 & proposal \\
Typed aperture, Definition~\ref{def:aperture} & 0079 & definition \\
Open and close, Proposal~\ref{prop:open-close} & 0079 & proposal \\
Aperture status, Definition~\ref{def:aperture-status} & 0079 & definition \\
Terminal worlds, Definition~\ref{def:terminal-worlds} & 0081 & definition \\
Seven faces, Proposition~\ref{prop:seven-faces} & 0096 & calibration \\
Mask direction, No-go~\ref{nogo:mask-direction} & 0096 & no-go \\
Filling fibre, Definition~\ref{def:filling-fibre} & 0081 & definition \\
Fibre connectives, Proposal~\ref{prop:fibre-connectives} & 0081 & proposal \\
Surface observer, Proposal~\ref{prop:surface-observer} & 0080 & proposal \\
Geometric warning, Warning~\ref{warn:geometric-conditional} & 0080 & warning \\
Threaded judgment, Definition~\ref{def:threaded-judgment} & 0084 & definition \\
Realisation, Definition~\ref{def:realisation} & 0082 & definition \\
Adequacy hypotheses, Definition~\ref{def:adequacy-hypotheses} & 0082 & hypotheses \\
Adequacy, Theorem~\ref{thm:adequacy} & 0082 & proved under the four hypotheses \\
Adequacy boundary, Calibration~\ref{cal:adequacy-boundary} & 0082 & calibration \\
Substitution natural, Theorem~\ref{thm:substitution-natural} & 0082 & proved \\
$\mathrm{TND}_0$, Definition~\ref{def:tnd0} & 0084 & definition \\
Ledger judgment, Definition~\ref{def:ledger-judgment} & 0084 & definition \\
Soundness, Theorem~\ref{thm:soundness} & 0084 & proved \\
Resource grading, Definition~\ref{def:grade} & 0084 & definition \\
Grading invariant, Proposition~\ref{prop:grade-invariant} & 0084 & proved \\
Graded non-derivability, Corollary~\ref{cor:graded-nonderivability} & 0084 & proved,
  with the explicit witness of Example~\ref{ex:graded-witness} \\
Incompleteness, Theorem~\ref{thm:incompleteness} & 0084 & proved \\
Conflict aperture, Proposal~\ref{prop:conflict-aperture} & 0083 & proposal \\
Common lift, Proposition~\ref{prop:common-lift} & 0083 & proved \\
One-hole context, Definition~\ref{def:one-hole-context} & 0085 & definition \\
Ordered substitution, Definition~\ref{def:ordered-substitution} & 0085 & definition \\
Ledger splice, Theorem~\ref{thm:substitution-splice} & 0085 & proved \\
Fresh cut, Corollary~\ref{cor:fresh-cut} & 0085 & conditional \\
Raw ledger fails, No-go~\ref{nogo:raw-ledger} & 0085 & no-go \\
Ledger transport, Proposal~\ref{prop:ledger-transport} & 0086 & proposal \\
Freshness reflects, Lemma~\ref{lem:freshness-reflects} & 0086 & proved \\
Root transport, Theorem~\ref{thm:root-transport} & 0086 & proved \\
Contextual closure, Corollary~\ref{cor:contextual-closure} & 0086 & proved \\
Proof-tree ledger and discharge count, Definition~\ref{def:proof-ledger} & 0086 &
  definition \\
Ledger descent, Theorem~\ref{thm:ledger-descent} & 0086 & proved \\
Strong normalization, Theorem~\ref{thm:strong-normalization} & 0086 & proved \\
Overlapping redexes, Definition~\ref{def:overlapping-redexes} & 0093 & definition \\
Four paths, Proposition~\ref{prop:four-paths} & 0093 & proved \\
Substitution and graft, Lemma~\ref{lem:substitution-graft} & 0093 & proved \\
Strong diamond, Theorem~\ref{thm:strong-diamond} & 0093 & proved \\
Confluence, Corollary~\ref{cor:confluence} & 0093 & proved per fixed proof \\
Confluence is not coherence, No-go~\ref{nogo:confluence-coherence} & 0093 & no-go \\
Normal and neutral, Definition~\ref{def:normal-neutral} & 0094 & definition \\
Subformula property, Theorem~\ref{thm:subformula} & 0094 & proved \\
Leftmost head, Proposition~\ref{prop:leftmost-head} & 0094 & proved \\
Decidability, Theorem~\ref{thm:decidability} & 0094 & proved \\
Skeleton positioning, Proposition~\ref{prop:skeleton-position} & 0094 & proved \\
No compactification, Warning~\ref{warn:no-compactification} & 0093, 0094 & warning \\
Whole-cut carrier, Definition~\ref{def:whole-cut-carrier} & 0101 & definition \\
Six-port contract, Definition~\ref{def:six-port-contract} & 0101 & declared contract \\
Counting bound, Theorem~\ref{thm:counting-bound} & 0101 & proved under the contract \\
Six is conditional, Warning~\ref{warn:six-conditional} & 0101 & warning \\
The carrier $\bW_6$, Definition~\ref{def:w6} & 0101 & definition \\
Line and circle as views, Proposal~\ref{prop:line-circle-views} & 0101, 0102 & proposal \\
Whole-cut record, Definition~\ref{def:whole-cut-record} & 0102 & definition \\
Collision conservation, Proposition~\ref{prop:collision-conservation} & 0102 & proved \\
Cell--carrier--view, Proposal~\ref{prop:cell-carrier-view} & 0103 & proposal \\
Envelope count, Proposition~\ref{prop:envelope} & 0103 & exact count \\
$M_6$ bridge, Theorem~\ref{thm:m6-bridge} & 0104 & proved under explicit hypotheses \\
Filler open, Calibration~\ref{cal:m6-open} & 0104, 0113 & calibration: no filler \\
Relative initiality, Proposal~\ref{prop:relative-initiality} & 0101 & proposal \\
Not universal, Warning~\ref{warn:not-universal} & 0101 & warning \\
Matching counts, Proposition~\ref{prop:matching-counts} & paper & elementary count, check of
  Theorem~\ref{thm:m6-bridge} \\
\bottomrule
\end{longtable}

\subsection{Source notes not absorbed as main text}

The following clusters of the source corpus are mathematical but belong to neither part of
this paper, and are recorded here so that the omission is visible rather than accidental.

\begin{itemize}
\item \textbf{0004--0007}: whole-cut program cells, the Alexandrov topology of completed
  causal pasts, frontier transport with its interchange square, and the occurrence-affine cut
  lift.  This is a theory of causal cuts with its own space--time duality, presupposing
  Paper~I's histories.
\item \textbf{0015--0018}: occurrence-resolved backward probe pairing, expression-valued
  backward transport, and the intensional cut-composition law.  Its natural home is with
  Paper~IV's continuation residuals.
\item \textbf{0020, 0022--0031}: the finite dual-cut surgery programme, the relational and
  polarized $PSP^{\ast}$ factorizations, and boundary polarization.  Much of this is stated
  relative to checked finite artifacts and would require its own careful separation of
  mathematics from fixture.
\item \textbf{0033, 0035--0037}: the $\Omega$-type effective boundary, the decorated handle
  operator, the triangular symbolic calculus, and the finite-observer tri-cusp reduction.
  Only Proposition~\ref{prop:effective-onesided} and
  Proposition~\ref{prop:no-finite-coding} are imported.
\item \textbf{0050--0058}: the triadic transformation monoid, context-generated proposition
  algebras, causal evidence layers, and the braid/gate comparison.  These are the finite
  models that motivate \S\ref{sec:satisfaction-logic} but are not restated.
\item \textbf{0059--0069, 0074--0078, 0087--0092, 0105--0127}: the bracket and shell
  calculi, the research-instrument reports, the multi-hole syntax, the venture and
  failure-frontier programme, and the later formation-gate notes.
\end{itemize}

Fifteen of the 133 surveyed files are predominantly engineering and contribute no main-text
mathematics: \texttt{0026}, \texttt{0050} (LABS search), \texttt{0054} (sorting networks),
\texttt{0074}, \texttt{0090}, \texttt{0092}, the three \texttt{0102} follow-ups,
\texttt{0108}, \texttt{0110}, \texttt{0115}, \texttt{0116}, \texttt{0123}, and one SVG
figure.


\bibliographystyle{plain}
\bibliography{aeg-paper-5}
\end{document}
